\chapter{E-Trike System Architecture}\label{e-trike-system-architecture}

Four-node distributed control: \textbf{Jetson Orin} (ROS 2 perception/planning),
\textbf{RT ESP32-S3} (realtime physics, steering, brake \& CAN gateway),
\textbf{SYS ESP32-S3} (safety \& body control), \textbf{MTR STM32} (motor
actuation).

\begin{quote}
\textbf{Variant:} A consolidated single-controller version exists at
\href{rt-aurix-lite/rt-aurix-lite-architecture.md}{\texttt{rt-aurix-lite/rt-aurix-lite-architecture.md}}
--- combines RT+SYS into one AURIX TC3xx on a single CAN bus. This document
describes the distributed reference architecture. Both variants share the same
CAN protocol and actuator interfaces.
\end{quote}

Two physical CAN buses at 500 kbit/s. RT bridges selected messages between
buses. Actuators are \textbf{SYNTREE} CAN modules: EPS-C (steer-by-wire) and SEB
(electro-hydraulic brake). \textbf{Mode-gated dual control (Option D):} RT
directly commands both EPS-C and SEB in AUTO mode; SYS commands SEB in MANUAL
mode (lever pass-through) and on ESTOP. This ensures no single MCU failure takes
both actuators, and AUTO-mode brake+steer are 1-hop from the kinematics engine.
Motor control is on a dedicated STM32 board (MTR) for safety isolation per ISO
26262 EGAS 3-level concept.

\begin{quote}
\textbf{Autoware.Auto v0.0.4-alpha:} A new ROS 2 node on the Jetson
(\texttt{autoware\_vehicle\_bridge}) provides Autoware.Auto topic compatibility
(\texttt{AckermannControlCommand} in,
\texttt{VelocityReport}/\texttt{SteeringReport} out). The CAN protocol across
both buses is unchanged. See §5.1.
\end{quote}

\begin{center}\rule{0.5\linewidth}{0.5pt}\end{center}

\section{1. Topology}\label{topology}

\begin{verbatim}
  ┌────────────────── High-Level CAN (500 kbit/s) ──────────────────┐
  │                                                                  │
  │  ┌──────────┐            ┌──────────────┐                       │
  │  │  Jetson  │            │  RT ESP32-S3 │                       │
  │  │  Orin │            │              │                       │
  │  │          │            │ Physics      │                       │
  │  │ ROS 2    │            │ Steering     │                       │
  │  │ Planning │            │ Gateway      │                       │
  │  └────┬─────┘            │              │                       │
  │       │                  └──────┬───────┘                       │
  │  TX:  0x300,0x301,    TX: 0x011,0x120, │                        │
  │       0x302,0x001           0x210,0x220,│                        │
  │                              0x400,0x600,│                       │
  │  RX:  0x011,0x120,          0x001,0x7FC │                        │
  │       0x210,0x220,      RX: 0x300,0x301,│                        │
  │       0x400,0x600,          0x302,0x001 │                        │
  │       0x001,0x7FD                       │                        │
  └─────────────────────────────────────────┘                        │
                                            │                        │
                 ┌──────────────────────────┘                        │
                 │                                                    │
  ┌──────────────▼─────── Low-Level CAN (500 kbit/s) ───────────────┐│
  │                                                                  ││
  │  ┌──────────────┐      ┌──────────────┐                         ││
  │  │  RT ESP32-S3 │      │ SYS ESP32-S3 │                         ││
  │  │  (gateway)   │      │              │                         ││
  │  └──────┬───────┘      │ Safety       │                         ││
  │         │              │ Motor        │                         ││
  │    TX:  0x169,0x204,   │ Brake        │                         ││
  │         0x302,0x001,   └──────┬───────┘                         ││
  │         0x7FD                 │                                  ││
  │                          TX:  0x7B9,0x011,                      ││
  │    RX:  0x001,0x011,          0x012,0x110,                      ││
  │         0x110,0x120,          0x120,0x600,                      ││
  │         0x201,0x600,          0x001,0x7FE                       ││
  │         0x7FD,0x7FE              RX:  0x001,0x204,                    ││
  │                               0x302,0x721,                      ││
  │                               0x7FD                             ││
  └─────────────────────────────────────────────────────────────────┘│
                                        │                            │
                ┌───────────────────────┼────────────────────┐       │
                │                       │                    │       │
          ┌─────▼─────┐          ┌─────▼─────┐        ┌─────▼─────┐ │
          │  SYNTREE  │          │  SYNTREE  │        │   DC-DC   │ │
          │  SEB      │          │  EPS-C    │        │ Converter │ │
          │  (Brake)  │          │ (Steering)│        │ 72V→12V   │ │
          │ 0x7B9 cmd │          │ 0x169 cmd │        │ (0x012)   │ │
          │ 0x721 stat│          │ 0x201 stat│        └───────────┘ │
          └───────────┘          └───────────┘                      │
                                                                    │
  ┌───────────┐                                                    │
  │  Motor    │  (analog: 0–5 V throttle via MCP4725,              │
  │Controller │   72 V gear lines via relay module, from SYS)      │
  └───────────┘                                                    │
\end{verbatim}

\begin{center}\rule{0.5\linewidth}{0.5pt}\end{center}

\section{2. CAN message catalog}\label{can-message-catalog}

\subsection{2.1 Low-level CAN}\label{low-level-can}

\begin{longtable}[]{@{}
  >{\raggedright\arraybackslash}p{(\columnwidth - 14\tabcolsep) * \real{0.0678}}
  >{\raggedright\arraybackslash}p{(\columnwidth - 14\tabcolsep) * \real{0.1017}}
  >{\raggedright\arraybackslash}p{(\columnwidth - 14\tabcolsep) * \real{0.1356}}
  >{\raggedright\arraybackslash}p{(\columnwidth - 14\tabcolsep) * \real{0.2203}}
  >{\raggedright\arraybackslash}p{(\columnwidth - 14\tabcolsep) * \real{0.0847}}
  >{\raggedright\arraybackslash}p{(\columnwidth - 14\tabcolsep) * \real{0.1525}}
  >{\raggedright\arraybackslash}p{(\columnwidth - 14\tabcolsep) * \real{0.1356}}
  >{\raggedright\arraybackslash}p{(\columnwidth - 14\tabcolsep) * \real{0.1017}}@{}}
\toprule\noalign{}
\begin{minipage}[b]{\linewidth}\raggedright
ID
\end{minipage} & \begin{minipage}[b]{\linewidth}\raggedright
Name
\end{minipage} & \begin{minipage}[b]{\linewidth}\raggedright
Sender
\end{minipage} & \begin{minipage}[b]{\linewidth}\raggedright
Receiver(s)
\end{minipage} & \begin{minipage}[b]{\linewidth}\raggedright
DLC
\end{minipage} & \begin{minipage}[b]{\linewidth}\raggedright
Payload
\end{minipage} & \begin{minipage}[b]{\linewidth}\raggedright
Period
\end{minipage} & \begin{minipage}[b]{\linewidth}\raggedright
Prio
\end{minipage} \\
\midrule\noalign{}
\endhead
\bottomrule\noalign{}
\endlastfoot
\texttt{0x001} & SAFETY\_ESTOP & Any & All (bridged to high) & 0 & (none) &
Event & Highest \\
\texttt{0x011} & SYS\_SAFETY\_STS & SYS & RT (→ Jetson) & 2 & u8 estop, u8
hb\_ok & 5 Hz & V.High \\
\texttt{0x012} & SYS\_DCDC\_CMD & SYS & DC-DC converter & 1 & u8 enable & Change
& V.High \\
\texttt{0x110} & SYS\_MODE\_CMD & SYS & RT & 1 & u8 mode (0=M, 1=A, 2=ESTOP) &
Change & High \\
\texttt{0x120} & SYS\_THROTTLE\_STS & MTR & RT (→ Jetson), SYS & 2 & i16
speed\_mmps & 100 Hz & Medium \\
\texttt{0x169} & VCU\_SES\_REQ & RT & EPS-C (steering) & 8 & Angle cmd +
security bytes & 50 Hz & Medium \\
\texttt{0x201} & SES\_STATUS & EPS-C & RT & 8 & Steering angle + status feedback
& 100 Hz & Medium \\
\texttt{0x202} & SES\_ErrInfo & EPS-C & RT & 8 & 25 fault flags (8× L3) & 10 Hz
& Medium \\
\texttt{0x203} & SES\_Version & EPS-C & RT & 8 & SW + HW version & 1 Hz &
Lowest \\
\texttt{0x204} & RT\_DRIVE\_CMD & RT & \textbf{MTR (STM32), SYS} & 5 & i32
speed\_mmps, u8 gear & 100 Hz & Medium \\
\texttt{0x205} & RT\_BRAKE\_CMD & RT & SYS & 4 & i32 brake\_pressure\_kpa & 50
Hz & Low \\
\texttt{0x206} & MTR\_MOTOR\_FBK & MTR & SYS, RT & 4 & i16 actual\_speed, u8
gear\_state, u8 fault\_flags & 50 Hz & Low \\
\texttt{0x302} & HOST\_LIGHT\_CMD & RT (fwd) & SYS & 1 & u8 lights bitfield &
Change & Medium \\
\texttt{0x600} & SYS\_DIAG\_RPT & SYS & RT (→ Jetson) & 8 & diag struct & 1 Hz &
Lowest \\
\texttt{0x6FA} & SES\_Test & EPS-C & RT & 8 & Motor current, ECU temp, supply
voltage & 100 Hz & Lowest \\
\texttt{0x6FB} & SEB\_Test & SEB & SYS & 8 & Motor current, ECU temp, supply
voltage & 100 Hz & Lowest \\
\texttt{0x721} & SEB\_STATUS & SEB & SYS & 8 & Brake stroke + status + error
level feedback & 100 Hz & Lowest \\
\texttt{0x731} & SEB\_ErrInfo & SEB & SYS & 8 & 23 fault flags (16× L3) & 10 Hz
& Lowest \\
\texttt{0x741} & SEB\_Version & SEB & SYS & 8 & SW + HW version & 1 Hz &
Lowest \\
\texttt{0x7B9} & VCU\_SEB\_REQ & RT (AUTO) / SYS (MANUAL, ESTOP) † & SEB (brake)
& 8 & Stroke/pressure cmd + security & 50 Hz & Lowest \\
\texttt{0x7FD} & RT\_HEARTBEAT & RT & SYS & 1 & u8 alive\_ctr & 2 Hz & Lowest \\
\texttt{0x7FE} & SYS\_HEARTBEAT & SYS & RT & 1 & u8 alive\_ctr & 10 Hz &
Lowest \\
\end{longtable}

\begin{quote}
\textbf{ID note}: SYNTREE units are preprogrammed and cannot be reconfigured.
EPS-C uses factory command \texttt{0x169} and status \texttt{0x201}, plus
diagnostic frames \texttt{0x202} (err info), \texttt{0x203} (version),
\texttt{0x6FA} (telemetry). SEB uses factory command \texttt{0x7B9} and status
\texttt{0x721}, plus diagnostic frames \texttt{0x731} (err info), \texttt{0x741}
(version), \texttt{0x6FB} (telemetry). \texttt{RT\_DRIVE\_CMD} is placed at
\texttt{0x204} to avoid collision with EPS-C \texttt{0x169}. \texttt{0x205}
RT\_BRAKE\_CMD avoids collision with \texttt{0x202} and \texttt{0x203}.

\textbf{† \texttt{0x7B9} dual-sender note:} Architecture §6.2 Option D specifies
RT sends \texttt{0x7B9} directly in AUTO (1-hop from kinematics), SYS in
MANUAL/ESTOP. \textbf{Current implementation (v0.0.4):} SYS is the sole
\texttt{0x7B9} sender in all modes --- RT sends \texttt{0x205\ RT\_BRAKE\_CMD}
to SYS, which converts to SEB protocol. Direct RT→SEB transmission is planned
(gap \#13).
\end{quote}

\subsection{2.2 High-level CAN}\label{high-level-can}

\begin{longtable}[]{@{}
  >{\raggedright\arraybackslash}p{(\columnwidth - 14\tabcolsep) * \real{0.0678}}
  >{\raggedright\arraybackslash}p{(\columnwidth - 14\tabcolsep) * \real{0.1017}}
  >{\raggedright\arraybackslash}p{(\columnwidth - 14\tabcolsep) * \real{0.1356}}
  >{\raggedright\arraybackslash}p{(\columnwidth - 14\tabcolsep) * \real{0.2203}}
  >{\raggedright\arraybackslash}p{(\columnwidth - 14\tabcolsep) * \real{0.0847}}
  >{\raggedright\arraybackslash}p{(\columnwidth - 14\tabcolsep) * \real{0.1525}}
  >{\raggedright\arraybackslash}p{(\columnwidth - 14\tabcolsep) * \real{0.1356}}
  >{\raggedright\arraybackslash}p{(\columnwidth - 14\tabcolsep) * \real{0.1017}}@{}}
\toprule\noalign{}
\begin{minipage}[b]{\linewidth}\raggedright
ID
\end{minipage} & \begin{minipage}[b]{\linewidth}\raggedright
Name
\end{minipage} & \begin{minipage}[b]{\linewidth}\raggedright
Sender
\end{minipage} & \begin{minipage}[b]{\linewidth}\raggedright
Receiver(s)
\end{minipage} & \begin{minipage}[b]{\linewidth}\raggedright
DLC
\end{minipage} & \begin{minipage}[b]{\linewidth}\raggedright
Payload
\end{minipage} & \begin{minipage}[b]{\linewidth}\raggedright
Period
\end{minipage} & \begin{minipage}[b]{\linewidth}\raggedright
Prio
\end{minipage} \\
\midrule\noalign{}
\endhead
\bottomrule\noalign{}
\endlastfoot
\texttt{0x001} & SAFETY\_ESTOP & RT (fwd), Jetson & Jetson, RT & 0 & (none) &
Event & Highest \\
\texttt{0x011} & SYS\_SAFETY\_STS & RT (fwd) & Jetson & 2 & u8 estop, u8 hb\_ok
& 5 Hz & V.High \\
\texttt{0x120} & SYS\_THROTTLE\_STS & RT (fwd) & Jetson & 2 & i16 speed\_mmps &
100 Hz & Medium \\
\texttt{0x210} & RT\_STATE\_RPT & RT & Jetson & 3 & u8 mode, u8 steer\_valid, u8
reversing & 10 Hz & Low \\
\texttt{0x220} & RT\_PID\_RPT & RT & Jetson & 6 & RESERVED --- future
closed-loop PID telemetry & --- (inactive) & Low \\
\texttt{0x300} & HOST\_DRIVE\_CMD & Jetson & RT & 8 & i32 speed\_mmps, i24
yaw\_rate\_mrad\_s, u8 gear & ≤100 Hz & Medium \\
\texttt{0x301} & HOST\_BRAKE\_REQ & Jetson & RT & 4 & i32 brake\_pressure\_kpa &
Demand & Medium \\
\texttt{0x302} & HOST\_LIGHT\_CMD & Jetson & RT (→ SYS) & 1 & u8 lights bitfield
& Change & Medium \\
\texttt{0x400} & HOST\_OBSTACLE\_DIST & Jetson & RT & 4 & u32 distance\_mm & 10
Hz & Low \\
\texttt{0x600} & SYS\_DIAG\_RPT & RT (fwd) & Jetson & 8 & diag struct & 1 Hz &
Lowest \\
\texttt{0x7FD} & RT\_HEARTBEAT & RT & Jetson & 1 & u8 alive\_ctr & 2 Hz &
Lowest \\
\texttt{0x7FC} & JETSON\_HEARTBEAT & Jetson & RT & 1 & u8 alive\_ctr & 2 Hz &
Lowest \\
\end{longtable}

\begin{quote}
Bit-level signal layouts in
\href{can-dictionary.md}{\texttt{can-dictionary.md}}.
\end{quote}

\subsection{2.3 RT CAN gateway --- message handling by
category}\label{rt-can-gateway-message-handling-by-category}

RT is the only dual-bus node. Every CAN message falls into exactly one of three
categories:

\subsubsection{Category 1: Transparent
forward}\label{category-1-transparent-forward}

RT copies the frame to the other bus unchanged --- same CAN ID, same payload,
same DLC. The receiving node cannot tell whether RT or the original sender
transmitted it.

\begin{longtable}[]{@{}
  >{\raggedright\arraybackslash}p{(\columnwidth - 4\tabcolsep) * \real{0.3235}}
  >{\raggedright\arraybackslash}p{(\columnwidth - 4\tabcolsep) * \real{0.4118}}
  >{\raggedright\arraybackslash}p{(\columnwidth - 4\tabcolsep) * \real{0.2647}}@{}}
\toprule\noalign{}
\begin{minipage}[b]{\linewidth}\raggedright
Direction
\end{minipage} & \begin{minipage}[b]{\linewidth}\raggedright
IDs forwarded
\end{minipage} & \begin{minipage}[b]{\linewidth}\raggedright
Example
\end{minipage} \\
\midrule\noalign{}
\endhead
\bottomrule\noalign{}
\endlastfoot
Low → High & \texttt{0x001}, \texttt{0x011}, \texttt{0x120}, \texttt{0x206},
\texttt{0x600} & SYS sends \texttt{0x011} on low → RT transmits \texttt{0x011}
on high → Jetson receives \\
High → Low & \texttt{0x001}, \texttt{0x302} & Jetson sends \texttt{0x302} on
high → RT transmits \texttt{0x302} on low → SYS receives \\
\end{longtable}

\begin{quote}
\texttt{0x001} is the only bidirectional forward --- ESTOP originates on either
bus and must reach all nodes on both buses.
\end{quote}

\subsubsection{Category 2: Consumed by RT → different message generated on other
bus}\label{category-2-consumed-by-rt-different-message-generated-on-other-bus}

RT receives a command on one bus, processes it internally, and transmits a
\textbf{different} CAN ID with \textbf{different} payload on the other bus.
These are not forwards --- they are translations.

\begin{longtable}[]{@{}
  >{\raggedright\arraybackslash}p{(\columnwidth - 8\tabcolsep) * \real{0.3175}}
  >{\raggedright\arraybackslash}p{(\columnwidth - 8\tabcolsep) * \real{0.0794}}
  >{\raggedright\arraybackslash}p{(\columnwidth - 8\tabcolsep) * \real{0.1746}}
  >{\raggedright\arraybackslash}p{(\columnwidth - 8\tabcolsep) * \real{0.3492}}
  >{\raggedright\arraybackslash}p{(\columnwidth - 8\tabcolsep) * \real{0.0794}}@{}}
\toprule\noalign{}
\begin{minipage}[b]{\linewidth}\raggedright
Inbound (consumed)
\end{minipage} & \begin{minipage}[b]{\linewidth}\raggedright
Bus
\end{minipage} & \begin{minipage}[b]{\linewidth}\raggedright
Processing
\end{minipage} & \begin{minipage}[b]{\linewidth}\raggedright
Outbound (generated)
\end{minipage} & \begin{minipage}[b]{\linewidth}\raggedright
Bus
\end{minipage} \\
\midrule\noalign{}
\endhead
\bottomrule\noalign{}
\endlastfoot
\texttt{0x300} HOST\_DRIVE\_CMD & High & Kinematics → \texttt{ResolvedSetpoint}
& \texttt{0x204} RT\_DRIVE\_CMD + \texttt{0x169} VCU\_SES\_REQ & Low \\
\texttt{0x301} HOST\_BRAKE\_REQ & High & Max-select arbitration → \texttt{0x205}
RT\_BRAKE\_CMD & \texttt{0x205} (RT→SYS, AUTO only); SYS converts to
\texttt{0x7B9} → SEB & Low \\
\end{longtable}

\subsubsection{Category 3: Bus-local (never forwarded, never
regenerated)}\label{category-3-bus-local-never-forwarded-never-regenerated}

These messages serve only nodes on a single bus. RT neither forwards nor
translates them.

\begin{longtable}[]{@{}
  >{\raggedright\arraybackslash}p{(\columnwidth - 4\tabcolsep) * \real{0.2778}}
  >{\raggedright\arraybackslash}p{(\columnwidth - 4\tabcolsep) * \real{0.2778}}
  >{\raggedright\arraybackslash}p{(\columnwidth - 4\tabcolsep) * \real{0.4444}}@{}}
\toprule\noalign{}
\begin{minipage}[b]{\linewidth}\raggedright
Bus
\end{minipage} & \begin{minipage}[b]{\linewidth}\raggedright
IDs
\end{minipage} & \begin{minipage}[b]{\linewidth}\raggedright
Reason
\end{minipage} \\
\midrule\noalign{}
\endhead
\bottomrule\noalign{}
\endlastfoot
Low only & \texttt{0x012}, \texttt{0x110}, \texttt{0x169}, \texttt{0x204},
\texttt{0x205}, \texttt{0x7B9} & DC-DC, mode, steering, drive, brake, and SEB
commands. RT-generated or SYS-generated --- never leave low bus. \\
Low only & \texttt{0x201}, \texttt{0x721} & SYNTREE status feedback --- RT/SYS
consume locally for boot sync and monitoring \\
High only & \texttt{0x210}, \texttt{0x220}, \texttt{0x400} & RT telemetry ---
generated by RT for Jetson consumption only \\
Both (independent) & \texttt{0x7FD}, \texttt{0x7FE}, \texttt{0x7FC} & Per-node
heartbeat with alive counter --- MUST NOT be bridged (see §8.6). Each bus has
its own liveness domain. \\
\end{longtable}

\begin{center}\rule{0.5\linewidth}{0.5pt}\end{center}

\section{3. Mode state machine}\label{mode-state-machine}

\begin{verbatim}
         ┌──────────┐
    ┌───▶│  MANUAL  │◀───┐
    │    └─────┬────┘    │
    │     push btn=AUTO  push btn=MANUAL
    │          │          │
    │    ┌─────▼────┐    │
    │    │   AUTO   │    │
    │    └─────┬────┘    │
    │          │          │
    │  ESTOP button / CAN 0x001 / HB timeout
    │          │          │
    │    ┌─────▼────┐    │
    │    │  ESTOP   │────┘
    │    └─────┬────┘
    │         │ START button (GPIO32)
    │         ▼
    └──────── MANUAL
\end{verbatim}

\begin{longtable}[]{@{}
  >{\raggedright\arraybackslash}p{(\columnwidth - 2\tabcolsep) * \real{0.3750}}
  >{\raggedright\arraybackslash}p{(\columnwidth - 2\tabcolsep) * \real{0.6250}}@{}}
\toprule\noalign{}
\begin{minipage}[b]{\linewidth}\raggedright
Mode
\end{minipage} & \begin{minipage}[b]{\linewidth}\raggedright
Behavior
\end{minipage} \\
\midrule\noalign{}
\endhead
\bottomrule\noalign{}
\endlastfoot
\textbf{MANUAL} & Rider steers / rides throttle. \textbf{MTR} reads throttle ADC
+ gear sense → pass-through via MCP4725 + relays. Brake lever → SYS GPIO → CAN
\texttt{0x7B9} → SEB. EPS-C standalone (RT idle). DC-DC on. Mode gated by SYS
\texttt{0x110}. \\
\textbf{AUTO} & Jetson \texttt{AckermannControlCommand} → high CAN
\texttt{0x300} → RT kinematics → low CAN \texttt{0x204} (MTR: speed+gear) +
\texttt{0x169} (EPS-C: angle). \textbf{MTR} drives MCP4725 + gear relays
following \texttt{0x204}. Lights from Jetson via \texttt{0x302} (RT fwd). Brake
via \texttt{0x7B9}. Mode gated by SYS \texttt{0x110}. \\
\textbf{ESTOP} & MCP4725 = 0 V, all gear outputs OFF, \texttt{0x7B9} stroke=max
(full brake), steering ramps to 0° at 20°/s via active \texttt{0x169} (unless
obstacle-triggered → hold then silent-stop), DC-DC ON (\texttt{0x012\ enable=1}
--- maintains 12V for MCUs, CAN transceivers, brake light), 12V accessory relay
OFF (kills headlight, turn signals, mode bulbs). Brake light ON (powered from
always-on DC-DC rail, not through accessory relay). Exit: \textbf{START button}
or \textbf{MODE long-press (3s)} → MANUAL. Steering ramp completes before 0x169
handoff (brake/motor/lights transition immediately; steering defers until
centered). Or power-cycle. \\
\end{longtable}

\begin{center}\rule{0.5\linewidth}{0.5pt}\end{center}

\section{4. Signal flow}\label{signal-flow}

\subsection{4.1 Manual mode}\label{manual-mode}

\begin{verbatim}
Throttle grip (0–5V) ──► MTR ADC ──► MTR MCP4725 (0–5V) ──► Motor controller
Gear selector (72V)  ──► TLP281 opto → MTR GPIO ──► relay module → 72V → ECU
Brake lever           ──► SYS GPIO ──► CAN 0x7B9 → SEB (stroke=MAX if pressed)
Steering wheel        ──► EPS-C standalone (RT idle, monitors 0x201)
Signal lights         ──► Turn: handlebar switches (SYS GPIO3/6). Head: toggle (SYS GPIO7). Brake: OR logic → SYS GPIO21
DC-DC converter       ──► SYS CAN 0x012 enable=1 → 12V rail on
SYS → CAN 0x110       ──► MTR receives mode=Manual → ADC pass-through
\end{verbatim}

\subsection{4.2 Auto mode}\label{auto-mode}

\begin{verbatim}
Jetson AckermannControlCommand ──► High CAN 0x300 ──► RT kinematics
                                          │
               ┌──────────────────────────┤
               ▼ (low CAN)                ▼ (low CAN)
   0x204 {speed, gear} → MTR       0x169 {angle} → EPS-C
               │                          │
               ├──► MCP4725 → Motor       │  RT listens 0x201 for feedback
               ├──► Relays → ECU gear     │  Dynamic angle clamp + following error
               └──► CAN 0x120, 0x206 →    │
                     RT (fwd to Host),     │
                     SYS (EGAS L2 monitor) │

SYS → CAN 0x110 ──► MTR receives mode=Auto → follows 0x204
Jetson ──► High CAN 0x301 ──► RT brake arbitration → Low CAN 0x7B9 → SEB (RT sends directly in AUTO)
Jetson ──► High CAN 0x302 ──► RT fwd → Low CAN 0x302 → SYS → light relays
SYS ────► Low CAN 0x7B9 → SEB (MANUAL/ESTOP only; RT sends in AUTO)
\end{verbatim}

\begin{center}\rule{0.5\linewidth}{0.5pt}\end{center}

\section{5. Responsibility split}\label{responsibility-split}

\begin{longtable}[]{@{}lllll@{}}
\toprule\noalign{}
Concern & Jetson & RT & SYS & MTR \\
\midrule\noalign{}
\endhead
\bottomrule\noalign{}
\endlastfoot
Perception / planning & ✓ & & & \\
ROS 2 → CAN bridge (\texttt{AckermannControlCommand}) & ✓ & & & \\
Autoware.Auto bridge (\texttt{autoware\_vehicle\_bridge}) & ✓ & & & \\
CAN gateway (low ↔ high) & & ✓ & & \\
Tricycle kinematics & & ✓ & & \\
Steering angle compute + CAN TX (\texttt{0x169}) & & ✓ & & \\
Steering boot sync (Listen-Before-Speaking) & & ✓ & & \\
Steering safety: dynamic angle clamp, hard-stops, following error & & ✓ & & \\
Obstacle speed limit & & ✓ & & \\
Command staleness watchdog & & ✓ & & \\
E-stop GPIO + button (wired to both) & & & ✓ & ✓ \\
Brake lever → CAN (\texttt{0x7B9}, 50 Hz continuous) & & & ✓ & \\
Brake boot sync (Listen-Before-Speaking) & & & ✓ & \\
Brake rolling counter + checksum & & & ✓ & \\
DC-DC converter CAN control (\texttt{0x012}) & & & ✓ & \\
Heartbeat monitoring & & ✓ (Jetson high + SYS low) & ✓ (RT, low) & \\
Mode switch reading & & & ✓ & \\
Throttle ADC read (0--5V) & & & & ✓ \\
Throttle MCP4725 DAC output (0--5V) & & & & ✓ \\
Gear 72V read (TLP281 opto) & & & & ✓ \\
Gear 72V output (relay module) & & & & ✓ \\
Motor feedback CAN TX (\texttt{0x206}) & & & & ✓ \\
12V accessory power relay & & & ✓ & \\
Mode indicator lights & & & ✓ & \\
Signal lights (turn, brake, head) & & & ✓ & \\
System diagnostics & & & ✓ & \\
\end{longtable}

\begin{center}\rule{0.5\linewidth}{0.5pt}\end{center}

\subsection{5.1 Jetson Orin --- Autoware.Auto ROS 2 Bridge
(v0.0.4-alpha)}\label{jetson-orin-autoware.auto-ros-2-bridge-v0.0.4-alpha}

A new ROS 2 lifecycle node (\texttt{autoware\_vehicle\_bridge}) at
\texttt{jetson/src/autoware\_vehicle\_bridge/} provides Autoware.Auto topic
compatibility. The CAN protocol on both buses is unchanged.

\textbf{Subscriptions:}
\texttt{autoware\_auto\_control\_msgs/msg/AckermannControlCommand} → CAN
\texttt{0x300}/\texttt{0x301}. \texttt{GearCommand},
\texttt{TurnIndicatorsCommand}, \texttt{HazardLightsCommand} → CAN
\texttt{0x302}. \texttt{Engage} → internal state (false suppresses commands).

\textbf{Publications:} \texttt{VelocityReport} ← CAN \texttt{0x120}.
\texttt{SteeringReport} ← CAN \texttt{0x210}. \texttt{GearReport} ← CAN
\texttt{0x210}. \texttt{ControlModeReport} ← CAN \texttt{0x210} +
\texttt{0x011}.

\textbf{Heartbeats:} Sends \texttt{0x7FC} at 2 Hz. Monitors \texttt{0x7FD}
(1500ms timeout → RCLCPP\_WARN\_THROTTLE). RT's own staleness watchdog (500ms)
stops the vehicle independently.

\textbf{Node parameters:} \texttt{wheel\_base:\ 1.5m},
\texttt{max\_speed\_forward:\ 3.0\ m/s},
\texttt{max\_steering\_angle:\ 0.698\ rad},
\texttt{max\_brake\_pressure\_kpa:\ 5000}, \texttt{loop\_rate:\ 100\ Hz},
\texttt{command\_timeout\_ms:\ 500}.

\begin{center}\rule{0.5\linewidth}{0.5pt}\end{center}

\section{6. Design principles}\label{design-principles}

\begin{enumerate}
\def\labelenumi{\arabic{enumi}.}
\tightlist
\item
  \textbf{Queues over shared state.} No mutexes, no semaphores. Thread-safe
  queue pipes.
\item
  \textbf{ESTOP bypasses queues.} Safety task preempts and writes directly to
  actuators.
\item
  \textbf{One CAN ID = one sender per bus.} Every CAN ID has exactly one
  originator on a given bus. Each node's heartbeat uses its own ID
  (\texttt{0x7FD} RT, \texttt{0x7FE} SYS, \texttt{0x7FC} Jetson).
\item
  \textbf{Lower CAN ID = higher bus priority.} Safety IDs (\texttt{0x00X}) win
  arbitration.
\item
  \textbf{All multi-byte CAN fields are big-endian (MSB first)} --- unless
  SYNTREE protocol specifies otherwise (see \texttt{0x169}, \texttt{0x7B9} in
  can-dictionary).
\item
  \textbf{Manual mode is pass-through, not dead.} SYS mirrors physical inputs to
  outputs.
\item
  \textbf{Actuators are standalone CAN modules.} EPS-C, SEB, and DC-DC are
  commanded via CAN.
\item
  \textbf{RT is the only dual-bus node.} No direct Jetson ↔ SYS path.
\item
  \textbf{Listen Before Speaking.} SYNTREE units require receiving status
  feedback before any command is sent. Boot state machines enforce this.
\item
  \textbf{EGAS 3-level safety separation for motor actuation.} The motor
  controller takes raw analog signals (0--5V throttle, 72V gear) with no
  internal intelligence. This is the only actuator without built-in CAN
  monitoring. Per ISO 26262, it is isolated on a dedicated STM32 (MTR) with
  three independent safety levels.
\item
  \textbf{Mode-gated dual control of SYNTREE actuators (Option D).} In AUTO, RT
  commands both EPS-C and SEB directly --- 1-hop from the kinematics engine. In
  MANUAL, SYS commands SEB based on lever position; EPS-C runs standalone. The
  \texttt{0x7B9} SEB command is dual-sender but mode-gated --- only one node
  transmits at a time, no collision. Per ISO 26262-5:2018 §7.4.4, this is
  redundant-controller practice. No single MCU failure can take both actuators
  offline.
\end{enumerate}

\begin{center}\rule{0.5\linewidth}{0.5pt}\end{center}

\subsection{6.1 EGAS 3-Level Motor Safety
Architecture}\label{egas-3-level-motor-safety-architecture}

The separation of \textbf{MTR (STM32)} from \textbf{SYS (ESP32-S3)} follows the
EGAS 3-level electronic throttle monitoring concept (ISO 26262 ASIL-C):

\begin{verbatim}
Level 3: Hardware — ESTOP button wired direct to both MTR and SYS
         TPS3850 external watchdog on each MCU. No software, no CAN.
         ESTOP press → MTR cuts throttle + gear instantly (local).
         
Level 2: Function Monitor — SYS ESP32-S3
         Monitors MTR via CAN: compares 0x204 setpoint vs 0x206 feedback.
         Mismatch > threshold → CAN 0x001 ESTOP.
         Also handles QM body functions (lights, DCDC, indicators).
         
Level 1: Function Controller — MTR STM32
         Normal actuation: reads sensors, drives MCP4725 DAC + gear relays.
         MANUAL: pass-through from grip/gear. AUTO: follows CAN 0x204.
         No wireless, no OS, minimal attack surface.
\end{verbatim}

\begin{longtable}[]{@{}
  >{\raggedright\arraybackslash}p{(\columnwidth - 2\tabcolsep) * \real{0.4231}}
  >{\raggedright\arraybackslash}p{(\columnwidth - 2\tabcolsep) * \real{0.5769}}@{}}
\toprule\noalign{}
\begin{minipage}[b]{\linewidth}\raggedright
Principle
\end{minipage} & \begin{minipage}[b]{\linewidth}\raggedright
Implementation
\end{minipage} \\
\midrule\noalign{}
\endhead
\bottomrule\noalign{}
\endlastfoot
\textbf{Freedom from interference} & Separate MCUs --- a SYS crash/hang cannot
block motor kill \\
\textbf{ASIL decomposition} & MTR (QM level sensor reads) + SYS monitor (ASIL-B
comparison) → combined ASIL-C \\
\textbf{Independent safe state path} & ESTOP button wired to both MCUs --- MTR
cuts throttle locally, zero CAN delay \\
\textbf{Diverse monitoring} & SYS compares commanded speed vs actual from 0x206
--- mismatch → ESTOP \\
\textbf{Why only motor needed MTR} & SYNTREE EPS-C and SEB already do EGAS Level
1 internally (CAN, PID, angle sensor, fault detection). Motor controller is a
dumb analog device --- it has no intelligence, no CAN, no internal monitoring.
That gap is filled by the dedicated MTR board. \\
\textbf{Body control is QM} & Lights, indicators, DCDC, mode bulbs on SYS are
non-safety --- safe to share MCU with Level 2 monitoring per ISO 26262 QM
classification \\
\end{longtable}

\subsection{6.2 Mode-Gated Dual Control (Option D) --- SYNTREE
Actuators}\label{mode-gated-dual-control-option-d-syntree-actuators}

The EPS-C and SEB are commanded by different nodes depending on mode. Four
options were evaluated (see §6.3). Option D was selected:

\textbf{Why Option D over A (distributed fixed):} In AUTO, RT computes steering
angle AND brake pressure from the same kinematics model. Option D sends both
from RT --- a single control tick produces both CAN frames. No cross-node sync
needed for obstacle response. In A, brake required RT→0x205→SYS→0x7B9 (3 hops);
in D it's RT→0x7B9 (1 hop).

\textbf{Why Option D over C (SYS owns both):} C makes SYS a single point of
failure for both actuators. D preserves independent failure modes --- RT failure
loses AUTO steering/brake but SYS can still brake in MANUAL; SYS failure loses
MANUAL brake but RT can still brake in AUTO. The EPS-C's internal timeout (20ms
comm loss → lock) and SEB's internal timeout provide hardware backup regardless.

\textbf{Why not B (private CAN):} Adds hardware cost (MCP2515 + transceiver +
bus wiring) and SYS bridging latency without safety improvement. The shared
low-level CAN with SYNTREE's own rolling counter + checksum validation already
prevents accidental actuation.

\begin{longtable}[]{@{}
  >{\raggedright\arraybackslash}p{(\columnwidth - 4\tabcolsep) * \real{0.2381}}
  >{\raggedright\arraybackslash}p{(\columnwidth - 4\tabcolsep) * \real{0.2619}}
  >{\raggedright\arraybackslash}p{(\columnwidth - 4\tabcolsep) * \real{0.5000}}@{}}
\toprule\noalign{}
\begin{minipage}[b]{\linewidth}\raggedright
Property
\end{minipage} & \begin{minipage}[b]{\linewidth}\raggedright
RT (AUTO)
\end{minipage} & \begin{minipage}[b]{\linewidth}\raggedright
SYS (MANUAL/ESTOP)
\end{minipage} \\
\midrule\noalign{}
\endhead
\bottomrule\noalign{}
\endlastfoot
EPS-C (0x169) & Sends angle from kinematics & Does NOT send --- EPS-C
standalone \\
SEB (0x7B9) & Sends arbitrated brake kPa→pressure & Sends lever stroke / ESTOP
max \\
0x201 (EPS-C status) & Monitors angle for following error & Does not monitor \\
0x721 (SEB status) & Does not monitor & Monitors stroke for boot sync \\
Mode switch & Receives 0x110 from SYS & Sends 0x110 on change \\
\end{longtable}

\begin{quote}
\textbf{Dual-sender exception}: \texttt{0x7B9} has two senders on the low-level
bus. But only one is active at a time --- mode-gated. In AUTO, RT sends
continuously at 50 Hz and SYS is silent. In MANUAL/ESTOP, SYS sends and RT is
silent. No simultaneous transmission in normal operation. The SEB accepts
whichever frame has a valid checksum and rolling counter. This is standard
automotive redundant-controller practice (ISO 26262-5:2018 §7.4.4).
\end{quote}

\subsection{6.3 Options Comparison
(Archived)}\label{options-comparison-archived}

\begin{longtable}[]{@{}
  >{\raggedright\arraybackslash}p{(\columnwidth - 8\tabcolsep) * \real{0.3793}}
  >{\raggedright\arraybackslash}p{(\columnwidth - 8\tabcolsep) * \real{0.1379}}
  >{\raggedright\arraybackslash}p{(\columnwidth - 8\tabcolsep) * \real{0.1724}}
  >{\raggedright\arraybackslash}p{(\columnwidth - 8\tabcolsep) * \real{0.1379}}
  >{\raggedright\arraybackslash}p{(\columnwidth - 8\tabcolsep) * \real{0.1724}}@{}}
\toprule\noalign{}
\begin{minipage}[b]{\linewidth}\raggedright
Criterion
\end{minipage} & \begin{minipage}[b]{\linewidth}\raggedright
A (distributed)
\end{minipage} & \begin{minipage}[b]{\linewidth}\raggedright
B (private CAN)
\end{minipage} & \begin{minipage}[b]{\linewidth}\raggedright
C (SYS only)
\end{minipage} & \begin{minipage}[b]{\linewidth}\raggedright
\textbf{D (mode-gated)}
\end{minipage} \\
\midrule\noalign{}
\endhead
\bottomrule\noalign{}
\endlastfoot
Single point takes both actuators? & No & Yes & Yes & \textbf{No} \\
AUTO brake latency & \textasciitilde30ms & \textasciitilde40ms &
\textasciitilde20ms & \textbf{\textasciitilde20ms} \\
AUTO steer latency & \textasciitilde20ms & \textasciitilde30ms &
\textasciitilde30ms & \textbf{\textasciitilde20ms} \\
SYS failure → brake lost? & Yes & Yes & Yes & \textbf{No (RT takes over)} \\
RT failure → steer lost? & Yes & Same & No (SYS bridges) & Yes (EPS-C timeout
backup) \\
Code complexity (1-5) & 3 & 4 & 4 & \textbf{3} \\
Hardware cost (1-5) & 4 & 2 & 4 & \textbf{4} \\
Safety score (1-5) & 4 & 2 & 2 & \textbf{5} \\
\textbf{Total (higher=better)} & 23 & 13 & 15 & \textbf{26} \\
\end{longtable}

\begin{center}\rule{0.5\linewidth}{0.5pt}\end{center}

\section{7. RT ESP32-S3 --- Realtime Physics, Steering \& CAN
Gateway}\label{rt-esp32-s3-realtime-physics-steering-can-gateway}

\subsection{7.1 Role}\label{role}

Converts ROS 2 motion commands (high CAN \texttt{0x300}) into: - \textbf{Speed +
gear} → low CAN \texttt{0x204} → SYS - \textbf{Steering angle} → low CAN
\texttt{0x169} → SYNTREE EPS-C

Bridges selected CAN messages (§2.3). Listens to \texttt{0x201\ SES\_STATUS} for
steering feedback and safety monitoring.

\textbf{8 FreeRTOS tasks} on ESP32-S3 @ 240 MHz, 1000 Hz tick.

\subsection{7.2 Dual CAN hardware}\label{dual-can-hardware}

\begin{longtable}[]{@{}
  >{\raggedright\arraybackslash}p{(\columnwidth - 8\tabcolsep) * \real{0.1087}}
  >{\raggedright\arraybackslash}p{(\columnwidth - 8\tabcolsep) * \real{0.2391}}
  >{\raggedright\arraybackslash}p{(\columnwidth - 8\tabcolsep) * \real{0.2391}}
  >{\raggedright\arraybackslash}p{(\columnwidth - 8\tabcolsep) * \real{0.1304}}
  >{\raggedright\arraybackslash}p{(\columnwidth - 8\tabcolsep) * \real{0.2826}}@{}}
\toprule\noalign{}
\begin{minipage}[b]{\linewidth}\raggedright
Bus
\end{minipage} & \begin{minipage}[b]{\linewidth}\raggedright
Controller
\end{minipage} & \begin{minipage}[b]{\linewidth}\raggedright
Interface
\end{minipage} & \begin{minipage}[b]{\linewidth}\raggedright
GPIO
\end{minipage} & \begin{minipage}[b]{\linewidth}\raggedright
Transceiver
\end{minipage} \\
\midrule\noalign{}
\endhead
\bottomrule\noalign{}
\endlastfoot
Low-level & Built-in TWAI & Direct & TX=5, RX=4 & SN65HVD230 \\
High-level & MCP2515 & SPI & SCK=36, MOSI=37, MISO=38, CS=39, INT=40 &
SN65HVD230 \\
\end{longtable}

\subsection{7.3 CAN messages received}\label{can-messages-received}

\begin{longtable}[]{@{}
  >{\raggedright\arraybackslash}p{(\columnwidth - 8\tabcolsep) * \real{0.1515}}
  >{\raggedright\arraybackslash}p{(\columnwidth - 8\tabcolsep) * \real{0.1515}}
  >{\raggedright\arraybackslash}p{(\columnwidth - 8\tabcolsep) * \real{0.1818}}
  >{\raggedright\arraybackslash}p{(\columnwidth - 8\tabcolsep) * \real{0.2727}}
  >{\raggedright\arraybackslash}p{(\columnwidth - 8\tabcolsep) * \real{0.2424}}@{}}
\toprule\noalign{}
\begin{minipage}[b]{\linewidth}\raggedright
Bus
\end{minipage} & \begin{minipage}[b]{\linewidth}\raggedright
ID
\end{minipage} & \begin{minipage}[b]{\linewidth}\raggedright
Name
\end{minipage} & \begin{minipage}[b]{\linewidth}\raggedright
Payload
\end{minipage} & \begin{minipage}[b]{\linewidth}\raggedright
Action
\end{minipage} \\
\midrule\noalign{}
\endhead
\bottomrule\noalign{}
\endlastfoot
Low & \texttt{0x001} & SAFETY\_ESTOP & --- & \texttt{mode\_set(Estop)}, forward
to high \\
Low & \texttt{0x011} & SYS\_SAFETY\_STS & \texttt{\{u8\ estop,\ u8\ hb\_ok\}} &
Forward to high \\
Low & \texttt{0x110} & SYS\_MODE\_CMD & \texttt{u8\ mode} &
\texttt{mode\_set(Manual/Auto)} \\
Low & \texttt{0x120} & SYS\_THROTTLE\_STS & \texttt{i16\ speed\_mmps} & Forward
to high \\
Low & \texttt{0x201} & SES\_STATUS &
\texttt{\{u8\ status,\ i16\ angle,\ u8\ torq,\ …\}} (8 bytes) & Steering
feedback: sync boot angle, following error check \\
Low & \texttt{0x202} & SES\_ErrInfo & 25 fault flags (8 bytes) & Log faults;
escalate any L3 flag to ESTOP via \texttt{0x001} \\
Low & \texttt{0x203} & SES\_Version & \texttt{\{u8\ sw\_ver,\ u8\ hw\_ver\}} &
Log on boot for compatibility check \\
Low & \texttt{0x206} & MTR\_MOTOR\_FBK &
\texttt{\{i16\ actual\_speed,\ u8\ gear\_state,\ u8\ fault\_flags\}} & Forward
to high; monitor motor feedback \\
Low & \texttt{0x600} & SYS\_DIAG\_RPT & 8 bytes & Forward to high \\
Low & \texttt{0x6FA} & SES\_Test &
\texttt{\{i16\ mtr\_curr,\ u16\ ecu\_temp,\ u16\ pow\_volt\}} & Monitor motor
current / ECU temp for degradation \\
Low & \texttt{0x7FE} & SYS\_HEARTBEAT & \texttt{u8\ alive\_ctr} & Feed SYS alive
counter (10 Hz); if no frame for \textgreater200ms → zero setpoints, RT takes
over brake \\
High & \texttt{0x001} & SAFETY\_ESTOP & --- & \texttt{mode\_set(Estop)}, forward
to low \\
High & \texttt{0x300} & HOST\_DRIVE\_CMD & \texttt{\{i32\ speed,\ i32\ yaw\}} &
→ \texttt{cmd\_queue} \\
High & \texttt{0x301} & HOST\_BRAKE\_REQ & \texttt{i32\ pressure\_kpa} & →
atomic store \\
High & \texttt{0x302} & HOST\_LIGHT\_CMD & \texttt{u8} bitfield & Forward to
low \\
High & \texttt{0x7FC} & JETSON\_HEARTBEAT & \texttt{u8\ alive\_ctr} & Feed
Jetson alive counter; frozen \textgreater1500ms → stale command \\
\end{longtable}

\subsection{7.4 CAN messages sent}\label{can-messages-sent}

\begin{longtable}[]{@{}
  >{\raggedright\arraybackslash}p{(\columnwidth - 8\tabcolsep) * \real{0.1613}}
  >{\raggedright\arraybackslash}p{(\columnwidth - 8\tabcolsep) * \real{0.1613}}
  >{\raggedright\arraybackslash}p{(\columnwidth - 8\tabcolsep) * \real{0.1935}}
  >{\raggedright\arraybackslash}p{(\columnwidth - 8\tabcolsep) * \real{0.2903}}
  >{\raggedright\arraybackslash}p{(\columnwidth - 8\tabcolsep) * \real{0.1935}}@{}}
\toprule\noalign{}
\begin{minipage}[b]{\linewidth}\raggedright
Bus
\end{minipage} & \begin{minipage}[b]{\linewidth}\raggedright
ID
\end{minipage} & \begin{minipage}[b]{\linewidth}\raggedright
Name
\end{minipage} & \begin{minipage}[b]{\linewidth}\raggedright
Payload
\end{minipage} & \begin{minipage}[b]{\linewidth}\raggedright
Rate
\end{minipage} \\
\midrule\noalign{}
\endhead
\bottomrule\noalign{}
\endlastfoot
Low & \texttt{0x001} & SAFETY\_ESTOP & --- & Event \\
Low & \texttt{0x204} & RT\_DRIVE\_CMD & \texttt{\{i32\ speed,\ u8\ gear\}} & 100
Hz \\
Low & \texttt{0x205} & RT\_BRAKE\_CMD & \texttt{i32\ brake\_pressure\_kpa} &
\textbf{50 Hz} \\
Low & \texttt{0x169} & VCU\_SES\_REQ &
\texttt{\{u8\ ctrl,\ i16\ angle,\ u8\ speed,\ u8\ sec,\ u8\ cnt+cksum,\ u8\ cksum\}}
(8 bytes) & \textbf{50 Hz} \\
Low & \texttt{0x302} & HOST\_LIGHT\_CMD (fwd) & \texttt{u8} bitfield & Change \\
Low & \texttt{0x7FD} & RT\_HEARTBEAT & \texttt{u8\ alive\_ctr} & 2 Hz \\
High & \texttt{0x001} & SAFETY\_ESTOP (fwd) & --- & Event \\
High & \texttt{0x011} & SYS\_SAFETY\_STS (fwd) &
\texttt{\{u8\ estop,\ u8\ hb\_ok\}} & 5 Hz \\
High & \texttt{0x120} & SYS\_THROTTLE\_STS (fwd) & \texttt{i16\ speed\_mmps} &
100 Hz \\
High & \texttt{0x210} & RT\_STATE\_RPT &
\texttt{\{u8\ mode,\ u8\ steer\_valid,\ u8\ reversing\}} & 10 Hz \\
High & \texttt{0x220} & RT\_PID\_RPT & RESERVED (inactive until encoders fitted)
& --- \\
High & \texttt{0x400} & RT\_OBSTACLE\_RPT & \texttt{u32\ distance\_mm} & 10
Hz \\
High & \texttt{0x600} & SYS\_DIAG\_RPT (fwd) & 8 bytes & 1 Hz \\
High & \texttt{0x7FD} & RT\_HEARTBEAT & \texttt{u8\ alive\_ctr} & 2 Hz \\
\end{longtable}

\subsection{7.5 Internal data types}\label{internal-data-types}

\begin{verbatim}
enum class Mode : uint8_t { Manual = 0, Auto = 1, Estop = 2 };
enum class Gear : uint8_t { N = 0, D = 1, S = 2, R = 3 };

enum class SteerState : uint8_t {
    STEER_BOOT_WAIT,          // 500ms power-on delay — do NOT transmit
    STEER_LISTEN_SYNC,        // Waiting for 0x201 SES_STATUS (angle + alignment), 5s timeout → FAULT
    STEER_ACTIVE,             // Normal operation — transmit 0x169 at 50 Hz
    ESTOP_RAMP_TO_ZERO,       // Non-obstacle ESTOP: ramp to 0° at 20°/s, continue transmitting
    ESTOP_HOLD_THEN_SILENT,   // Obstacle ESTOP: hold current angle 500ms, then silent-stop
    STEER_FAULT               // Timeout or silent-stop — stop transmitting
};

struct DriveCmd {
    int32_t speed_mmps      = 0;   // [-500, 3000]
    int32_t yaw_rate_mrad_s = 0;   // [-3000, 3000]
};

struct ResolvedSetpoint {
    int32_t motor_speed_mmps = 0;
    int32_t steer_angle_mdeg = 0;  // ±45000, +right (internal; convert to decideg for SYNTREE)
    uint8_t gear             = 0;
    bool    steer_valid      = false;
    bool    reversing        = false;
};
\end{verbatim}

\subsection{7.6 Control mechanisms}\label{control-mechanisms}

\subsubsection{Tricycle kinematics}\label{tricycle-kinematics}

\[\delta = \arctan\left(\frac{L \cdot \omega}{|v|}\right) \quad L = 1500\text{ mm}\]

\begin{verbatim}
physics_resolve(cmd):
  1. Convert mm/s→m/s, mrad/s→rad/s
  2. If |v| > 50 mm/s: δ = atan2(L·ω, |v|), steer_valid = true
     Else: δ = steer_hold · 0.8 (decay), steer_valid = false
  3. Clamp δ to ±steer_limit (dynamic, see below)
  4. Clamp v to [-500, 3000] mm/s
  5. reversing = v < 0
  6. gear: v > 0 → D, v == 0 → N, v < 0 → R
\end{verbatim}

\subsubsection{\texorpdfstring{Steering --- SYNTREE EPS-C via CAN
(\texttt{0x169})}{Steering --- SYNTREE EPS-C via CAN (0x169)}}\label{steering-syntree-eps-c-via-can-0x169}

\textbf{Boot sequence --- ``Listen Before Speaking'':}

\begin{verbatim}
State machine (steer_state_machine_loop):

STEER_BOOT_WAIT:
  - 500ms delay after power-on
  - DO NOT transmit any 0x169 frames
  - → STEER_LISTEN_SYNC

STEER_LISTEN_SYNC:
  - Wait for 0x201 SES_STATUS frame (timeout: 5s)
      - On timeout → STEER_FAULT (rider can retry: short-press START → STEER_LISTEN_SYNC)
  - Extract SES_StrAngle (int16, scale 0.1°/bit → convert to internal mdeg)
  - CRITICAL: Set active_target_angle = current_physical_angle
  - Wait for SES_INF_Angle_Status == 1 (aligned)
  - → STEER_ACTIVE

STEER_ACTIVE:
  - Transmit 0x169 at 50 Hz
  - First frame commands wheels to stay exactly where they are
  - Then follow Jetson targets with dynamic clamp
  - Monitor following error: if |cmd - actual| > threshold for > timeout → ESTOP
      - ESTOP behavior depends on trigger:
          Obstacle-triggered: hold current angle (clamped to dynamic limit for speed),
            silent-stop after 500ms. If current angle exceeds dynamic limit, ramp to
            limit at 20°/s first, then hold. Prevents rollover during cornering + braking.
          Non-obstacle (heartbeat loss, cmd stale, manual button): ramp to 0° at 20°/s,
            continue transmitting 0x169 during ramp, then hold at 0°
      - If following error persists during ESTOP centering ramp (>5° for >1s):
          fall back to silent-stop (linkage likely mechanically jammed)

STEER_FAULT:
  - Stop transmitting 0x169
  - EPS-C will timeout-fault internally
      - Short-press START button → reset to STEER_LISTEN_SYNC (manual retry)
      - Long-press START (3s) + throttle at zero → force-activate with target=0°
        (MANUAL mode only; AUTO bulb blinks 2 Hz = degraded steering, AUTO mode locked out)
\end{verbatim}

\textbf{Unit conversion} (internal mdeg ↔ SYNTREE decideg):

\begin{verbatim}
SYNTREE raw = internal_angle_mdeg / 100   (45500 mdeg → 455 raw → 45.5°)
internal_mdeg = SYNTREE raw * 100         (455 raw → 45500 mdeg)
\end{verbatim}

\textbf{SYNTREE protocol specifics:}

\begin{longtable}[]{@{}
  >{\raggedright\arraybackslash}p{(\columnwidth - 2\tabcolsep) * \real{0.6111}}
  >{\raggedright\arraybackslash}p{(\columnwidth - 2\tabcolsep) * \real{0.3889}}@{}}
\toprule\noalign{}
\begin{minipage}[b]{\linewidth}\raggedright
Parameter
\end{minipage} & \begin{minipage}[b]{\linewidth}\raggedright
Value
\end{minipage} \\
\midrule\noalign{}
\endhead
\bottomrule\noalign{}
\endlastfoot
Command ID & \texttt{0x169} (factory default --- SYNTREE preprogrammed, not
reconfigurable. \texttt{0x7B9} is sent by RT in AUTO, SYS in MANUAL/ESTOP ---
one sender active at a time, no collision. \texttt{0x205} removed (RT now sends
SEB directly)) \\
Rate & 50 Hz (20 ms period) --- continuous transmission required \\
Control mode & 1 = Angle Mode \\
Angle range & ±700 raw (±70.0°, unit limit; software clamp tighter) \\
Angle resolution & 0.1°/bit (int16) \\
Slew rate & \texttt{VCU\_SES\_Tgt\_StrAngleSpd} {[}°/s{]} --- speed-dependent \\
Rolling counter & 4-bit, increment 0→15 every frame \\
Checksum & XOR of bytes 0--6, then \texttt{\^{}\ 0xFF} (verify against spec) \\
Security enables & Byte 5: \texttt{roll\_cnt\_enable=1},
\texttt{checksum\_enable=1} \\
\end{longtable}

\textbf{Safety mechanisms:}

\begin{longtable}[]{@{}
  >{\raggedright\arraybackslash}p{(\columnwidth - 2\tabcolsep) * \real{0.4583}}
  >{\raggedright\arraybackslash}p{(\columnwidth - 2\tabcolsep) * \real{0.5417}}@{}}
\toprule\noalign{}
\begin{minipage}[b]{\linewidth}\raggedright
Mechanism
\end{minipage} & \begin{minipage}[b]{\linewidth}\raggedright
Description
\end{minipage} \\
\midrule\noalign{}
\endhead
\bottomrule\noalign{}
\endlastfoot
\textbf{Software hard-stops} & Clamp commanded angle to ±40° (inside physical
end-stops). Reject any Jetson command exceeding this regardless of unit's ±78°
capability. \\
\textbf{Dynamic angle clamp} & Max allowable angle inversely proportional to
\texttt{RT\_PidMeasured} speed. At 25 km/h → max \textasciitilde5°. At 2 km/h →
max \textasciitilde40°. Prevents rollover. \\
\textbf{Following error} & Compare commanded angle (\texttt{0x169}) vs actual
(\texttt{SES\_StrAngle} from \texttt{0x201}). If abs(error) \textgreater{} 5°
for \textgreater{} 300 ms → trigger ESTOP (stuck linkage / rock jam). \\
\textbf{Timeout fault} & If \texttt{0x169} stops for \textgreater20 ms, EPS-C
triggers internal comm fault. RT must maintain 50 Hz in AUTO. \\
\textbf{Alignment check} & \texttt{SES\_INF\_Angle\_Status} must be 1 before
AUTO mode engages. Drive motor locked out until aligned. \\
\end{longtable}

\textbf{Mode behavior:}

\begin{longtable}[]{@{}
  >{\raggedright\arraybackslash}p{(\columnwidth - 2\tabcolsep) * \real{0.2500}}
  >{\raggedright\arraybackslash}p{(\columnwidth - 2\tabcolsep) * \real{0.7500}}@{}}
\toprule\noalign{}
\begin{minipage}[b]{\linewidth}\raggedright
Mode
\end{minipage} & \begin{minipage}[b]{\linewidth}\raggedright
Steering behavior
\end{minipage} \\
\midrule\noalign{}
\endhead
\bottomrule\noalign{}
\endlastfoot
MANUAL & RT does NOT send \texttt{0x169}. EPS-C standalone. RT still listens
\texttt{0x201} for telemetry. \\
AUTO & RT sends \texttt{0x169} at 50 Hz with resolved angle, dynamic clamp +
slew rate applied. \\
ESTOP & Obstacle: hold current angle clamped to dynamic limit (ramp to limit if
exceeding), silent-stop after 500ms. Non-obstacle: ramp to 0° at 20°/s via
active \texttt{0x169}. Fall back to silent-stop if following error persists
\textgreater1s. Both paths non-interruptible by START button (steering
transition defers until ramp/hold complete). \\
\end{longtable}

\subsubsection{Speed control --- open-loop today,
PID-ready}\label{speed-control-open-loop-today-pid-ready}

Motor speed control is \textbf{open-loop}: \texttt{0x204} desired speed → MTR
MCP4725 DAC = \texttt{speed/3000\ ×\ 4095} (fixed voltage mapping). No feedback
compensation. Acceptable for flat-ground steady-state operation but will not
compensate for hills, headwinds, or load changes.

\textbf{Why open-loop?} PID gains are defined in config.h (Kp=1.0, Ki=0.1,
Kd=0.05) but the implementation is deferred. The rear motor encoder hasn't been
physically fitted to the trike yet. The PCNT encoder inputs (GPIO1/2) currently
read zero. Running PID against \texttt{measured=0} would command full throttle
--- it MUST NOT be active until the encoder is wired.

\textbf{Integration plan (gap \#5):} Once the rear motor encoder is fitted: 1.
\texttt{control\_task} on RT feeds
\texttt{speed\_pid\_update(desired,\ measured,\ dt)} with real encoder data 2.
PID output (effort correction) is sent to SYS via an added field in
\texttt{0x204} or a new CAN ID 3. MTR adds the PID trim to the base MCP4725
voltage --- closing the loop 4. Until then, the PID runs as a \textbf{shadow
controller} in RT, outputting to \texttt{0x220\ RT\_PID\_RPT} (telemetry only)
for validation against the open-loop command --- the two should track closely on
flat ground and diverge under load, telling us the PID gains are reasonable
before wiring it in.

\subsubsection{Obstacle speed limit}\label{obstacle-speed-limit}

Jetson perception (LiDAR/camera/stereo) provides minimum obstacle distance via
CAN \texttt{0x400} at 10 Hz. RT applies linear clamp: 300 mm→0, 3000 mm→full
speed. Jetson's own MPPI planner produces safe speed targets; RT's limiter is a
safety backstop, not the primary avoidance mechanism.

\subsubsection{Command staleness watchdog}\label{command-staleness-watchdog}

500 ms, checked @ 10 Hz. Zero \texttt{0x204} + stop \texttt{0x169}.

\subsubsection{Brake arbitration
(max-select)}\label{brake-arbitration-max-select}

\begin{verbatim}
brake_kpa = max(rt_computed_obstacle, jetson_request_0x301)
send 0x205 {brake_kpa} at 50 Hz on low bus → SYS brake_task
\end{verbatim}

RT computes obstacle-emergency brake pressure from distance sensor. Jetson sends
deceleration requests via \texttt{0x301}. The worse (higher) pressure wins.
Result goes to SYS via \texttt{0x205\ RT\_BRAKE\_CMD}.

SYS forwards to SEB in Pressure Mode when \texttt{0x205\ \textgreater{}\ 0},
falling back to Stroke Mode for lever/ESTOP binary triggers.

\subsection{7.7 RTOS task layout}\label{rtos-task-layout}

\begin{verbatim}
Pri 5  can_rx_low   ── TWAI → can_rx_low_queue (16)
      can_rx_high  ── MCP2515 SPI → can_rx_high_queue (16)

Pri 4  dispatch     ◀── both RX queues
           Routes: high 0x300→cmd_queue, 0x301→atomic, 0x302→gw_tx_low
                   low 0x011→gw_tx_high, 0x120→gw_tx_high, 0x600→gw_tx_high
                   low 0x201→steer_feedback (sync angle, following error)
                   any 0x001→mode_set(Estop)+gateway, low 0x110→mode_set

Pri 4  control      ◀── cmd_queue (4, overwrite)
           100 Hz: kinematics  + dynamic angle clamp + obstacle + brake → setpoint_queue

Pri 3  can_tx_low   ◀── setpoint_queue + gw_tx_low_queue
           → 0x204 (100 Hz), 0x205 (50 Hz, drive setpoint), 0x169 (50 Hz, steer state machine), 0x302 (change)
      can_tx_high  ◀── telemetry + gw_tx_high_queue → 0x011,0x120,0x210,0x220,0x400,0x600

Pri 1  watchdog     ── 10 Hz staleness check
      heartbeat    ── 2 Hz 0x7FD on both buses
\end{verbatim}

\begin{longtable}[]{@{}
  >{\raggedright\arraybackslash}p{(\columnwidth - 8\tabcolsep) * \real{0.1622}}
  >{\raggedright\arraybackslash}p{(\columnwidth - 8\tabcolsep) * \real{0.1622}}
  >{\raggedright\arraybackslash}p{(\columnwidth - 8\tabcolsep) * \real{0.1892}}
  >{\raggedright\arraybackslash}p{(\columnwidth - 8\tabcolsep) * \real{0.2162}}
  >{\raggedright\arraybackslash}p{(\columnwidth - 8\tabcolsep) * \real{0.2703}}@{}}
\toprule\noalign{}
\begin{minipage}[b]{\linewidth}\raggedright
Task
\end{minipage} & \begin{minipage}[b]{\linewidth}\raggedright
Prio
\end{minipage} & \begin{minipage}[b]{\linewidth}\raggedright
Stack
\end{minipage} & \begin{minipage}[b]{\linewidth}\raggedright
Period
\end{minipage} & \begin{minipage}[b]{\linewidth}\raggedright
Behavior
\end{minipage} \\
\midrule\noalign{}
\endhead
\bottomrule\noalign{}
\endlastfoot
\texttt{can\_rx\_low} & 5 & 4096 B & Event & \texttt{twai\_receive()} → queue \\
\texttt{can\_rx\_high} & 5 & 4096 B & Event & MCP2515 SPI → queue \\
\texttt{dispatch} & 4 & 4096 B & Event & Route both RX queues + gateway + steer
feedback \\
\texttt{control} & 4 & 4096 B & 100 Hz & Kinematics, dynamic angle clamp,
obstacle, brake arbitration, gear derivation \\
\texttt{can\_tx\_low} & 3 & 3072 B & Event & 0x204@100Hz, 0x169@50Hz (steer SM
gated), 0x302 \\
\texttt{can\_tx\_high} & 3 & 3072 B & Event & Telemetry → MCP2515 SPI \\
\texttt{watchdog} & 1 & 2048 B & 10 Hz & Staleness → zero setpoints + stop
steer \\
\texttt{heartbeat} & 1 & 2048 B & 2 Hz & \texttt{0x7FD} on both buses (per-bus,
not bridged): \texttt{alive\_ctr++\ \&\ 0xFF}, DLC=1 \\
\end{longtable}

\subsection{7.8 Hardware pin assignments --- Board: RT
ESP32-S3}\label{hardware-pin-assignments-board-rt-esp32-s3}

\begin{quote}
⚠️ \textbf{Board identity:} This is the \textbf{RT} ESP32-S3 (realtime physics,
steering, CAN gateway). Do not confuse with SYS ESP32-S3 pin assignments in
§8.8. Same GPIO numbers on different boards are different physical pins ---
connect to the board labeled ``RT,'' not the one labeled ``SYS.''
\end{quote}

\begin{longtable}[]{@{}
  >{\raggedright\arraybackslash}p{(\columnwidth - 6\tabcolsep) * \real{0.2500}}
  >{\raggedright\arraybackslash}p{(\columnwidth - 6\tabcolsep) * \real{0.1875}}
  >{\raggedright\arraybackslash}p{(\columnwidth - 6\tabcolsep) * \real{0.3438}}
  >{\raggedright\arraybackslash}p{(\columnwidth - 6\tabcolsep) * \real{0.2188}}@{}}
\toprule\noalign{}
\begin{minipage}[b]{\linewidth}\raggedright
Signal
\end{minipage} & \begin{minipage}[b]{\linewidth}\raggedright
GPIO
\end{minipage} & \begin{minipage}[b]{\linewidth}\raggedright
Direction
\end{minipage} & \begin{minipage}[b]{\linewidth}\raggedright
Notes
\end{minipage} \\
\midrule\noalign{}
\endhead
\bottomrule\noalign{}
\endlastfoot
CAN TX (low) & 5 & Out & SN65HVD230 \\
CAN RX (low) & 4 & In & SN65HVD230 \\
SPI SCK & 36 & Out & MCP2515 (high CAN) \\
SPI MOSI & 37 & Out & MCP2515 \\
SPI MISO & 38 & In & MCP2515 \\
SPI CS & 39 & Out & MCP2515 \\
MCP INT & 40 & In & MCP2515 interrupt \\
Encoder A (rear motor) & 1 & In & Speed feedback (PCNT), quadrature \\
Encoder B (rear motor) & 2 & In & \\
Encoder A (front wheel) & 3 & In & Speed/angle feedback (PCNT), quadrature ---
\textbf{sensor TBD} \\
Encoder B (front wheel) & 6 & In & \\
Encoder A (rear left wheel) & 9 & In & Differential speed feedback (PCNT),
quadrature --- \textbf{sensor TBD} \\
Encoder B (rear left wheel) & 12 & In & \\
Encoder A (rear right wheel) & 13 & In & Differential speed feedback (PCNT),
quadrature --- \textbf{sensor TBD} \\
Encoder B (rear right wheel) & 14 & In & \\
I2C SDA & 10 & I/O & IMU (optional) \\
I2C SCL & 11 & Out & IMU (optional) \\
WDT toggle & 21 & Out & External watchdog IC (TPS3850). Toggled by
\texttt{control\_task} at 100 Hz. \\
\end{longtable}

\subsection{7.9 Configuration constants}\label{configuration-constants}

\begin{verbatim}
namespace rt {
// Vehicle
constexpr float kWheelbaseMM = 1500.0f;
// Steering (SYNTREE EPS-C)
constexpr float kSteerHardLimitDeg = 40.0f;       // software hard-stop inside mechanical limit
constexpr float kSteerFollowingErrMinDeg = 2.0f;   // floor threshold (was fixed 5.0)
constexpr float kSteerFollowingErrFactor = 0.25f;  // × dynamic_limit → threshold
constexpr int   kSteerFollowingErrMs = 300;        // duration before ESTOP
constexpr int   kSteerCmdRateHz = 50;              // SYNTREE requires 20ms period
constexpr int   kSteerBootWaitMs = 500;            // power-on delay before listening
// Dynamic angle clamp (0.0.4): limit_deg = 40.0 − (speed_kmh − 2.0) × (35.0/23.0), clamped [5.0, 40.0]
constexpr float kAngleClampBaseDeg    = 40.0f;   // max at 2 km/h
constexpr float kAngleClampMinDeg     =  5.0f;   // min at ≥25 km/h
constexpr float kAngleClampRangeDeg   = 35.0f;   // base − min
constexpr float kAngleClampSpeedRange = 23.0f;   // 25 − 2 km/h
// Speed limits
constexpr int kMaxSpeedFwdMmps = 3000, kMaxSpeedRevMmps = 500;
constexpr int kLowSpeedThreshMmps = 50;
// PID (deferred — see gap #5; gains TBD once encoders fitted)
// Obstacle
constexpr unsigned kObstacleStopMM = 300, kObstacleClearMM = 3000;
// Timing
constexpr int kControlLoopHz = 100, kCmdStaleTimeoutMs = 500;
constexpr int kHeartbeatIntervalMs = 500;
constexpr int kHeartbeatTimeoutMsSys = 200;       // SYS heartbeat loss (2 missed at 10 Hz). RT takes over 0x7B9.
constexpr int kHeartbeatTimeoutMsJetson = 1500;  // Jetson heartbeat loss (3 missed frames) → assisted stop
// CAN IDs are in shared/can/can_protocol.h (namespace can):
//   kIdRtHeartbeatLow (0x7FD), kIdSysHeartbeat (0x7FE), kIdJetsonHeartbeat (0x7FC),
//   kIdRtBrakeCmd (0x205), kIdSyntreeEpsCmd (0x169), kIdSyntreeSebCmd (0x7B9), etc.
constexpr int kWdtToggleGpio = 21;
// CAN
constexpr int kCanLowBitrateHz = 500000, kCanHighBitrateHz = 500000;
// Encoders (quadrature, PCNT)
constexpr int kEncRearMotorA = 1, kEncRearMotorB = 2;
constexpr int kEncFrontWheelA = 3, kEncFrontWheelB = 6;     // sensor TBD
constexpr int kEncRearLeftA = 9, kEncRearLeftB = 12;        // sensor TBD
constexpr int kEncRearRightA = 13, kEncRearRightB = 14;     // sensor TBD
} // namespace rt
\end{verbatim}

\subsection{7.10 Error handling}\label{error-handling}

\begin{longtable}[]{@{}
  >{\raggedright\arraybackslash}p{(\columnwidth - 4\tabcolsep) * \real{0.3000}}
  >{\raggedright\arraybackslash}p{(\columnwidth - 4\tabcolsep) * \real{0.3667}}
  >{\raggedright\arraybackslash}p{(\columnwidth - 4\tabcolsep) * \real{0.3333}}@{}}
\toprule\noalign{}
\begin{minipage}[b]{\linewidth}\raggedright
Failure
\end{minipage} & \begin{minipage}[b]{\linewidth}\raggedright
Detection
\end{minipage} & \begin{minipage}[b]{\linewidth}\raggedright
Response
\end{minipage} \\
\midrule\noalign{}
\endhead
\bottomrule\noalign{}
\endlastfoot
Low CAN bus-off & TWAI TEC \textgreater{} 255 & Log, auto-recover; ESTOP if
persistent \\
High CAN bus-off & MCP2515 error flags & Log, auto-recover; zero setpoints until
restored \\
Command stale & Watchdog 500 ms & Zero \texttt{0x204} + stop \texttt{0x169} \\
Obstacle timeout & Echo \textgreater{} 30 ms & Distance = UINT32\_MAX \\
Rear motor encoder missing & Speed = 0 & open-loop - no encoder feedback (I-term
saturates) \\
Wheel encoder missing (any) & No pulses for \textgreater1s at known speed & Log
warning; differential odometry degraded. Does NOT trigger ESTOP. \\
Steering CAN TX fail & TWAI TX errors & Log, EPS-C will timeout-fault \\
Steering following error & abs(cmd − actual) \textgreater{} 5° for 300 ms &
\texttt{mode\_set(Estop)} \\
Steering sync timeout & No \texttt{0x201} within 5s of boot & → STEER\_FAULT;
rider can retry via START button short-press \\
Gateway queue full & \texttt{xQueueSend} fail & Drop (except 0x001 --- direct
TX) \\
\end{longtable}

\subsection{7.11 Startup}\label{startup}

\begin{verbatim}
1. can_low_init() → TWAI, low CAN
2. can_high_init() → SPI + MCP2515, high CAN
3. steer SM → STEER_BOOT_WAIT (500ms) → STEER_LISTEN_SYNC (await 0x201) → STEER_ACTIVE
4. watchdog_init() → timestamp
5. heartbeat_init() → alive counter
6. Create queues (6), Create 8 tasks
7. ESP_LOGI("Ready")
\end{verbatim}

\begin{center}\rule{0.5\linewidth}{0.5pt}\end{center}

\section{8. SYS ESP32-S3 --- Safety, Motor Actuation \& Body
Control}\label{sys-esp32-s3-safety-motor-actuation-body-control}

\subsection{8.1 Role}\label{role-1}

Safety (E-stop, brake lever, RT heartbeat), motor actuation (0--5V via MCP4725,
72V gear via relays), brake control (SYNTREE SEB via CAN \texttt{0x7B9}), DC-DC
converter, signal lights, mode indicators, 12V power, diagnostics.

Low-level CAN only. Jetson communication via RT.

\textbf{15 FreeRTOS tasks} on ESP32-S3 @ 240 MHz, 1000 Hz tick.

\subsection{8.2 CAN interface}\label{can-interface}

Built-in TWAI, GPIO 4/5, 500 kbit/s, SN65HVD230.

\subsection{8.3 CAN messages received}\label{can-messages-received-1}

\begin{longtable}[]{@{}
  >{\raggedright\arraybackslash}p{(\columnwidth - 8\tabcolsep) * \real{0.1143}}
  >{\raggedright\arraybackslash}p{(\columnwidth - 8\tabcolsep) * \real{0.1714}}
  >{\raggedright\arraybackslash}p{(\columnwidth - 8\tabcolsep) * \real{0.2571}}
  >{\raggedright\arraybackslash}p{(\columnwidth - 8\tabcolsep) * \real{0.2286}}
  >{\raggedright\arraybackslash}p{(\columnwidth - 8\tabcolsep) * \real{0.2286}}@{}}
\toprule\noalign{}
\begin{minipage}[b]{\linewidth}\raggedright
ID
\end{minipage} & \begin{minipage}[b]{\linewidth}\raggedright
Name
\end{minipage} & \begin{minipage}[b]{\linewidth}\raggedright
Payload
\end{minipage} & \begin{minipage}[b]{\linewidth}\raggedright
Source
\end{minipage} & \begin{minipage}[b]{\linewidth}\raggedright
Action
\end{minipage} \\
\midrule\noalign{}
\endhead
\bottomrule\noalign{}
\endlastfoot
\texttt{0x001} & SAFETY\_ESTOP & --- & RT or any & \texttt{mode\_set(Estop)} \\
\texttt{0x204} & RT\_DRIVE\_CMD & \texttt{\{i32\ speed,\ u8\ gear\}} & RT & →
\texttt{setpoint\_queue}; stale \textgreater200ms → zero speed + N \\
\texttt{0x205} & RT\_BRAKE\_CMD & \texttt{i32\ brake\_pressure\_kpa} & RT & →
\texttt{g\_brake\_pressure\_kpa} atomic; \textgreater0 → SEB Pressure Mode \\
\texttt{0x206} & MTR\_MOTOR\_FBK &
\texttt{\{i16\ actual\_speed,\ u8\ gear\_state,\ u8\ fault\_flags\}} & MTR &
EGAS L2: compare speed setpoint vs actual; mismatch → ESTOP \\
\texttt{0x302} & HOST\_LIGHT\_CMD (fwd) & \texttt{u8} bitfield & RT & →
\texttt{g\_light\_state} \\
\texttt{0x6FB} & SEB\_Test &
\texttt{\{i16\ mtr\_curr,\ u16\ ecu\_temp,\ u16\ pow\_volt\}} & SEB & Monitor
motor current / ECU temp trends for degradation early warning \\
\texttt{0x721} & SEB\_STATUS &
\texttt{\{u8\ status,\ u16\ stroke,\ u16\ angle,\ u8\ press,\ …\}} (8 bytes) &
SEB & Sync boot stroke, brake feedback, error level \\
\texttt{0x731} & SEB\_ErrInfo & 23 fault flags (8 bytes) & SEB & Log faults;
escalate any L3 flag to ESTOP via \texttt{0x001} \\
\texttt{0x741} & SEB\_Version & \texttt{\{u8\ sw\_ver,\ u8\ hw\_ver\}} & SEB &
Log on boot for compatibility check \\
\texttt{0x7FD} & RT\_HEARTBEAT & \texttt{u8\ alive\_ctr} & RT & Feed RT alive
counter; timeout \textgreater1000ms → ESTOP. Faster detection: \texttt{0x204}
staleness at 200ms → zero speed. \\
\end{longtable}

\subsection{8.4 CAN messages sent}\label{can-messages-sent-1}

\begin{longtable}[]{@{}
  >{\raggedright\arraybackslash}p{(\columnwidth - 8\tabcolsep) * \real{0.1250}}
  >{\raggedright\arraybackslash}p{(\columnwidth - 8\tabcolsep) * \real{0.1875}}
  >{\raggedright\arraybackslash}p{(\columnwidth - 8\tabcolsep) * \real{0.2812}}
  >{\raggedright\arraybackslash}p{(\columnwidth - 8\tabcolsep) * \real{0.1875}}
  >{\raggedright\arraybackslash}p{(\columnwidth - 8\tabcolsep) * \real{0.2188}}@{}}
\toprule\noalign{}
\begin{minipage}[b]{\linewidth}\raggedright
ID
\end{minipage} & \begin{minipage}[b]{\linewidth}\raggedright
Name
\end{minipage} & \begin{minipage}[b]{\linewidth}\raggedright
Payload
\end{minipage} & \begin{minipage}[b]{\linewidth}\raggedright
Rate
\end{minipage} & \begin{minipage}[b]{\linewidth}\raggedright
Notes
\end{minipage} \\
\midrule\noalign{}
\endhead
\bottomrule\noalign{}
\endlastfoot
\texttt{0x001} & SAFETY\_ESTOP & --- & Event & → All nodes (ESTOP on L3 fault,
RT heartbeat loss) \\
\texttt{0x011} & SYS\_SAFETY\_STS & \texttt{\{u8\ estop,\ u8\ hb\_ok\}} & 5 Hz &
→ RT (fwd to Jetson) \\
\texttt{0x012} & SYS\_DCDC\_CMD & \texttt{u8\ enable} & Change & → DC-DC
converter \\
\texttt{0x110} & SYS\_MODE\_CMD & \texttt{u8\ mode} & Change & → RT \\
\texttt{0x600} & SYS\_DIAG\_RPT & 8 bytes & 1 Hz & → RT (fwd to Jetson) \\
\texttt{0x7B9} & VCU\_SEB\_REQ &
\texttt{\{u8\ ctrl{[}2{]},\ u16\ stroke,\ u16\ press,\ u8\ sec,\ u8\ cksum\}} (8
bytes) & \textbf{50 Hz} & → SYNTREE SEB \\
\texttt{0x7FE} & SYS\_HEARTBEAT & \texttt{u8\ alive\_ctr} & 10 Hz & → RT \\
\end{longtable}

\subsection{8.5 Internal data types}\label{internal-data-types-1}

\begin{verbatim}
enum class SysMode : uint8_t { Manual = 0, Auto = 1, Estop = 2 };
enum class Gear : uint8_t { N = 0, D = 1, S = 2, R = 3 };

enum class BrakeState : uint8_t {
    BRAKE_BOOT_WAIT,     // 500ms — do NOT transmit
    BRAKE_LISTEN_SYNC,   // Wait for 0x721 SEB_STATUS, read current stroke
    BRAKE_ACTIVE,        // Transmit 0x7B9 at 50 Hz, synced to SEB
    BRAKE_DEGRADED       // LISTEN_SYNC timeout — transmit with lever defaults, no sync
};

struct ActuatorSetpoint {
    int32_t motor_speed_mmps = 0;
    Gear    gear             = Gear::N;
};
\end{verbatim}

\subsection{8.6 Control mechanisms}\label{control-mechanisms-1}

\subsubsection{Throttle --- MCP4725 I2C DAC
(0--5V)}\label{throttle-mcp4725-i2c-dac-05v}

\begin{longtable}[]{@{}ll@{}}
\toprule\noalign{}
Parameter & Value \\
\midrule\noalign{}
\endhead
\bottomrule\noalign{}
\endlastfoot
DAC device & MCP4725, I2C addr 0x60, SDA=GPIO15, SCL=GPIO16 \\
Resolution & 12-bit (0--4095), VCC=5V → 0--5V output (no op-amp) \\
ADC read & ADC1\_CH5, GPIO10, 12-bit, voltage divider 5V→3.3V \\
Dead zone & 200 (raw ADC) \\
Max speed & 3000 mm/s \\
Update rate & 100 Hz \\
\end{longtable}

\begin{longtable}[]{@{}ll@{}}
\toprule\noalign{}
Mode & Behavior \\
\midrule\noalign{}
\endhead
\bottomrule\noalign{}
\endlastfoot
MANUAL & ADC read → MCP4725 write (pass-through) \\
AUTO & \texttt{setpoint.speed} → \texttt{abs(speed)/3000\ ×\ 4095} → MCP4725 \\
ESTOP & MCP4725 = 0 \\
\end{longtable}

Direction via gear lines --- MCP4725 outputs 0--5V proportional to speed
magnitude only.

\subsubsection{Gear --- TLP281 input + relay output
(72V)}\label{gear-tlp281-input-relay-output-72v}

\textbf{Input} (manual sense, galvanic isolation):

\begin{longtable}[]{@{}lll@{}}
\toprule\noalign{}
Signal & GPIO & Conditioning \\
\midrule\noalign{}
\endhead
\bottomrule\noalign{}
\endlastfoot
D sense & 12 & TLP281 optoisolator ch1 (72V→3.3V) \\
S sense & 13 & TLP281 optoisolator ch2 \\
R sense & 14 & TLP281 optoisolator ch3 \\
\end{longtable}

\textbf{Output} (mimic 72V to ECU):

\begin{longtable}[]{@{}lll@{}}
\toprule\noalign{}
Signal & GPIO & Path \\
\midrule\noalign{}
\endhead
\bottomrule\noalign{}
\endlastfoot
D out & 33 & Relay ch1: GPIO→IN, 72V→1A fuse→COM→NO→ECU, TVS SMCJ90CA→GND \\
S out & 34 & Relay ch2: same path \\
R out & 35 & Relay ch3: same path \\
\end{longtable}

\begin{longtable}[]{@{}ll@{}}
\toprule\noalign{}
Mode & Behavior \\
\midrule\noalign{}
\endhead
\bottomrule\noalign{}
\endlastfoot
MANUAL & Read TLP281 → mirror to relays \\
AUTO & Gear from CAN \texttt{0x204} → energize relay \\
ESTOP & All OFF (N) \\
\end{longtable}

\subsubsection{\texorpdfstring{DC-DC converter --- CAN
\texttt{0x012}}{DC-DC converter --- CAN 0x012}}\label{dc-dc-converter-can-0x012}

\begin{longtable}[]{@{}
  >{\raggedright\arraybackslash}p{(\columnwidth - 4\tabcolsep) * \real{0.2075}}
  >{\raggedright\arraybackslash}p{(\columnwidth - 4\tabcolsep) * \real{0.2264}}
  >{\raggedright\arraybackslash}p{(\columnwidth - 4\tabcolsep) * \real{0.5660}}@{}}
\toprule\noalign{}
\begin{minipage}[b]{\linewidth}\raggedright
Condition
\end{minipage} & \begin{minipage}[b]{\linewidth}\raggedright
CAN \texttt{0x012}
\end{minipage} & \begin{minipage}[b]{\linewidth}\raggedright
12V accessory relay (GPIO27)
\end{minipage} \\
\midrule\noalign{}
\endhead
\bottomrule\noalign{}
\endlastfoot
MANUAL or AUTO & \texttt{enable\ =\ 1} & ON (all accessories powered) \\
ESTOP & \textbf{\texttt{enable\ =\ 1}} (maintains 12V for MCUs, CAN
transceivers, and brake light) & OFF (cuts headlight, turn signals, mode
bulbs) \\
\end{longtable}

Sent on state change. \textbf{DC-DC stays ON during ESTOP} --- MCUs need
12V→3.3V to run, CAN transceivers need 5V, and the brake light must illuminate
during ESTOP (fail-visible). The 12V accessory relay (GPIO27) provides the
secondary cut for non-safety loads. The brake light is wired to the always-on
DC-DC output, not through the accessory relay, so it remains powered during
ESTOP.

\begin{quote}
\textbf{Note:} \href{can-dictionary.md}{\texttt{can-dictionary.md}} §0x012 still
shows DC-DC → OFF on ESTOP (line 81). This paragraph documents the deliberate
design change: DC-DC stays ON during ESTOP to maintain MCU power and brake
light. The can-dictionary entry is outdated and should be updated to match this
section.
\end{quote}

\subsubsection{Signal lights}\label{signal-lights}

\begin{longtable}[]{@{}lll@{}}
\toprule\noalign{}
Signal & GPIO & Notes \\
\midrule\noalign{}
\endhead
\bottomrule\noalign{}
\endlastfoot
Left turn & 18 & Blink 500ms on/off while active \\
Right turn & 19 & Blink 500ms on/off while active \\
Brake light & 21 & \textbf{OR of all braking sources} (see below) \\
Headlight & 22 & On/off \\
\end{longtable}

\textbf{Brake light logic --- OR of all braking sources:}

\begin{verbatim}
brake_light_on = safety_brake_lever_pressed()   // GPIO2 — physical lever
              OR (mode == Estop)                // ESTOP — full brake
              OR g_light_state.brake_light;     // Jetson CAN 0x302 — predictive / hazard
              // Future: OR (brake_stroke_mm > 0) — SEB feedback confirms actual braking
\end{verbatim}

All four sources are local to SYS. \texttt{g\_light\_state.brake\_light} from
Jetson is a \textbf{supplemental} trigger --- useful for predictive illumination
(Jetson sees obstacle before pressure builds) or hazard flashing --- but can
never be the \emph{only} trigger. The physical braking state always wins.

\textbf{Manual mode --- handlebar switches:}

\begin{longtable}[]{@{}
  >{\raggedright\arraybackslash}p{(\columnwidth - 6\tabcolsep) * \real{0.2857}}
  >{\raggedright\arraybackslash}p{(\columnwidth - 6\tabcolsep) * \real{0.2143}}
  >{\raggedright\arraybackslash}p{(\columnwidth - 6\tabcolsep) * \real{0.2143}}
  >{\raggedright\arraybackslash}p{(\columnwidth - 6\tabcolsep) * \real{0.2857}}@{}}
\toprule\noalign{}
\begin{minipage}[b]{\linewidth}\raggedright
Switch
\end{minipage} & \begin{minipage}[b]{\linewidth}\raggedright
GPIO
\end{minipage} & \begin{minipage}[b]{\linewidth}\raggedright
Type
\end{minipage} & \begin{minipage}[b]{\linewidth}\raggedright
Action
\end{minipage} \\
\midrule\noalign{}
\endhead
\bottomrule\noalign{}
\endlastfoot
Left turn & 3 & Momentary, active-low, pull-up & Press → left turn blinks (500ms
on/off). Press again → cancel. \\
Right turn & 6 & Momentary, active-low, pull-up & Press → right turn blinks.
Press again → cancel. \\
Headlight & 7 & Toggle, active-low, pull-up & Each press toggles headlight
on/off. \\
\end{longtable}

\textbf{Mode-dependent behavior:}

\begin{longtable}[]{@{}
  >{\raggedright\arraybackslash}p{(\columnwidth - 6\tabcolsep) * \real{0.1395}}
  >{\raggedright\arraybackslash}p{(\columnwidth - 6\tabcolsep) * \real{0.3023}}
  >{\raggedright\arraybackslash}p{(\columnwidth - 6\tabcolsep) * \real{0.2558}}
  >{\raggedright\arraybackslash}p{(\columnwidth - 6\tabcolsep) * \real{0.3023}}@{}}
\toprule\noalign{}
\begin{minipage}[b]{\linewidth}\raggedright
Mode
\end{minipage} & \begin{minipage}[b]{\linewidth}\raggedright
Turn signals
\end{minipage} & \begin{minipage}[b]{\linewidth}\raggedright
Headlight
\end{minipage} & \begin{minipage}[b]{\linewidth}\raggedright
Brake light
\end{minipage} \\
\midrule\noalign{}
\endhead
\bottomrule\noalign{}
\endlastfoot
MANUAL & Handlebar switches → \texttt{lights\_task} & Handlebar switch → GPIO22
& \textbf{OR logic} --- lever + Jetson bit \\
AUTO & \texttt{g\_light\_state} from CAN \texttt{0x302} &
\texttt{g\_light\_state.headlight} & \textbf{OR logic} --- lever + ESTOP +
Jetson bit \\
ESTOP & OFF & OFF & \textbf{ON} (forced, overrides all) \\
\end{longtable}

\textbf{Turn signal blink pattern} (\texttt{lights\_task}): - 500 ms ON, 500 ms
OFF, repeating while \texttt{left\_turn} or \texttt{right\_turn} is active -
Pressing the same switch again cancels (toggle behavior) - Pressing the opposite
switch cancels the current side and starts the new side - Both pressed
simultaneously → hazard flashers (both blink in sync)

\textbf{Headlight toggle} (\texttt{lights\_task}): - Each press of GPIO7 toggles
\texttt{headlight\_on\ =\ !headlight\_on} - AUTO mode: override from
\texttt{g\_light\_state.headlight} via CAN \texttt{0x302}

\subsubsection{Mode switch --- push button
toggle}\label{mode-switch-push-button-toggle}

A momentary push button on GPIO11 (active-low, internal pull-up, debounced).
Each press toggles the mode: MANUAL → AUTO → MANUAL. ESTOP cannot be exited via
the MODE button --- use the START button (GPIO32) or power-cycle.

\begin{verbatim}
mode_task @ 10 Hz:
  // Mode toggle button (GPIO11)
  read GPIO11
  if falling edge (prev_mode==HIGH, now==LOW) and debounce == 0:
      if current mode == MANUAL → mode_set(Auto)
      elif current mode == AUTO  → mode_set(Manual)
      debounce = kDebounceMs / 100

  // Start button (GPIO32) — exit ESTOP
  read GPIO32
  if falling edge (prev_start==HIGH, now==LOW) and debounce == 0:
      if current mode == ESTOP → mode_set(Manual)
      debounce = kDebounceMs / 100

  if debounce > 0: debounce--
  prev_mode = GPIO11; prev_start = GPIO32
\end{verbatim}

\begin{quote}
A toggle switch would require the rider to physically change switch position. A
push button is simpler to operate while riding --- one tap to switch modes.
\end{quote}

\subsubsection{Mode indicator bulbs}\label{mode-indicator-bulbs}

Two relay-driven bulbs (visible in sunlight, not just PCB LEDs):

\begin{longtable}[]{@{}lll@{}}
\toprule\noalign{}
Signal & GPIO & Active for \\
\midrule\noalign{}
\endhead
\bottomrule\noalign{}
\endlastfoot
AUTO bulb & 25 & AUTO mode \\
MANUAL bulb & 26 & MANUAL mode \\
(both OFF) & --- & ESTOP \\
\end{longtable}

\begin{quote}
GPIO → relay coil → bulb. Bulbs are powered from the 12V accessory rail --- they
go dark on ESTOP regardless of MCU state.
\end{quote}

\subsubsection{12V relay --- unchanged}\label{v-relay-unchanged}

GPIO27 (HIGH=ON, ESTOP→OFF).

\subsubsection{Heartbeat --- automotive liveness
supervision}\label{heartbeat-automotive-liveness-supervision}

Each node sends a \textbf{1-byte alive counter} on its own CAN ID at 2 Hz. A
frozen counter = a frozen node, even if the CAN controller is still DMA-ing from
a hardware buffer.

\textbf{Heartbeat IDs (one sender per ID):}

\begin{longtable}[]{@{}llll@{}}
\toprule\noalign{}
ID & Sender & Bus & Receiver \\
\midrule\noalign{}
\endhead
\bottomrule\noalign{}
\endlastfoot
\texttt{0x7FD} & RT & Both (per-bus, NOT bridged) & SYS (low), Jetson (high) \\
\texttt{0x7FE} & SYS & Low only & RT \\
\texttt{0x7FC} & Jetson & High only & RT \\
\end{longtable}

\textbf{Why heartbeats are NOT bridged:}

Each CAN bus is an independent liveness domain. RT sends \texttt{0x7FD}
independently on each bus (same ID, separate alive counters per bus). SYS and
Jetson heartbeats never leave their respective buses. This means every receiver
sees exactly one sender per heartbeat ID --- no ambiguity, no software demuxing
needed.

\textbf{Liveness matrix (who monitors whom):}

\begin{verbatim}
       SYS ──── low CAN ──── RT ──── high CAN ──── Jetson
              0x7FE (SYS)          0x7FC (Jetson)
              0x7FD (RT)           0x7FD (RT)

SYS watches:   RT (0x7FD on low, 1000ms timeout)
RT watches:    SYS (0x7FE on low, 200ms timeout — 2 missed frames at 10 Hz) AND Jetson (0x7FC on high, 1500ms timeout)
Jetson watches: RT (0x7FD on high, 1500ms timeout)

SYS does NOT watch Jetson — RT handles Jetson failure (zero setpoints)
Jetson does NOT watch SYS — RT handles SYS failure (0x001 ESTOP)
\end{verbatim}

\begin{longtable}[]{@{}
  >{\raggedright\arraybackslash}p{(\columnwidth - 6\tabcolsep) * \real{0.1642}}
  >{\raggedright\arraybackslash}p{(\columnwidth - 6\tabcolsep) * \real{0.2687}}
  >{\raggedright\arraybackslash}p{(\columnwidth - 6\tabcolsep) * \real{0.2836}}
  >{\raggedright\arraybackslash}p{(\columnwidth - 6\tabcolsep) * \real{0.2836}}@{}}
\toprule\noalign{}
\begin{minipage}[b]{\linewidth}\raggedright
Parameter
\end{minipage} & \begin{minipage}[b]{\linewidth}\raggedright
RT↔SYS (low bus)
\end{minipage} & \begin{minipage}[b]{\linewidth}\raggedright
RT→Jetson (high)
\end{minipage} & \begin{minipage}[b]{\linewidth}\raggedright
Jetson→RT (high)
\end{minipage} \\
\midrule\noalign{}
\endhead
\bottomrule\noalign{}
\endlastfoot
RT heartbeat ID & \texttt{0x7FD} & \texttt{0x7FD} & --- \\
Peer heartbeat ID & \texttt{0x7FE} (SYS) & --- & \texttt{0x7FC} (Jetson) \\
DLC & 1 & 1 & 1 \\
Payload & \texttt{u8\ alive\_ctr} (0--255, wraps) & \texttt{u8\ alive\_ctr} &
\texttt{u8\ alive\_ctr} \\
Period & RT: 500 ms (2 Hz) / SYS: 100 ms (10 Hz) & 500 ms (2 Hz) & 500 ms (2
Hz) \\
Timeout & \textbf{200ms} (2 missed frames) & 1500ms (3 missed frames) &
\textbf{1500ms} (3 missed frames) \\
Action on loss & \textbf{RT takes over 0x7B9 (stroke=max) + CAN 0x001.} MTR
kills motor/gear locally. Brake gap ≤220ms. & Jetson: \textbf{assisted stop} ---
zero 0x204 + stop 0x169 + 0x205 brake=2000 kPa + SYS→MANUAL. Rider can override.
& RT: zero \texttt{0x204} + stop \texttt{0x169} \\
\end{longtable}

\textbf{Why 200ms for SYS heartbeat (10 Hz):}

SYS controls the SEB brake actuator --- there is no independent fast-path check
for brake command loss (unlike steering which has the 300ms following-error
check, and motor which has the 200ms 0x204 staleness check). The heartbeat IS
the brake fast-path check. At 25 km/h, 200ms is \textasciitilde1.4m of travel
--- within FTTI for brake loss. At 10 Hz (100ms period), 2 missed frames = 200ms
worst case. This is tight enough for brake safety while providing a 100ms
debounce against CAN jitter.

\textbf{RT brake takeover on SYS heartbeat loss:} Architecture §6.2 Option D
explicitly requires RT to take over brake on SYS failure: ``SYS failure → brake
lost? No (RT takes over).'' When RT detects SYS heartbeat loss, RT immediately
begins transmitting 0x7B9 with stroke=max (full brake) at 50 Hz, regardless of
current mode. RT already knows the full SYNTREE SEB protocol --- this is the
same transmission it performs in AUTO mode, with stroke=max substituted for the
computed pressure target. The mode gate opens on emergency: both RT and SYS may
briefly transmit 0x7B9 during the takeover transition, which is within the
dual-sender exception documented in §6.2. Total brake gap: 200ms detection +
20ms first frame = 220ms worst case.

\textbf{Jetson→RT at 1500ms --- assisted stop:}

Jetson runs the perception stack (obstacle detection). When it dies, the vehicle
must decelerate actively --- pure coast is insufficient in traffic. Three missed
frames at 2 Hz = 1500ms protects against false triggers from Jetson's
non-realtime Linux CAN stack (a Linux machine that can't send a single CAN frame
in 1.5s is genuinely dead). On timeout, RT commands: zero 0x204 speed, stop
0x169 steering, 0x205 brake = 2000 kPa (\textasciitilde2 MPa, moderate
deceleration without wheel lockup), and transitions SYS to MANUAL mode. Brake
light illuminates. Rider can override with lever. This is ``assisted stop'' ---
between pure coast and full ESTOP --- providing active deceleration while
preserving lights and rider agency.

\textbf{MTR-side \texttt{0x204} staleness check (fast path):}

MTR's control loop checks data freshness independently of heartbeat. If
\texttt{0x204\ RT\_DRIVE\_CMD} stops arriving in AUTO mode (RT control loop
crashed, CAN TX failed):

\begin{verbatim}
motor_task @ 100Hz:
  if (time_since_last_0x204 > 200ms):  // 2 missed frames at 100 Hz
      setpoint.speed = 0
      setpoint.gear = N
\end{verbatim}

This is a data-quality check, not a node-liveness check. It triggers in 200ms
(vs 1000ms for heartbeat) because stale data is immediately dangerous. Combined
with the heartbeat, SYS now has two independent checks:

\begin{longtable}[]{@{}llll@{}}
\toprule\noalign{}
Check & Detects & Timeout & Action \\
\midrule\noalign{}
\endhead
\bottomrule\noalign{}
\endlastfoot
\texttt{0x204} staleness & RT control loop stopped & 200ms & Zero speed +
neutral \\
RT heartbeat loss (\texttt{0x7FD}) & RT node dead & 1000ms & ESTOP (full
stop) \\
\end{longtable}

\textbf{Startup grace period:}

At boot, heartbeats haven't been established. \texttt{safety\_heartbeat\_ok()}
returns \texttt{true} for the first \textbf{3 seconds} if
\texttt{last\_hb\_timestamp\ ==\ 0}. After the grace period, real heartbeat
checking begins. This prevents false ESTOP during boot.

\begin{verbatim}
bool safety_heartbeat_ok() {
    int64_t now = esp_timer_get_time();
    if (last_hb_rt_us == 0) {
        return (now < kStartupGracePeriodUs);  // 3_000_000
    }
    return ((now - last_hb_rt_us) / 1000) < kHeartbeatTimeoutMs;  // 1000
}
\end{verbatim}

\textbf{Alive counter validation:}

A heartbeat frame with the same counter value as the previous frame = stuck CAN
controller (MCU hung, controller still DMA-ing from buffer). Treated as a missed
heartbeat.

\begin{verbatim}
bool heartbeat_is_fresh(uint8_t new_ctr) {
    if (new_ctr != last_alive_ctr) { last_alive_ctr = new_ctr; return true; }
    return false;  // frozen — same counter twice
}
\end{verbatim}

\subsubsection{External watchdog}\label{external-watchdog}

Each ESP32 toggles a dedicated \textbf{external watchdog GPIO} every iteration
of its highest-priority task. A hardware window watchdog IC (e.g., TPS3850)
resets the MCU if toggling stops for \textgreater100 ms. On reset, all outputs
default to safe state.

\begin{longtable}[]{@{}lllll@{}}
\toprule\noalign{}
Node & GPIO & Toggled by & Period & Watchdog IC \\
\midrule\noalign{}
\endhead
\bottomrule\noalign{}
\endlastfoot
RT & \textbf{21} & \texttt{control\_task} & 100 Hz & TPS3850 or equiv, 100ms
window \\
SYS & \textbf{23} & \texttt{safety\_task} & 20 Hz & TPS3850 or equiv, 100ms
window \\
\end{longtable}

\begin{quote}
This is independent of CAN heartbeat. A hung MCU with a frozen CAN controller is
invisible to heartbeat --- but the external watchdog catches it.
\end{quote}

\subsubsection{Physical controls}\label{physical-controls}

Three buttons on the dashboard:

\begin{longtable}[]{@{}
  >{\raggedright\arraybackslash}p{(\columnwidth - 6\tabcolsep) * \real{0.2857}}
  >{\raggedright\arraybackslash}p{(\columnwidth - 6\tabcolsep) * \real{0.2143}}
  >{\raggedright\arraybackslash}p{(\columnwidth - 6\tabcolsep) * \real{0.2143}}
  >{\raggedright\arraybackslash}p{(\columnwidth - 6\tabcolsep) * \real{0.2857}}@{}}
\toprule\noalign{}
\begin{minipage}[b]{\linewidth}\raggedright
Button
\end{minipage} & \begin{minipage}[b]{\linewidth}\raggedright
GPIO
\end{minipage} & \begin{minipage}[b]{\linewidth}\raggedright
Type
\end{minipage} & \begin{minipage}[b]{\linewidth}\raggedright
Action
\end{minipage} \\
\midrule\noalign{}
\endhead
\bottomrule\noalign{}
\endlastfoot
\textbf{ESTOP} & 1 & Big red mushroom, NC, active-low, hardware ISR & Instant
ESTOP --- motor kill, brake engage, DCDC off \\
\textbf{START} & 32 & Green momentary, active-low, pull-up, debounced & Exit
ESTOP → MANUAL. No effect in AUTO/MANUAL. \\
\textbf{MODE} & 11 & Momentary, active-low, pull-up, debounced & Toggle MANUAL ↔
AUTO. Ignored in ESTOP. \\
\end{longtable}

Plus brake lever on GPIO2 (active-low, pull-up). Safety task polls ESTOP + brake
lever at 20 Hz. Mode task handles MODE + START buttons at 10 Hz.

\begin{quote}
Industrial safety pattern: separate STOP (red mushroom) and START (green)
buttons. STOP is NC (normally-closed) --- a cut wire triggers ESTOP, not a
failure-silent state.
\end{quote}

\subsubsection{\texorpdfstring{Brake --- SYNTREE SEB via CAN
(\texttt{0x7B9})}{Brake --- SYNTREE SEB via CAN (0x7B9)}}\label{brake-syntree-seb-via-can-0x7b9}

\textbf{Boot sequence --- ``Listen Before Speaking'':}

\begin{verbatim}
BRAKE_BOOT_WAIT:
  - 500ms delay after power-on
  - DO NOT transmit any 0x7B9 frames
  - → BRAKE_LISTEN_SYNC

BRAKE_LISTEN_SYNC:
  - Wait for 0x721 SEB_STATUS frame
  - Extract SEB_Stroke_Value (u16, scale 0.05, offset -30)
  - Set initial command target = current stroke (hold position)
  - Wait for SEB_Alignment_Status == 1
  - → BRAKE_ACTIVE

BRAKE_ACTIVE:
  - Transmit 0x7B9 at 50 Hz continuously
  - Rolling counter increments 0→15 every frame
  - Checksum = XOR(bytes 0–6) ^ 0xFF (verify against spec)

BRAKE_DEGRADED:
  - Entered when LISTEN_SYNC times out (no 0x721 within 2s)
  - Transmit 0x7B9 at 50 Hz with conservative defaults:
      Lever pressed (GPIO2 LOW) → stroke = kBrakeManualStroke (~15 mm)
      Lever released → stroke = 0 mm
  - Rolling counter still increments, checksum still computed
  - When first valid 0x721 arrives: sync current stroke, → BRAKE_ACTIVE
  - Diagnostic flag set in 0x600 SYS_DIAG_RPT
  - Brake lever remains functional — this is not a terminal fault state
\end{verbatim}

\textbf{SYNTREE SEB protocol:}

\begin{longtable}[]{@{}
  >{\raggedright\arraybackslash}p{(\columnwidth - 2\tabcolsep) * \real{0.6111}}
  >{\raggedright\arraybackslash}p{(\columnwidth - 2\tabcolsep) * \real{0.3889}}@{}}
\toprule\noalign{}
\begin{minipage}[b]{\linewidth}\raggedright
Parameter
\end{minipage} & \begin{minipage}[b]{\linewidth}\raggedright
Value
\end{minipage} \\
\midrule\noalign{}
\endhead
\bottomrule\noalign{}
\endlastfoot
Command ID & \texttt{0x7B9} \\
Rate & 50 Hz (20 ms) --- continuous transmission required \\
Control mode & 0 = Stroke Mode, 1 = Pressure Mode (1-bit per CSV) \\
Stroke range & -5 to 27 mm (raw: 500--1140, scale 0.05, offset -30) \\
Pressure range & 0 to 5 MPa (raw: 0--100, u8, scale 0.05 MPa/bit, offset 0) \\
Pressure conversion & \texttt{seb\_raw\ =\ kPa\ ×\ 0.02} (1 MPa = 1000 kPa; 1
bit = 0.05 MPa) \\
Rolling counter & 4-bit, increment every frame \\
Checksum & XOR of bytes 0--6, then \texttt{\^{}\ 0xFF} (verify against spec) \\
\end{longtable}

\begin{quote}
\textbf{Pressure Mode (1) --- verified kPa→raw conversion:}

RT sends \texttt{0x205\ \{i32\ brake\_pressure\_kpa\}}. SYS converts using
verified SYNTREE SEB spec: \texttt{seb\_raw\ =\ (uint8\_t)(kpa\ *\ 0.02f)}.
Scale: 0.05 MPa/bit, range 0--5 MPa (raw 0--100). At 5000 kPa (5 MPa) → raw=100.
Clamp to \texttt{kSebMaxPressureRaw} (100).

\textbf{Mode-switching:} 0x205 transitions 0→positive: hold current stroke,
switch mode bit to 1 (Pressure), ramp pressure from current measured to target.
0x205 drops to 0: switch mode to 1 (Stroke), stroke=0. Prevents pressure
transients.
\end{quote}

\textbf{Mode-dependent behavior:}

\begin{longtable}[]{@{}
  >{\raggedright\arraybackslash}p{(\columnwidth - 2\tabcolsep) * \real{0.2857}}
  >{\raggedright\arraybackslash}p{(\columnwidth - 2\tabcolsep) * \real{0.7143}}@{}}
\toprule\noalign{}
\begin{minipage}[b]{\linewidth}\raggedright
Mode
\end{minipage} & \begin{minipage}[b]{\linewidth}\raggedright
Brake behavior
\end{minipage} \\
\midrule\noalign{}
\endhead
\bottomrule\noalign{}
\endlastfoot
MANUAL & Brake lever GPIO2 LOW → stroke = \texttt{kBrakeManualStroke}
(\textasciitilde15 mm). Released → stroke = 0 mm. Transmit at 50 Hz in Stroke
Mode. \\
AUTO & Lever pressed → \texttt{kBrakeManualStroke} (driver override always
wins). No lever + \texttt{0x205\ \textgreater{}\ 0} → Pressure Mode with kPa→MPa
mapping. No lever + \texttt{0x205\ ==\ 0} → stroke = 0. ESTOP → max. \\
ESTOP & Stroke = \texttt{kBrakeMaxStroke} (full brake, \textasciitilde27 mm).
Transmit at 50 Hz in Stroke Mode. \\
\end{longtable}

\textbf{Stroke value calculation:}

\begin{verbatim}
Physical stroke [mm] → raw = (physical + 30.0) / 0.05
Example: 0 mm → (0+30)/0.05 = 600
         15 mm → (15+30)/0.05 = 900
         27 mm → (27+30)/0.05 = 1140
\end{verbatim}

\subsection{8.7 RTOS task layout}\label{rtos-task-layout-1}

\begin{verbatim}
Pri 5  can_rx      ── TWAI → can_rx_queue (16)
       safety      ── GPIO poll @ 20 Hz → ESTOP / HB check

Pri 4  dispatch    ◀── can_rx_queue: 0x204→setpoint, 0x302→light, 0x001→ESTOP, 0x721→brake_feedback
       mode        ── Push button (GPIO11) @ 10 Hz → toggle MANUAL↔AUTO, CAN 0x110
       motor       ◀── setpoint_queue (4, overwrite)
             100 Hz: AUTO→MCP4725+gear, MANUAL→pass-through, ESTOP→all off

Pri 3  throttle    ── ADC @ 100 Hz → CAN 0x120
       gear        ── Gear FSM @ 50 Hz
       brake       ── Brake FSM @ 50 Hz → CAN 0x7B9 (continuous, rolling ctr + cksum)
       lights      ── Light FSM @ 20 Hz (blink timing, ESTOP=brake ON)
       dcdc        ── DCDC FSM @ 5 Hz → CAN 0x012

Pri 2  indicator   ── Mode bulbs @ 5 Hz
       power       ── 12V relay @ 5 Hz
       can_tx      ── Safety status @ 5 Hz → CAN 0x011

Pri 1  diag        ── System health @ 1 Hz → CAN 0x600
       hb          ── 0x7FE @ 10 Hz
\end{verbatim}

\begin{longtable}[]{@{}
  >{\raggedright\arraybackslash}p{(\columnwidth - 8\tabcolsep) * \real{0.1622}}
  >{\raggedright\arraybackslash}p{(\columnwidth - 8\tabcolsep) * \real{0.1622}}
  >{\raggedright\arraybackslash}p{(\columnwidth - 8\tabcolsep) * \real{0.1892}}
  >{\raggedright\arraybackslash}p{(\columnwidth - 8\tabcolsep) * \real{0.2162}}
  >{\raggedright\arraybackslash}p{(\columnwidth - 8\tabcolsep) * \real{0.2703}}@{}}
\toprule\noalign{}
\begin{minipage}[b]{\linewidth}\raggedright
Task
\end{minipage} & \begin{minipage}[b]{\linewidth}\raggedright
Prio
\end{minipage} & \begin{minipage}[b]{\linewidth}\raggedright
Stack
\end{minipage} & \begin{minipage}[b]{\linewidth}\raggedright
Period
\end{minipage} & \begin{minipage}[b]{\linewidth}\raggedright
Behavior
\end{minipage} \\
\midrule\noalign{}
\endhead
\bottomrule\noalign{}
\endlastfoot
\texttt{can\_rx} & 5 & 4096 B & Event & \texttt{twai\_receive()}, copy to
queue \\
\texttt{safety} & 5 & 2048 B & 20 Hz & ESTOP GPIO, RT HB timeout, EGAS L2:
compare 0x204 vs 0x206 \\
\texttt{dispatch} & 4 & 3072 B & Event & Route 0x206, 0x302, 0x001, 0x721 \\
\texttt{mode} & 4 & 2048 B & 10 Hz & MODE btn toggle + START btn (ESTOP→MANUAL),
CAN 0x110 → RT + MTR \\
\texttt{brake} & 3 & 2048 B & \textbf{50 Hz} & Brake SM (boot sync) + CAN 0x7B9
with rolling ctr + checksum \\
\texttt{lights} & 3 & 1536 B & 20 Hz & Light GPIOs + blink \\
\texttt{dcdc} & 3 & 1024 B & 5 Hz & DCDC FSM, CAN 0x012 \\
\texttt{indicator} & 2 & 1024 B & 5 Hz & Mode bulbs (AUTO/MANUAL) \\
\texttt{power} & 2 & 1024 B & 5 Hz & 12V relay \\
\texttt{can\_tx} & 2 & 3072 B & 5 Hz & CAN \texttt{0x011} (mode task sends
\texttt{0x110}) \\
\texttt{diag} & 1 & 2048 B & 1 Hz & CAN 0x600 \\
\texttt{hb} & 1 & 2048 B & 10 Hz & CAN \texttt{0x7FE} (SYS heartbeat to RT, fast
path for brake loss detection) \\
\end{longtable}

\begin{quote}
\textbf{Note:} Motor actuation (MCP4725 DAC), throttle ADC, and gear relay
control run on SYS in MANUAL/AUTO modes. MTR STM32 handles only the bare-metal
motor controller interface (EGAS L1). In ESTOP, SYS zeros the DAC and disengages
all relays.
\end{quote}

\subsection{8.8 Hardware pin assignments --- Board: SYS
ESP32-S3}\label{hardware-pin-assignments-board-sys-esp32-s3}

\begin{quote}
⚠️ \textbf{Board identity:} This is the \textbf{SYS} ESP32-S3 (safety, brake
control, body control). Motor actuation is on the dedicated \textbf{MTR} STM32
board (§5.0). Do not confuse with RT ESP32-S3 pin assignments in §7.8. Same GPIO
numbers on different boards are different physical pins --- connect to the board
labeled ``SYS,'' not the one labeled ``RT.''

\textbf{Shared GPIO numbers (safe --- different chips):} ESTOP button is shared
between SYS GPIO1 and MTR (hardwired to both MCUs for Level 3 kill). No other
signals are split. GPIO 3,6,7 appear in both RT and SYS tables but are
physically separate pins on separate ESP32-S3 packages --- no electrical
conflict. Throttle ADC, gear sense, MCP4725 I2C, and gear outputs are on MTR
only (§5.0).
\end{quote}

\begin{longtable}[]{@{}
  >{\raggedright\arraybackslash}p{(\columnwidth - 6\tabcolsep) * \real{0.2105}}
  >{\raggedright\arraybackslash}p{(\columnwidth - 6\tabcolsep) * \real{0.1579}}
  >{\raggedright\arraybackslash}p{(\columnwidth - 6\tabcolsep) * \real{0.2895}}
  >{\raggedright\arraybackslash}p{(\columnwidth - 6\tabcolsep) * \real{0.3421}}@{}}
\toprule\noalign{}
\begin{minipage}[b]{\linewidth}\raggedright
Signal
\end{minipage} & \begin{minipage}[b]{\linewidth}\raggedright
GPIO
\end{minipage} & \begin{minipage}[b]{\linewidth}\raggedright
Direction
\end{minipage} & \begin{minipage}[b]{\linewidth}\raggedright
Conditioning
\end{minipage} \\
\midrule\noalign{}
\endhead
\bottomrule\noalign{}
\endlastfoot
CAN TX (low) & 5 & Out & SN65HVD230 \\
CAN RX (low) & 4 & In & SN65HVD230 \\
E-stop button & 1 & In & Big red mushroom, NC (active-low), pull-up, hardware
ISR. Shared with MTR. \\
Brake lever & 2 & In & Active-low, pull-up → CAN 0x7B9 → SEB \\
Start button & \textbf{32} & In & Momentary, active-low, pull-up, debounced.
Exits ESTOP → MANUAL. \\
Mode button & 11 & In & Push button, active-low, pull-up, debounced (momentary
toggle MANUAL↔AUTO). Publishes CAN 0x110 to RT + MTR. \\
Left turn switch & \textbf{3} & In & Handlebar switch, active-low, pull-up \\
Right turn switch & \textbf{6} & In & Handlebar switch, active-low, pull-up \\
Headlight switch & \textbf{7} & In & Handlebar toggle, active-low, pull-up \\
Left turn & 18 & Out & Relay → lamp \\
Right turn & 19 & Out & Relay → lamp \\
Brake light & 21 & Out & Relay → lamp \\
Headlight & 22 & Out & Relay → lamp \\
AUTO bulb & 25 & Out & Relay → bulb (12V rail) \\
MANUAL bulb & 26 & Out & Relay → bulb (12V rail) \\
12V relay & 27 & Out & Secondary cut on ESTOP \\
WDT toggle & \textbf{23} & Out & External watchdog IC (TPS3850). Toggled by
\texttt{safety\_task} every 50 ms. \\
\end{longtable}

\begin{quote}
\textbf{Motor I/O (throttle, gear, MCP4725) is on MTR STM32 only} --- see §5.0
for MTR pin assignments.
\end{quote}

\subsection{8.9 Configuration constants}\label{configuration-constants-1}

\begin{verbatim}
namespace sys {
// CAN
constexpr int kCanBitrateHz = 500000, kCanTxGpio = 5, kCanRxGpio = 4;
// Safety
constexpr int kEstopGpio = 1, kBrakeLeverGpio = 2, kModeSwitchGpio = 11;
constexpr int kWdtToggleGpio = 23;
// Throttle/gear I/O is on MTR (mtr-stm32/src/config.h), not SYS
// Light inputs (handlebar switches, MANUAL mode)
constexpr int kSwitchLeftTurn = 3, kSwitchRightTurn = 6, kSwitchHeadlight = 7;
// Light outputs
constexpr int kLightLeftTurn = 18, kLightRightTurn = 19;
constexpr int kLightBrake = 21, kLightHead = 22;
// Indicators & power
constexpr int kBulbAuto = 25, kBulbManual = 26, kPower12vRelay = 27;
constexpr int kModeBtnGpio = 11, kStartBtnGpio = 32;
constexpr int kDebounceMs = 500;         // push button debounce period
// Turn blink
constexpr int kTurnBlinkOnMs = 500, kTurnBlinkOffMs = 500;
// Timing
constexpr int kControlLoopHz = 100;
constexpr int kHeartbeatIntervalMs    = 100;   // SYS sends 0x7FE at 10 Hz (fast path for brake loss detection)
constexpr int kHeartbeatTimeoutMsRt   = 1000;  // RT heartbeat loss (0x7FD at 2 Hz, 2 missed frames = 1000ms). Faster 0x204 staleness (200ms) catches RT crash first.
// CAN IDs are in shared/can/can_protocol.h (namespace can):
//   kIdRtHeartbeatLow (0x7FD), kIdSysHeartbeat (0x7FE), kIdRtBrakeCmd (0x205),
//   kIdSyntreeSebCmd (0x7B9), etc.
// Shared constants are in shared/shared_config.h (namespace shared):
//   kSebMaxPressureRaw (100), kStartupGracePeriodMs (3000),
//   kBrakeStrokeScale (0.05), kBrakeStrokeOffset (-30.0)
constexpr int kSetpointStaleMs = 200;           // 0x204 staleness → zero speed
constexpr int kSafetyCheckHz = 20, kGearCheckHz = 50;
// Brake (SYNTREE SEB)
constexpr int kBrakeCmdRateHz = 50, kBrakeBootWaitMs = 500;
constexpr float kBrakeManualStroke = 15.0f, kBrakeMaxStroke = 27.0f;
} // namespace sys
\end{verbatim}

\subsection{8.10 Error handling}\label{error-handling-1}

\begin{longtable}[]{@{}
  >{\raggedright\arraybackslash}p{(\columnwidth - 4\tabcolsep) * \real{0.3000}}
  >{\raggedright\arraybackslash}p{(\columnwidth - 4\tabcolsep) * \real{0.3667}}
  >{\raggedright\arraybackslash}p{(\columnwidth - 4\tabcolsep) * \real{0.3333}}@{}}
\toprule\noalign{}
\begin{minipage}[b]{\linewidth}\raggedright
Failure
\end{minipage} & \begin{minipage}[b]{\linewidth}\raggedright
Detection
\end{minipage} & \begin{minipage}[b]{\linewidth}\raggedright
Response
\end{minipage} \\
\midrule\noalign{}
\endhead
\bottomrule\noalign{}
\endlastfoot
E-stop pressed & GPIO1 LOW & ESTOP → DAC=0, gears off, brake=max, \textbf{DCDC
on} (maintains 12V for MCUs, CAN transceivers, brake light), 12V accessory relay
off \\
CAN bus-off & TWAI TEC \textgreater{} 255 & Log, auto-recover \\
RT HB timeout & \texttt{0x7FD} alive ctr frozen \textgreater1000ms & ESTOP (AUTO
only) \\
\texttt{0x204} stale & No \texttt{0x204} for \textgreater200ms & Zero speed + N
gear (controlled stop) \\
Brake lever & GPIO2 LOW & Engage brake \\
ADC fail & \texttt{adc1\_get\_raw()==0} & Throttle = 0 \\
Gear sense conflict & Multiple lines HIGH & Treat as N (fail-safe) \\
DCDC CAN TX fail & TWAI TX errors & 12V relay backup cut \\
SEB sync timeout & No \texttt{0x721} within 2s of boot &
\texttt{BRAKE\_DEGRADED} --- transmits 0x7B9 with lever-based defaults (0mm
released, 15mm pressed). Recovers when 0x721 arrives. Lever always
functional. \\
SEB checksum fail & SEB rejects frame & Frame dropped; counter still
increments \\
SEB stroke following error & abs(0x7B9 cmd − 0x721 actual) \textgreater{} 3mm
for \textgreater100ms & Log brake fault to 0x600 diag. Set persistent fault
flag. Cannot escalate (ESTOP is already max brake) but critical for
post-incident analysis. \\
SEB Level 3 fault & \texttt{SEB\_Error\_Status\ ≥\ 3} in 0x721 or 0x731 & Log
SEB severe fault to 0x600 diag. Alert rider via brake light flash pattern (if
12V available). \\
0x721 staleness & No 0x721 for \textgreater100ms (10 missed at 100 Hz) & SEB
comms lost --- log fault. If sustained, treat as brake system offline. \\
MTR ESTOP ACK timeout & \texttt{ESTOP\_ACTIVE} bit not set in 0x206 fault\_flags
within 100ms of ESTOP & Log MTR failure to acknowledge. If 0x206 also stale, MTR
comms lost. \\
External WDT timeout & TPS3850 MR pin & MCU hardware reset → all outputs safe
state \\
Queue full & \texttt{xQueueSend} fail & Frame dropped \\
\end{longtable}

\subsection{8.11 Startup}\label{startup-1}

\begin{verbatim}
 1. can_driver_init()     → TWAI, low-level CAN
 2. mode_manager_init()   → GPIO11 (MODE), GPIO32 (START)
 3. safety_monitor_init() → GPIO1 (ESTOP), GPIO2 (brake lever), WDT GPIO23
 4. throttle_init()       → ADC1_CH5 + I2C + MCP4725 (output=0)
 5. gear_init()           → GPIO12-14 (IN), GPIO33-35 (OUT, LOW)
 6. lights_init()         → GPIO3,6,7 (IN, switches) + GPIO18-22,25-26 (OUT)
 7. power_init()          → GPIO27 (OUT, LOW)
 8. brake_init()          → BRAKE_BOOT_WAIT (500ms) → LISTEN_SYNC (await 0x721) → ACTIVE
 9. dcdc_init()           → CAN 0x012 enable=0
10. Create queues         → can_rx(16), setpoint(4)
11. Create 15 tasks
12. power_task → 12V relay ON (if not ESTOP)
13. dcdc_task → CAN 0x012 enable=1 (if not ESTOP)
14. safety_task starts WDT toggle → external watchdog armed
15. ESP_LOGI("Ready")
\end{verbatim}

\begin{center}\rule{0.5\linewidth}{0.5pt}\end{center}

\section{9. CAN bus device maps}\label{can-bus-device-maps}

\subsection{Low-level}\label{low-level}

\begin{verbatim}
 Low-Level CAN (500 kbit/s)
  ├── RT ESP32-S3 (TWAI)        TX: 0x169,0x204,0x205,0x302,0x001,0x7FD
  │                              RX: 0x001,0x011,0x110,0x120,0x201,0x202,0x203,0x206,0x600,0x6FA,0x7FD,0x7FE
  ├── SYS ESP32-S3               TX: 0x011,0x012,0x110,0x600,0x7B9,0x001,0x7FE
  │                              RX: 0x001,0x204,0x205,0x206,0x302,0x6FB,0x721,0x731,0x741,0x7FD
  ├── SYNTREE EPS-C (steering)   TX: 0x201 | RX: 0x169
  ├── SYNTREE SEB (brake)        TX: 0x721 | RX: 0x7B9
  └── DC-DC converter (72→12V)  RX: 0x012
\end{verbatim}

\subsection{High-level}\label{high-level}

\begin{verbatim}
 High-Level CAN (500 kbit/s)
  ├── Jetson Orin             TX: 0x300,0x301,0x302,0x001,0x7FC
  │                              RX: 0x001,0x011,0x120,0x210,0x220,0x400,0x600,0x7FD
  └── RT ESP32-S3 (MCP2515)      TX: 0x011,0x120,0x210,0x220,0x400,0x600,0x001,0x7FD
                                  RX: 0x001,0x300,0x301,0x302,0x7FC
\end{verbatim}

\begin{center}\rule{0.5\linewidth}{0.5pt}\end{center}

\section{10. Hardware summary}\label{hardware-summary}

\begin{longtable}[]{@{}
  >{\raggedright\arraybackslash}p{(\columnwidth - 8\tabcolsep) * \real{0.1250}}
  >{\raggedright\arraybackslash}p{(\columnwidth - 8\tabcolsep) * \real{0.2292}}
  >{\raggedright\arraybackslash}p{(\columnwidth - 8\tabcolsep) * \real{0.1042}}
  >{\raggedright\arraybackslash}p{(\columnwidth - 8\tabcolsep) * \real{0.2292}}
  >{\raggedright\arraybackslash}p{(\columnwidth - 8\tabcolsep) * \real{0.3125}}@{}}
\toprule\noalign{}
\begin{minipage}[b]{\linewidth}\raggedright
Node
\end{minipage} & \begin{minipage}[b]{\linewidth}\raggedright
Controller
\end{minipage} & \begin{minipage}[b]{\linewidth}\raggedright
MCU
\end{minipage} & \begin{minipage}[b]{\linewidth}\raggedright
Framework
\end{minipage} & \begin{minipage}[b]{\linewidth}\raggedright
CAN interfaces
\end{minipage} \\
\midrule\noalign{}
\endhead
\bottomrule\noalign{}
\endlastfoot
Jetson & Orin & --- & ROS 2 & 1× CAN (high) \\
RT & ESP32-S3 @ 240 MHz & Xtensa LX7 & ESP-IDF + FreeRTOS & TWAI (low) + MCP2515
SPI (high) \\
SYS & ESP32-S3 @ 240 MHz & Xtensa LX7 & ESP-IDF + FreeRTOS & TWAI (low only) \\
\end{longtable}

\begin{longtable}[]{@{}ll@{}}
\toprule\noalign{}
Parameter & Value \\
\midrule\noalign{}
\endhead
\bottomrule\noalign{}
\endlastfoot
CAN bitrate (both) & 500 kbit/s \\
CAN transceiver & SN65HVD230 \\
FreeRTOS tick & 1000 Hz \\
RT tasks & 8 \\
SYS tasks & 15 \\
\end{longtable}

\begin{center}\rule{0.5\linewidth}{0.5pt}\end{center}

\section{11. Build}\label{build}

\begin{verbatim}
cd rt-esp32 && pio run && pio run -t upload && pio device monitor
cd sys-esp32 && pio run && pio run -t upload && pio device monitor
\end{verbatim}

\begin{center}\rule{0.5\linewidth}{0.5pt}\end{center}

\section{12. Known design gaps}\label{known-design-gaps}

\begin{longtable}[]{@{}
  >{\raggedright\arraybackslash}p{(\columnwidth - 6\tabcolsep) * \real{0.1071}}
  >{\raggedright\arraybackslash}p{(\columnwidth - 6\tabcolsep) * \real{0.1786}}
  >{\raggedright\arraybackslash}p{(\columnwidth - 6\tabcolsep) * \real{0.2857}}
  >{\raggedright\arraybackslash}p{(\columnwidth - 6\tabcolsep) * \real{0.4286}}@{}}
\toprule\noalign{}
\begin{minipage}[b]{\linewidth}\raggedright
\#
\end{minipage} & \begin{minipage}[b]{\linewidth}\raggedright
Gap
\end{minipage} & \begin{minipage}[b]{\linewidth}\raggedright
Impact
\end{minipage} & \begin{minipage}[b]{\linewidth}\raggedright
Resolution
\end{minipage} \\
\midrule\noalign{}
\endhead
\bottomrule\noalign{}
\endlastfoot
1 & \st{Brake arbitration has no CAN path to SYS} & \st{Jetson
\mbox{\texttt{0x301}} + RT obstacle braking never actuated} & \textbf{RESOLVED:}
\texttt{0x205\ RT\_BRAKE\_CMD} (RT→SYS, DLC=4, i32 kPa, 50 Hz). SYS maps kPa→SEB
Pressure Mode. Mode-switching protocol defined in §8.6. \\
2 & \st{No CAN message for Jetson to request S (Sport) gear} & \st{AUTO can only
select D/N/R} & \textbf{RESOLVED:} \texttt{0x300\ HOST\_DRIVE\_CMD} repacked:
i32 speed + i24 yaw + u8 gear (N=0,D=1,S=2,R=3). RT passes gear through to
\texttt{0x204}. \\
3 & \st{EPS-C timeout-fault behavior unknown} & \st{On ESTOP or comm loss,
steering may lock, center, or freewheel} & \textbf{RESOLVED:} Two-tier ESTOP
steering (§7.6). Active-zero centering for non-obstacle ESTOP with silent-stop
fallback on mechanical jam. Obstacle-triggered ESTOP holds current angle. EPS-C
timeout-fault is now a last-resort fallback (CAN bus dead), not the primary
ESTOP mechanism. \\
4 & \st{SEB pressure control mode not defined} & \st{SYS currently uses stroke
mode only} & \textbf{RESOLVED:} Verified SYNTREE SEB spec:
\texttt{VCU\_SEB\_Pre\_Value\_Req} is u8 at bit 32, scale 0.05 MPa/bit, range
0--5 MPa (raw 0--100). Conversion: \texttt{seb\_raw\ =\ kPa\ ×\ 0.02}. Clamp to
100. \\
5 & Speed control open-loop --- PID exists but encoder not fitted &
\texttt{speed\_pid.cpp} is correct but runs with \texttt{measured=0} (encoder
not physically installed). PID must NOT be in the active control path until
GPIO1/2 PCNT is wired. Currently runs as shadow controller → \texttt{0x220}
telemetry for validation. & 1) Fit rear motor encoder to GPIO1/2. 2) Verify
encoder pulses on PCNT. 3) Enable \texttt{pid\_update(desired,\ measured,\ dt)}
with real data. 4) Route PID effort trim to SYS via new field in \texttt{0x204}
or dedicated ID. 5) Validate closed-loop response against open-loop baseline. \\
6 & \textbf{ESTOP exit race condition --- START button during steering ramp} &
Pressing START before non-obstacle centering ramp completes stops 0x169
mid-ramp, leaving EPS-C to comm-fault at an off-center angle. A 30° steering
lock in MANUAL mode (where EPS-C is standalone with no active centering) makes
the vehicle unrideable. & \textbf{Software fix:} RT defers ESTOP→MANUAL steering
transition until ramp completes. Brake/motor/lights transition immediately. No
new CAN messages needed. See \texttt{docs/emergency-safety-analysis.md} §1. \\
7 & \textbf{SEB BRAKE\_FAULT makes physical brake lever inoperative} & Original
BRAKE\_FAULT state permanently stops 0x7B9 transmission on boot sync timeout.
Since SEB is by-wire (electro-hydraulic, no mechanical linkage), this disables
the brake lever --- contradicting the claim ``brake lever always works.'' &
\textbf{Software fix:} Replaced BRAKE\_FAULT with BRAKE\_DEGRADED (§8.6). On
sync timeout, transmit 0x7B9 with lever-based defaults (0mm/15mm) without
waiting for sync. Recovers when 0x721 arrives. Lever always functional. See
\texttt{docs/emergency-safety-analysis.md} §2. \\
8 & \textbf{Watchdog reset unbraked window --- SEB comm-fault behavior
unverified} & When SYS watchdog resets, SEB enters comm-fault after 20ms. If SEB
releases on timeout, the vehicle coasts without brake for \textasciitilde2.5s
(SYS reboot + brake LBS). If SEB holds, the window is only 20ms. Behavior is
empirically unverified. & \textbf{Test SEB comm-fault behavior} (stroke=27mm,
stop CAN, measure pressure over 5s). If release: add NC brake-hold relay gated
by TPS3850 RST line. If hold: document as verified safety property. \textbf{Also
mitigated by gap \#11 --- RT brake takeover closes most of the window.} See
\texttt{docs/emergency-safety-analysis.md} §3. \\
9 & \textbf{Obstacle ESTOP ``hold angle'' can cause rollover during cornering} &
Obstacle-triggered ESTOP holds current steering angle regardless of vehicle
speed. A cornering vehicle under hard braking experiences lateral load transfer
that the dynamic angle clamp was designed to prevent --- but the clamp is only
applied to commanded angles, not to the ESTOP hold angle. Physics model §8
defines the rollover threshold: a\_y = v²/L·tan(δ) \textgreater{} g·w/(2h). &
\textbf{Software fix:} During obstacle-triggered ESTOP, clamp the hold angle to
the dynamic angle clamp limit for the current speed. If current angle exceeds
the limit, ramp down to the limit at 20°/s. Straight-line cases unchanged. See
\texttt{docs/emergency-safety-analysis.md} §4. \\
10 & \textbf{Jetson heartbeat loss → pure coast is insufficient for perception
failure} & Jetson runs obstacle detection. When it dies, the vehicle coasts with
no active brake, no perception, and no steering. At 25 km/h, \textgreater50m
coast-to-stop. The rider must recognize the failure and manually brake. The
heartbeat timeout (1500ms) is long enough that false positives from Linux CAN
jitter are unlikely. & \textbf{Software fix:} On Jetson heartbeat loss, RT
commands 0x205 = 2000 kPa (\textasciitilde2 MPa moderate brake) + transitions
SYS to MANUAL. Brake light ON. Rider can override with lever. DC-DC stays on
(lights work). This is ``assisted stop'' --- between coast and full ESTOP. See
\texttt{docs/emergency-safety-analysis.md} §5. \\
11 & \textbf{No secondary ESTOP exit path --- START button is a single point of
failure} & GPIO32 is the only ESTOP exit. Stuck HIGH (broken wire) = can never
exit ESTOP. The only backup is power-cycle, which restarts all nodes and
requires brake/steering LBS at roadside. & \textbf{Software fix:} MODE button
(GPIO11) long-press (3s) exits ESTOP → MANUAL as secondary path. Two independent
GPIOs on separate physical buttons. Add START button health monitoring in
diag\_task. See \texttt{docs/emergency-safety-analysis.md} §6. \\
12 & \textbf{SYS crash --- 1000ms brake gap before RT responds, and RT doesn't
take over 0x7B9} & Architecture §6.2 Option D explicitly claims ``SYS failure →
brake lost? No (RT takes over).'' But RT only broadcasts CAN 0x001 on SYS
heartbeat loss --- which SEB ignores (it only responds to 0x7B9). The 1000ms
heartbeat timeout creates a \textasciitilde980ms brake gap. The heartbeat
justification (1000ms acceptable because fast-path checks cover actuators) omits
the brake --- there is no fast-path check for SEB command loss. &
\textbf{Software fix:} (1) Increase SYS heartbeat rate from 2 Hz to 10 Hz; RT
timeout shortened to 200ms (2 missed frames at 100ms period). (2) On SYS
heartbeat loss, RT immediately takes over 0x7B9 transmission with stroke=max,
regardless of current mode (mode-gate opens on emergency). Brake gap: 220ms
worst case. \textbf{Also mitigates gap \#8} by reducing watchdog brake window
from \textasciitilde2.5s to \textasciitilde220ms. See
\texttt{docs/emergency-safety-analysis.md} §7. \\
13 & \textbf{No independent brake monitor --- ESTOP could silently have no
brakes} & All 8 safety layers protect motor+steering; none verify SEB actually
applied braking force. If SEB's CAN receiver is faulted, ESTOP
\texttt{0x7B9\ stroke=max} is never received and the system never knows.
SEB\_STATUS (\texttt{0x721}, 100 Hz) already provides SEB\_Stroke\_Value,
SEB\_Pressure\_Value, and SEB\_Error\_Status (Level 3 = severe/shutdown). This
is sufficient for a brake following-error monitor with zero new sensors. &
\textbf{Software fix:} Add brake following-error check in SYS \texttt{dispatch}:
compare 0x7B9 cmd stroke vs 0x721 actual stroke (threshold 3mm, debounce 100ms).
Monitor SEB\_Error\_Status for Level 3 faults. Add 0x721 staleness check (100ms
timeout). Faults logged via 0x600 diagnostic --- cannot escalate beyond ESTOP
(already max brake) but essential for incident analysis. See
\texttt{docs/emergency-safety-analysis.md} §8. \\
14 & \textbf{CAN 0x001 spoofable --- no authentication, DoS-vulnerable} & 0x001
is DLC=0 with no sender ID, checksum, or rolling counter. A corrupted node can
flood 0x001, triggering persistent ESTOP and saturating the CAN bus. CAN error
containment (bus-off at TEC\textgreater255) eventually isolates the node, but
the flood window is disruptive. & \textbf{Software fix:} (1) Rate-limit 0x001
processing: max 2 frames per 500ms window per node (the first one already
triggered ESTOP). (2) Change 0x001 DLC from 0 to 1 with sender-ID byte (0=SYS,
1=RT, 2=Jetson, 3=MTR) for diagnostics and per-sender rate limiting. See
\texttt{docs/emergency-safety-analysis.md} §9. \\
15 & \textbf{No ESTOP acknowledgment from MTR STM32} & SYS/RT have no
confirmation that MTR received and acted on ESTOP. MTR has no dedicated
heartbeat. If MTR freezes at zero output, the SYS monitor sees 0x204=0, 0x206=0
and assumes MTR responded correctly. No 0x206 staleness check exists. &
\textbf{Software fix:} MTR sets \texttt{ESTOP\_ACTIVE} bit in 0x206 fault\_flags
when it has locally cut throttle+gear. SYS checks for ACK within 100ms of ESTOP.
Add 0x206 staleness check (\textgreater200ms → MTR comms lost). See
\texttt{docs/emergency-safety-analysis.md} §10. \\
16 & \textbf{Startup grace period masks heartbeat but not 0x204 staleness} & On
cold boot, 0x204 is absent for \textgreater200ms before RT comes online. SYS
0x204 staleness check triggers (harmlessly --- forces speed=0 which is already
default). But the startup grace concept is applied inconsistently: heartbeat
checks are masked, staleness is not. & \textbf{Trivial fix:} Gate 0x204
staleness check with same 3-second startup grace period as heartbeat. See
\texttt{docs/emergency-safety-analysis.md} §11. \\
17 & \textbf{ESTOP HMI ambiguous --- ``both bulbs OFF'' identical to powered-off
vehicle} & During ESTOP, DC-DC is OFF → 12V rail dead → brake light, mode bulbs,
and indicators all dark. A rider returning to the vehicle cannot distinguish
ESTOP from power-off. The OR logic claim ``brake light ON during ESTOP'' is
physically impossible with DC-DC off. & \textbf{Software + wiring fix:} Keep
DC-DC ON during ESTOP (needed for MCU power anyway). Cut only the 12V accessory
relay (GPIO27). Rewire brake light to always-on DC-DC rail (not accessory relay
output). Result: ESTOP = brake light ON + mode bulbs OFF. Power-off = everything
OFF. See \texttt{docs/emergency-safety-analysis.md} §12. \\
18 & \textbf{EPS-C mechanical jam silent-stop recovery path unclear} & When
mechanical jam triggers silent-stop during ESTOP centering, the architecture
says ``fall back to silent-stop'' without specifying steer SM state transition.
It's unclear whether the jam is recoverable via START short-press (STEER\_FAULT
path) or requires power-cycle. & \textbf{Documentation fix:} Explicitly
transition to STEER\_FAULT on mechanical jam during ESTOP centering. Existing
STEER\_FAULT recovery paths apply: START short-press → LISTEN\_SYNC retry; START
long-press 3s + throttle=0 → force-activate at 0° (MANUAL only). See
\texttt{docs/emergency-safety-analysis.md} §13. \\
19 & \textbf{Fixed 5° steering following error threshold wrong at high speed} &
At 25 km/h, dynamic clamp limits steering to \textasciitilde5°. A fixed 5°
threshold represents a 100\% error --- the EPS-C must have ZERO response to
trigger. A 4° error at speed (80\% authority loss) would NOT trigger ESTOP. At 2
km/h (40° limit), 5° is only 12.5\% --- potentially too sensitive for parking
maneuvers. & \textbf{Software fix:} Speed-scaled threshold:
\texttt{max(2°,\ 0.25\ ×\ dynamic\_limit)}. Result: 2° at 25 km/h (tight), 4.5°
at 10 km/h, 10° at 2 km/h (tolerant). One-line change in RT following-error
check. See \texttt{docs/emergency-safety-analysis.md} §14. \\
\end{longtable}

\begin{center}\rule{0.5\linewidth}{0.5pt}\end{center}

\section{13. Reference documents}\label{reference-documents}

\begin{longtable}[]{@{}
  >{\raggedright\arraybackslash}p{(\columnwidth - 2\tabcolsep) * \real{0.4000}}
  >{\raggedright\arraybackslash}p{(\columnwidth - 2\tabcolsep) * \real{0.6000}}@{}}
\toprule\noalign{}
\begin{minipage}[b]{\linewidth}\raggedright
File
\end{minipage} & \begin{minipage}[b]{\linewidth}\raggedright
Content
\end{minipage} \\
\midrule\noalign{}
\endhead
\bottomrule\noalign{}
\endlastfoot
\href{can-dictionary.md}{\texttt{can-dictionary.md}} & Bit-level CAN signal
layouts for all IDs on both buses \\
\href{docs/steering-unit.md}{\texttt{docs/steering-unit.md}} & SYNTREE EPS-C
protocol reference \\
\href{docs/brake-unit.md}{\texttt{docs/brake-unit.md}} & SYNTREE SEB protocol
reference \\
\href{rt-esp32/README.md}{\texttt{rt-esp32/README.md}} & RT build \& test \\
\href{sys-esp32/README.md}{\texttt{sys-esp32/README.md}} & SYS build \& test \\
\href{notes/can-protocol.md}{\texttt{notes/can-protocol.md}} & CAN protocol
theory --- arbitration, frame types, standards \\
\href{notes/can-hardware-basics.md}{\texttt{notes/can-hardware-basics.md}} & CAN
physical layer --- termination, topology, transceivers \\
\href{legacy/notes/can-addressing-for-etrike.md}{\texttt{legacy/notes/can-addressing-for-etrike.md}}
& DEPRECATED --- CAN addressing scheme (single-ESP32 design, historical
reference) \\
\href{notes/can-troubleshooting.md}{\texttt{notes/can-troubleshooting.md}} & CAN
debugging --- common mistakes, error states, tools \\
\href{notes/rtos-task-design.md}{\texttt{notes/rtos-task-design.md}} & RTOS task
design --- priorities, preemptive scheduling, queues, tick rate \\
\href{notes/pid-control.md}{\texttt{notes/pid-control.md}} & PID control
fundamentals --- P/I/D terms, integral windup, tuning procedure \\
\href{notes/state-machine-design.md}{\texttt{notes/state-machine-design.md}} &
State machine design --- FSMs, hierarchical states, C++ implementation
patterns \\
\href{notes/heartbeat-monitoring.md}{\texttt{notes/heartbeat-monitoring.md}} &
Heartbeat \& liveness --- alive counters, FTTI, timeout selection, startup
grace \\
\href{notes/distributed-safety-patterns.md}{\texttt{notes/distributed-safety-patterns.md}}
& Distributed safety --- defense in depth, ESTOP, fail-safe, NC wiring,
debounce \\
\href{notes/endianness-binary-protocols.md}{\texttt{notes/endianness-binary-protocols.md}}
& Endianness \& binary protocols --- big/little endian, bit numbering, struct
packing \\
\href{notes/can-gateway-design.md}{\texttt{notes/can-gateway-design.md}} & CAN
gateway design --- multi-bus architecture, forwarding categories, isolation \\
\href{notes/analog-interfacing.md}{\texttt{notes/analog-interfacing.md}} &
Analog interfacing --- ADC, DAC, optoisolators, relays, protection circuits \\
\href{docs/physics-model.md}{\texttt{docs/physics-model.md}} & Tricycle
kinematics --- forward/inverse, rollover, slip angles \\
\href{docs/listen-before-speaking.md}{\texttt{docs/listen-before-speaking.md}} &
CAN actuator safe bootstrapping pattern \\
\href{docs/can-gateway-bridging.md}{\texttt{docs/can-gateway-bridging.md}} & CAN
gateway forwarding rules and implementation \\
\href{docs/emergency-system.md}{\texttt{docs/emergency-system.md}} &
\textbf{Primary:} ESTOP triggers, 8-layer defense, emergency response matrix,
rider's guide, EGAS 3-level architecture, testing procedures \\
\href{docs/emergency-safety-analysis.md}{\texttt{docs/emergency-safety-analysis.md}}
& \textbf{Safety analysis:} ESTOP exit race condition, SEB brake lever
contradiction, watchdog unbraked window --- causal traces, risk quantification,
recommended fixes \\
\href{docs/defense-in-depth-safety.md}{\texttt{docs/defense-in-depth-safety.md}}
& Layered safety --- ESTOP, following error, dynamic clamp, OR logic \\
\href{docs/syntree-security-protocol.md}{\texttt{docs/syntree-security-protocol.md}}
& Rolling counter + XOR checksum for SYNTREE actuators \\
\href{docs/high-voltage-isolation.md}{\texttt{docs/high-voltage-isolation.md}} &
72V galvanic isolation --- TLP281 optos, relays, fuses, TVS \\
\href{docs/distributed-architecture.md}{\texttt{docs/distributed-architecture.md}}
& Three-node rationale --- Jetson/RT/SYS split, dual-CAN on RT \\
\href{docs/actuator-interfacing.md}{\texttt{docs/actuator-interfacing.md}} &
MCP4725 DAC throttle, gear pass-through, relay logic \\
\href{docs/external-watchdog.md}{\texttt{docs/external-watchdog.md}} & External
watchdog IC --- timeout, safe state, testing \\
\href{docs/pid-speed-control.md}{\texttt{docs/pid-speed-control.md}} & PID speed
control theory, tuning, anti-windup \\
\end{longtable}


\chapter{CAN Signal Dictionary --- E-Trike}\label{can-signal-dictionary-e-trike}

Two physical CAN buses at 500 kbit/s. All fields big-endian (MSB first) unless
noted (SYNTREE protocol uses Motorola LSB).

\subsection{Type Notation}\label{type-notation}

\begin{longtable}[]{@{}lll@{}}
\toprule\noalign{}
Notation & C Type & Meaning \\
\midrule\noalign{}
\endhead
\bottomrule\noalign{}
\endlastfoot
\texttt{i8\ /\ i16\ /\ i32} & \texttt{int8\_t\ /\ int16\_t\ /\ int32\_t} &
Signed integer \\
\texttt{i24} & 24-bit signed (packed) & Non-standard width, CAN frame only \\
\texttt{u8\ /\ u16\ /\ u32} & \texttt{uint8\_t\ /\ uint16\_t\ /\ uint32\_t} &
Unsigned integer \\
\texttt{u8\ bool} & \texttt{uint8\_t} & Boolean in a byte (0 or 1) \\
\texttt{u8\ enum} & \texttt{uint8\_t} & Enumeration in a byte \\
\texttt{u8\ bitmask} & \texttt{uint8\_t} & Bitfield, each bit is a flag \\
\texttt{DLC=0} & --- & Zero-length CAN frame (event signal, no payload) \\
\end{longtable}

\subsection{Physical Units}\label{physical-units}

\begin{longtable}[]{@{}lll@{}}
\toprule\noalign{}
Unit & Meaning & Example \\
\midrule\noalign{}
\endhead
\bottomrule\noalign{}
\endlastfoot
\texttt{mm/s} & Millimeters per second & Speed: 3000 = 3.0 m/s \\
\texttt{mrad/s} & Milliradians per second & Yaw rate \\
\texttt{kPa} & Kilopascals & Brake pressure \\
\texttt{mm} & Millimeters & Distance \\
\texttt{0.1°/bit} & Tenths of a degree per LSB & Steer angle: 455 = 45.5° \\
\texttt{0.05\ mm/bit} & 0.05 mm per LSB & Brake stroke \\
\texttt{0.05\ MPa/bit} & 0.05 MPa per LSB & Brake pressure \\
\texttt{°/s} & Degrees per second & Steer angle speed \\
\texttt{Nm} & Newton-meters & Torque \\
\end{longtable}

\begin{center}\rule{0.5\linewidth}{0.5pt}\end{center}

\section{1. Low-Level CAN Bus}\label{low-level-can-bus}

Nodes: RT ESP32-S3, SYS ESP32-S3, SYNTREE EPS-C (steering), SYNTREE SEB (brake),
DC-DC converter.

\begin{center}\rule{0.5\linewidth}{0.5pt}\end{center}

\subsection{0x001 --- SAFETY\_ESTOP}\label{x001-safety_estop}

\begin{longtable}[]{@{}ll@{}}
\toprule\noalign{}
Property & Value \\
\midrule\noalign{}
\endhead
\bottomrule\noalign{}
\endlastfoot
\textbf{Sender} & Any (RT, SYS) \\
\textbf{Receiver(s)} & All nodes \\
\textbf{DLC} & 0 \\
\textbf{Period} & On event \\
\end{longtable}

Presence of this frame = emergency stop. Motor stop, brake engage, steering
disable, DCDC off.

\begin{center}\rule{0.5\linewidth}{0.5pt}\end{center}

\subsection{0x011 --- SYS\_SAFETY\_STS}\label{x011-sys_safety_sts}

\begin{longtable}[]{@{}ll@{}}
\toprule\noalign{}
Property & Value \\
\midrule\noalign{}
\endhead
\bottomrule\noalign{}
\endlastfoot
\textbf{Sender} & SYS \\
\textbf{Receiver(s)} & RT (→ Jetson) \\
\textbf{DLC} & 2 \\
\textbf{Period} & 5 Hz \\
\end{longtable}

\begin{longtable}[]{@{}lllllllll@{}}
\toprule\noalign{}
Signal & Start bit & Len & Type & Scale & Offset & Min & Max & Unit \\
\midrule\noalign{}
\endhead
\bottomrule\noalign{}
\endlastfoot
\texttt{SYS\_EstopActive} & 0 & 8 & u8 & 1 & 0 & 0 & 1 & --- \\
\texttt{SYS\_HeartbeatOk} & 8 & 8 & u8 & 1 & 0 & 0 & 1 & --- \\
\end{longtable}

\begin{center}\rule{0.5\linewidth}{0.5pt}\end{center}

\subsection{0x012 --- SYS\_DCDC\_CMD}\label{x012-sys_dcdc_cmd}

\begin{longtable}[]{@{}ll@{}}
\toprule\noalign{}
Property & Value \\
\midrule\noalign{}
\endhead
\bottomrule\noalign{}
\endlastfoot
\textbf{Sender} & SYS \\
\textbf{Receiver(s)} & DC-DC converter (72V→12V) \\
\textbf{DLC} & 1 \\
\textbf{Period} & On change \\
\end{longtable}

\begin{longtable}[]{@{}lllllllll@{}}
\toprule\noalign{}
Signal & Start bit & Len & Type & Scale & Offset & Min & Max & Unit \\
\midrule\noalign{}
\endhead
\bottomrule\noalign{}
\endlastfoot
\texttt{SYS\_DcdcEnable} & 0 & 8 & u8 & 1 & 0 & 0 & 1 & --- \\
\end{longtable}

ESTOP → \textbf{1 (on)} --- maintains 12V for MCUs, CAN transceivers, and brake
light. The 12V accessory relay (GPIO27) provides the secondary cut for
non-safety loads (see architecture §8.6). All other modes → 1 (on).

\begin{center}\rule{0.5\linewidth}{0.5pt}\end{center}

\subsection{0x110 --- SYS\_MODE\_CMD}\label{x110-sys_mode_cmd}

\begin{longtable}[]{@{}ll@{}}
\toprule\noalign{}
Property & Value \\
\midrule\noalign{}
\endhead
\bottomrule\noalign{}
\endlastfoot
\textbf{Sender} & SYS \\
\textbf{Receiver(s)} & RT \\
\textbf{DLC} & 1 \\
\textbf{Period} & On change \\
\end{longtable}

\begin{longtable}[]{@{}lllllllll@{}}
\toprule\noalign{}
Signal & Start bit & Len & Type & Scale & Offset & Min & Max & Unit \\
\midrule\noalign{}
\endhead
\bottomrule\noalign{}
\endlastfoot
\texttt{SYS\_Mode} & 0 & 8 & u8 & 1 & 0 & 0 & 2 & enum (0=M, 1=A, 2=ESTOP) \\
\end{longtable}

\begin{center}\rule{0.5\linewidth}{0.5pt}\end{center}

\subsection{0x120 --- SYS\_THROTTLE\_STS}\label{x120-sys_throttle_sts}

\begin{longtable}[]{@{}ll@{}}
\toprule\noalign{}
Property & Value \\
\midrule\noalign{}
\endhead
\bottomrule\noalign{}
\endlastfoot
\textbf{Sender} & MTR \\
\textbf{Receiver(s)} & RT (→ Jetson) \\
\textbf{DLC} & 2 \\
\textbf{Period} & 100 Hz \\
\end{longtable}

\begin{longtable}[]{@{}lllllllll@{}}
\toprule\noalign{}
Signal & Start bit & Len & Type & Scale & Offset & Min & Max & Unit \\
\midrule\noalign{}
\endhead
\bottomrule\noalign{}
\endlastfoot
\texttt{SYS\_ThrottleSpeed} & 0 & 16 & i16 & 1 & 0 & -500 & 3000 & mm/s \\
\end{longtable}

\begin{center}\rule{0.5\linewidth}{0.5pt}\end{center}

\subsection{0x204 --- RT\_DRIVE\_CMD}\label{x204-rt_drive_cmd}

\begin{longtable}[]{@{}ll@{}}
\toprule\noalign{}
Property & Value \\
\midrule\noalign{}
\endhead
\bottomrule\noalign{}
\endlastfoot
\textbf{Sender} & RT \\
\textbf{Receiver(s)} & SYS, \textbf{MTR (STM32)} \\
\textbf{DLC} & 5 \\
\textbf{Period} & 100 Hz \\
\end{longtable}

\begin{longtable}[]{@{}lllllllll@{}}
\toprule\noalign{}
Signal & Start bit & Len & Type & Scale & Offset & Min & Max & Unit \\
\midrule\noalign{}
\endhead
\bottomrule\noalign{}
\endlastfoot
\texttt{RT\_MotorSpeed} & 0 & 32 & i32 & 1 & 0 & -500 & 3000 & mm/s \\
\texttt{RT\_Gear} & 32 & 8 & u8 & 1 & 0 & 0 & 3 & enum (0=N,1=D,2=S,3=R) \\
\end{longtable}

Byte layout (big-endian): Byte 0-3 = speed {[}31:0{]}, Byte 4 = gear.

MTR receives 0x204 directly for motor actuation (speed → MCP4725 DAC, gear →
relay module). SYS also receives 0x204 for EGAS Level 2 monitoring (compares
setpoint vs 0x206 feedback).

\begin{quote}
Placed at \texttt{0x204} to avoid collision with EPS-C error info frame at
\texttt{0x202}. SYNTREE units are preprogrammed and cannot be reconfigured.
\end{quote}

\begin{center}\rule{0.5\linewidth}{0.5pt}\end{center}

\subsection{0x205 --- RT\_BRAKE\_CMD}\label{x205-rt_brake_cmd}

\begin{longtable}[]{@{}ll@{}}
\toprule\noalign{}
Property & Value \\
\midrule\noalign{}
\endhead
\bottomrule\noalign{}
\endlastfoot
\textbf{Sender} & RT \\
\textbf{Receiver(s)} & SYS \\
\textbf{DLC} & 4 \\
\textbf{Period} & 50 Hz (20 ms) \\
\end{longtable}

\begin{longtable}[]{@{}lllllllll@{}}
\toprule\noalign{}
Signal & Start bit & Len & Type & Scale & Offset & Min & Max & Unit \\
\midrule\noalign{}
\endhead
\bottomrule\noalign{}
\endlastfoot
\texttt{RT\_BrakePressure} & 0 & 32 & i32 & 1 & 0 & 0 & 20000 & kPa \\
\end{longtable}

Byte layout (big-endian): Bytes 0-3 = brake pressure {[}kPa{]}.

RT max-select: \texttt{brake\_kpa\ =\ max(rt\_obstacle,\ jetson\_0x301)}. SYS
converts: \texttt{seb\_raw\ =\ (uint8\_t)(kpa\ *\ 0.02f)} (verified SYNTREE
spec: \texttt{VCU\_SEB\_Pre\_Value\_Req} is u8, scale 0.05 MPa/bit, range 0--5
MPa). When \texttt{0x205\ \textgreater{}\ 0}, SYS switches SEB to Pressure Mode
(mode=2). When \texttt{0x205\ ==\ 0}, falls back to Stroke Mode for lever/ESTOP
triggers.

\subsection{0x206 --- MTR\_MOTOR\_FBK}\label{x206-mtr_motor_fbk}

\begin{longtable}[]{@{}ll@{}}
\toprule\noalign{}
Property & Value \\
\midrule\noalign{}
\endhead
\bottomrule\noalign{}
\endlastfoot
\textbf{Sender} & MTR STM32 \\
\textbf{Receiver(s)} & SYS, RT \\
\textbf{DLC} & 4 \\
\textbf{Period} & 50 Hz (20 ms) \\
\end{longtable}

\begin{longtable}[]{@{}
  >{\raggedright\arraybackslash}p{(\columnwidth - 16\tabcolsep) * \real{0.1311}}
  >{\raggedright\arraybackslash}p{(\columnwidth - 16\tabcolsep) * \real{0.1803}}
  >{\raggedright\arraybackslash}p{(\columnwidth - 16\tabcolsep) * \real{0.0820}}
  >{\raggedright\arraybackslash}p{(\columnwidth - 16\tabcolsep) * \real{0.0984}}
  >{\raggedright\arraybackslash}p{(\columnwidth - 16\tabcolsep) * \real{0.1148}}
  >{\raggedright\arraybackslash}p{(\columnwidth - 16\tabcolsep) * \real{0.1311}}
  >{\raggedright\arraybackslash}p{(\columnwidth - 16\tabcolsep) * \real{0.0820}}
  >{\raggedright\arraybackslash}p{(\columnwidth - 16\tabcolsep) * \real{0.0820}}
  >{\raggedright\arraybackslash}p{(\columnwidth - 16\tabcolsep) * \real{0.0984}}@{}}
\toprule\noalign{}
\begin{minipage}[b]{\linewidth}\raggedright
Signal
\end{minipage} & \begin{minipage}[b]{\linewidth}\raggedright
Start bit
\end{minipage} & \begin{minipage}[b]{\linewidth}\raggedright
Len
\end{minipage} & \begin{minipage}[b]{\linewidth}\raggedright
Type
\end{minipage} & \begin{minipage}[b]{\linewidth}\raggedright
Scale
\end{minipage} & \begin{minipage}[b]{\linewidth}\raggedright
Offset
\end{minipage} & \begin{minipage}[b]{\linewidth}\raggedright
Min
\end{minipage} & \begin{minipage}[b]{\linewidth}\raggedright
Max
\end{minipage} & \begin{minipage}[b]{\linewidth}\raggedright
Unit
\end{minipage} \\
\midrule\noalign{}
\endhead
\bottomrule\noalign{}
\endlastfoot
\texttt{MTR\_ActualSpeed} & 0 & 16 & i16 & 1 & 0 & -500 & 3000 & mm/s \\
\texttt{MTR\_GearState} & 16 & 8 & u8 enum & 1 & 0 & 0 & 3 &
\{0=N,1=D,2=S,3=R\} \\
\texttt{MTR\_FaultFlags} & 24 & 8 & u8 bitmask & --- & --- & --- & --- &
bit0=ESTOP active, bit1=CMD timeout (0x204 stale), bit2=ADC fault, bit3=gear
conflict \\
\end{longtable}

Byte layout (big-endian): Bytes 0-1=speed, Byte 2=gear, Byte 3=faults.

\subsection{0x201 --- SES\_STATUS (SYNTREE EPS-C
Feedback)}\label{x201-ses_status-syntree-eps-c-feedback}

\begin{longtable}[]{@{}ll@{}}
\toprule\noalign{}
Property & Value \\
\midrule\noalign{}
\endhead
\bottomrule\noalign{}
\endlastfoot
\textbf{Sender} & SYNTREE EPS-C (steering module) \\
\textbf{Receiver(s)} & RT \\
\textbf{DLC} & 8 \\
\textbf{Period} & 10 ms (100 Hz) \\
\textbf{Endianness} & Motorola LSB (little-endian) \\
\end{longtable}

\begin{longtable}[]{@{}
  >{\raggedright\arraybackslash}p{(\columnwidth - 18\tabcolsep) * \real{0.1081}}
  >{\raggedright\arraybackslash}p{(\columnwidth - 18\tabcolsep) * \real{0.1486}}
  >{\raggedright\arraybackslash}p{(\columnwidth - 18\tabcolsep) * \real{0.0676}}
  >{\raggedright\arraybackslash}p{(\columnwidth - 18\tabcolsep) * \real{0.0811}}
  >{\raggedright\arraybackslash}p{(\columnwidth - 18\tabcolsep) * \real{0.0946}}
  >{\raggedright\arraybackslash}p{(\columnwidth - 18\tabcolsep) * \real{0.1081}}
  >{\raggedright\arraybackslash}p{(\columnwidth - 18\tabcolsep) * \real{0.0676}}
  >{\raggedright\arraybackslash}p{(\columnwidth - 18\tabcolsep) * \real{0.0676}}
  >{\raggedright\arraybackslash}p{(\columnwidth - 18\tabcolsep) * \real{0.0811}}
  >{\raggedright\arraybackslash}p{(\columnwidth - 18\tabcolsep) * \real{0.1757}}@{}}
\toprule\noalign{}
\begin{minipage}[b]{\linewidth}\raggedright
Signal
\end{minipage} & \begin{minipage}[b]{\linewidth}\raggedright
Start bit
\end{minipage} & \begin{minipage}[b]{\linewidth}\raggedright
Len
\end{minipage} & \begin{minipage}[b]{\linewidth}\raggedright
Type
\end{minipage} & \begin{minipage}[b]{\linewidth}\raggedright
Scale
\end{minipage} & \begin{minipage}[b]{\linewidth}\raggedright
Offset
\end{minipage} & \begin{minipage}[b]{\linewidth}\raggedright
Min
\end{minipage} & \begin{minipage}[b]{\linewidth}\raggedright
Max
\end{minipage} & \begin{minipage}[b]{\linewidth}\raggedright
Unit
\end{minipage} & \begin{minipage}[b]{\linewidth}\raggedright
Description
\end{minipage} \\
\midrule\noalign{}
\endhead
\bottomrule\noalign{}
\endlastfoot
\texttt{SES\_INF\_Angle\_Status} & 0 & 1 & bool & 1 & 0 & 0 & 1 & --- & Center
Finding Status. 0=Center Finding, 1=Found. (CSV Row 11) \\
\texttt{SES\_Control\_Mode\_Status} & 1 & 2 & u8 & 1 & 0 & 0 & 3 & enum &
Control Mode Feedback. 0=Manual, 1=Automatic. (CSV Row 12) \\
(unaccounted) & 3 & 3 & --- & --- & --- & --- & --- & --- & Byte 0 bits 3--5 ---
not enumerated in CSV \\
\texttt{SES\_Error\_Status} & 6 & 2 & u8 & 1 & 0 & 0 & 3 & enum & Error Status.
0=Normal, 1=L1 Warning, 2=L2, 3=L3. (CSV Row 13) \\
(unaccounted) & --- & 8 & --- & --- & --- & --- & --- & --- & Byte 1 --- not
enumerated in CSV \\
\texttt{SES\_StrAngle} & 16 & 16 & u16 & 0.1 & -3000 & -700 & 700 & ° & Steering
Angle. Unsigned per CSV. Raw 0→-3000°, raw 30000→0°, raw 23000→-700°, raw
37000→700°. (CSV Row 14) \\
\texttt{SES\_Tgt\_StrAngleSpd} & 32 & 16 & i16 & 0.5 & 0 & 0 & 1480 & °/s &
Target Angle Speed. 16-bit signed per CSV. Overlaps Torq at byte 5. (CSV Row
15) \\
\texttt{EPS\_SteeringWheel\_Torq} & 40 & 8 & u8 & 0.1 & -12.1 & -12 & 12 & Nm &
Steering Wheel Torque Feedback. Overlaps StrAngleSpd{[}15:8{]} at byte 5. Init
0x79 (121 raw = 0 Nm). (CSV Row 16) \\
\texttt{SES\_RollCnt\_Enable\_Status} & 48 & 1 & bool & 1 & 0 & 0 & 1 & --- &
Life Signal Enable Feedback. 0=Invalid, 1=Valid. (CSV Row 17) \\
\texttt{SES\_CheckSum\_Enable\_Status} & 49 & 1 & bool & 1 & 0 & 0 & 1 & --- &
Checksum Enable Feedback. 0=Invalid, 1=Valid. (CSV Row 18) \\
(unaccounted) & 50 & 2 & --- & --- & --- & --- & --- & --- & Byte 6 bits 2--3
--- not enumerated in CSV \\
\texttt{SES\_RollCnt\_Status} & 52 & 4 & u8 & 1 & 0 & 0 & 15 & --- & Life Signal
Feedback. Rolling counter 0--15. (CSV Row 19) \\
\texttt{SES\_CheckSum\_Status} & 56 & 8 & u8 & 1 & 0 & 0 & 255 & --- & Checksum
Feedback = XOR(bytes 0--6) \^{} 0xFF. (CSV Row 20) \\
\end{longtable}

\textbf{Byte layout} (little-endian, CSV as source of truth):

\begin{longtable}[]{@{}
  >{\raggedright\arraybackslash}p{(\columnwidth - 16\tabcolsep) * \real{0.2000}}
  >{\raggedright\arraybackslash}p{(\columnwidth - 16\tabcolsep) * \real{0.1000}}
  >{\raggedright\arraybackslash}p{(\columnwidth - 16\tabcolsep) * \real{0.1000}}
  >{\raggedright\arraybackslash}p{(\columnwidth - 16\tabcolsep) * \real{0.1000}}
  >{\raggedright\arraybackslash}p{(\columnwidth - 16\tabcolsep) * \real{0.1000}}
  >{\raggedright\arraybackslash}p{(\columnwidth - 16\tabcolsep) * \real{0.1000}}
  >{\raggedright\arraybackslash}p{(\columnwidth - 16\tabcolsep) * \real{0.1000}}
  >{\raggedright\arraybackslash}p{(\columnwidth - 16\tabcolsep) * \real{0.1000}}
  >{\raggedright\arraybackslash}p{(\columnwidth - 16\tabcolsep) * \real{0.1000}}@{}}
\toprule\noalign{}
\begin{minipage}[b]{\linewidth}\raggedright
Byte
\end{minipage} & \begin{minipage}[b]{\linewidth}\raggedright
0
\end{minipage} & \begin{minipage}[b]{\linewidth}\raggedright
1
\end{minipage} & \begin{minipage}[b]{\linewidth}\raggedright
2
\end{minipage} & \begin{minipage}[b]{\linewidth}\raggedright
3
\end{minipage} & \begin{minipage}[b]{\linewidth}\raggedright
4
\end{minipage} & \begin{minipage}[b]{\linewidth}\raggedright
5
\end{minipage} & \begin{minipage}[b]{\linewidth}\raggedright
6
\end{minipage} & \begin{minipage}[b]{\linewidth}\raggedright
7
\end{minipage} \\
\midrule\noalign{}
\endhead
\bottomrule\noalign{}
\endlastfoot
Content & AngleSts{[}0{]}+ModeSts{[}1:2{]}+(gap)+Error{[}6:7{]} & (unacc.) &
StrAngle {[}7:0{]} & StrAngle {[}15:8{]} & Speed {[}7:0{]} & Speed{[}15:8{]} /
Torq {[}7:0{]} (overlap) & RollCntEn{[}0{]}+CksEn{[}1{]}+(gap)+RollCnt{[}4:7{]}
& CksSum\_Stat \\
\end{longtable}

\begin{quote}
\textbf{StrAngle conversion (CSV Unsigned, offset=-3000):}
\texttt{physical\_deg\ =\ raw\ ×\ 0.1\ −\ 3000}. The EPS-C encodes steering
angle as an unsigned 16-bit value with -3000 offset. Raw 30000 → 0° (straight).
Raw 23000 → -700° (full left). Raw 37000 → 700° (full right). In practice, RT
uses \texttt{internal\_angle\_mdeg\ /\ 100} to produce the raw value; adjust per
actual calibration.

\textbf{Torq encoding (CSV scale=0.1, offset=-12.1):}
\texttt{physical\_Nm\ =\ raw\ ×\ 0.1\ −\ 12.1}. Raw 121 (0x79) → 0.0 Nm (no
torque). Raw 0 → -12.1 Nm. Raw 241 → 12.0 Nm. The EPS-C biases torque readings
so zero torque is at raw 121.

\textbf{Byte 5 overlap:} CSV declares \texttt{SES\_Tgt\_StrAngleSpd} as 16-bit
(bytes 4--5, Signed) AND \texttt{EPS\_SteeringWheel\_Torq} as 8-bit (byte 5).
Both are listed --- the EPS-C may report these in alternate frames or the CSV
represents signals available across firmware versions.
\end{quote}

\begin{center}\rule{0.5\linewidth}{0.5pt}\end{center}

\subsection{0x169 --- VCU\_SES\_REQ (SYNTREE EPS-C
Command)}\label{x169-vcu_ses_req-syntree-eps-c-command}

\begin{longtable}[]{@{}
  >{\raggedright\arraybackslash}p{(\columnwidth - 2\tabcolsep) * \real{0.5882}}
  >{\raggedright\arraybackslash}p{(\columnwidth - 2\tabcolsep) * \real{0.4118}}@{}}
\toprule\noalign{}
\begin{minipage}[b]{\linewidth}\raggedright
Property
\end{minipage} & \begin{minipage}[b]{\linewidth}\raggedright
Value
\end{minipage} \\
\midrule\noalign{}
\endhead
\bottomrule\noalign{}
\endlastfoot
\textbf{Sender} & RT ESP32-S3 \\
\textbf{Receiver(s)} & SYNTREE EPS-C (steering module) \\
\textbf{DLC} & 8 \\
\textbf{Period} & 20 ms (50 Hz) --- \textbf{continuous, every frame} \\
\textbf{Endianness} & Motorola LSB (little-endian) \\
\textbf{Note} & Factory default \texttt{0x169}. SYNTREE unit is preprogrammed
and not reconfigurable. \texttt{RT\_DRIVE\_CMD} placed at \texttt{0x204} to
avoid collision. \\
\end{longtable}

\begin{longtable}[]{@{}
  >{\raggedright\arraybackslash}p{(\columnwidth - 18\tabcolsep) * \real{0.1081}}
  >{\raggedright\arraybackslash}p{(\columnwidth - 18\tabcolsep) * \real{0.1486}}
  >{\raggedright\arraybackslash}p{(\columnwidth - 18\tabcolsep) * \real{0.0676}}
  >{\raggedright\arraybackslash}p{(\columnwidth - 18\tabcolsep) * \real{0.0811}}
  >{\raggedright\arraybackslash}p{(\columnwidth - 18\tabcolsep) * \real{0.0946}}
  >{\raggedright\arraybackslash}p{(\columnwidth - 18\tabcolsep) * \real{0.1081}}
  >{\raggedright\arraybackslash}p{(\columnwidth - 18\tabcolsep) * \real{0.0676}}
  >{\raggedright\arraybackslash}p{(\columnwidth - 18\tabcolsep) * \real{0.0676}}
  >{\raggedright\arraybackslash}p{(\columnwidth - 18\tabcolsep) * \real{0.0811}}
  >{\raggedright\arraybackslash}p{(\columnwidth - 18\tabcolsep) * \real{0.1757}}@{}}
\toprule\noalign{}
\begin{minipage}[b]{\linewidth}\raggedright
Signal
\end{minipage} & \begin{minipage}[b]{\linewidth}\raggedright
Start bit
\end{minipage} & \begin{minipage}[b]{\linewidth}\raggedright
Len
\end{minipage} & \begin{minipage}[b]{\linewidth}\raggedright
Type
\end{minipage} & \begin{minipage}[b]{\linewidth}\raggedright
Scale
\end{minipage} & \begin{minipage}[b]{\linewidth}\raggedright
Offset
\end{minipage} & \begin{minipage}[b]{\linewidth}\raggedright
Min
\end{minipage} & \begin{minipage}[b]{\linewidth}\raggedright
Max
\end{minipage} & \begin{minipage}[b]{\linewidth}\raggedright
Unit
\end{minipage} & \begin{minipage}[b]{\linewidth}\raggedright
Description
\end{minipage} \\
\midrule\noalign{}
\endhead
\bottomrule\noalign{}
\endlastfoot
\texttt{VCU\_SES\_Alignment\_Enable} & 0 & 1 & bool & 1 & 0 & 0 & 1 & --- & SES
Angle Initial Alignment Enable. 0=disabled, 1=centering. (CSV Row 2) \\
\texttt{VCU\_SES\_Control\_Enable} & 1 & 1 & bool & 1 & 0 & 0 & 1 & --- & VCU
Direction Control Enable. 0=Disabled (Default Assist), 1=Rising Edge Enable
(Angle Control Mode). (CSV Row 3) \\
(unaccounted) & 2 & 6 & --- & --- & --- & --- & --- & --- & Byte 0 bits 2--7 ---
not enumerated in CSV \\
(unaccounted) & --- & 8 & --- & --- & --- & --- & --- & --- & Byte 1 --- not
enumerated in CSV \\
\texttt{VCU\_SES\_Tgt\_StrAngle} & 16 & 16 & i16 & 0.1 & -3000 & -700 & 700 & °
& Target Steering Angle. Negative = left. (CSV Row 4). Note: CSV offset=-3000
(see conversion note below). \\
\texttt{VCU\_SES\_Tgt\_StrAngleSpd} & 32 & 16 & u16 & 1 & 0 & 125 & 525 & °/s &
Target Steering Angle Speed. 16-bit per CSV. Overlaps security signals at byte
5. (CSV Row 5) \\
\texttt{VCU\_SES\_RollCnt\_Enable} & 40 & 1 & bool & 1 & 0 & 0 & 1 & --- & Life
Signal Enable --- \textbf{Must be 1}. Overlaps StrAngleSpd{[}15:8{]}. (CSV Row
6) \\
\texttt{VCU\_SES\_CheckSum\_Enable} & 41 & 1 & bool & 1 & 0 & 0 & 1 & --- &
Checksum Enable --- \textbf{Must be 1}. Overlaps StrAngleSpd{[}15:8{]}. (CSV Row
7) \\
(unaccounted) & 42 & 2 & --- & --- & --- & --- & --- & --- & Byte 5 bits 2--3
--- not enumerated in CSV \\
\texttt{VCU\_SES\_RollCnt} & 44 & 4 & u8 & 1 & 0 & 0 & 15 & --- & Life Signal
rolling counter. Increment every frame. Overlaps StrAngleSpd{[}15:8{]}. (CSV Row
8) \\
\texttt{VCU\_Veh\_Spd\_Value} & 48 & 8 & u8 & 1 & 0 & 0 & 255 & --- & Vehicle
Speed. RT must populate with current speed. (CSV Row 9) \\
\texttt{VCU\_SES\_CheckSum} & 56 & 8 & u8 & 1 & 0 & 0 & 255 & --- & Checksum =
XOR(bytes 0--6) \^{} 0xFF (CSV Row 10) \\
\end{longtable}

\textbf{Byte layout} (little-endian, CSV as source of truth):

\begin{longtable}[]{@{}
  >{\raggedright\arraybackslash}p{(\columnwidth - 16\tabcolsep) * \real{0.2000}}
  >{\raggedright\arraybackslash}p{(\columnwidth - 16\tabcolsep) * \real{0.1000}}
  >{\raggedright\arraybackslash}p{(\columnwidth - 16\tabcolsep) * \real{0.1000}}
  >{\raggedright\arraybackslash}p{(\columnwidth - 16\tabcolsep) * \real{0.1000}}
  >{\raggedright\arraybackslash}p{(\columnwidth - 16\tabcolsep) * \real{0.1000}}
  >{\raggedright\arraybackslash}p{(\columnwidth - 16\tabcolsep) * \real{0.1000}}
  >{\raggedright\arraybackslash}p{(\columnwidth - 16\tabcolsep) * \real{0.1000}}
  >{\raggedright\arraybackslash}p{(\columnwidth - 16\tabcolsep) * \real{0.1000}}
  >{\raggedright\arraybackslash}p{(\columnwidth - 16\tabcolsep) * \real{0.1000}}@{}}
\toprule\noalign{}
\begin{minipage}[b]{\linewidth}\raggedright
Byte
\end{minipage} & \begin{minipage}[b]{\linewidth}\raggedright
0
\end{minipage} & \begin{minipage}[b]{\linewidth}\raggedright
1
\end{minipage} & \begin{minipage}[b]{\linewidth}\raggedright
2
\end{minipage} & \begin{minipage}[b]{\linewidth}\raggedright
3
\end{minipage} & \begin{minipage}[b]{\linewidth}\raggedright
4
\end{minipage} & \begin{minipage}[b]{\linewidth}\raggedright
5
\end{minipage} & \begin{minipage}[b]{\linewidth}\raggedright
6
\end{minipage} & \begin{minipage}[b]{\linewidth}\raggedright
7
\end{minipage} \\
\midrule\noalign{}
\endhead
\bottomrule\noalign{}
\endlastfoot
Content & Align{[}0{]}+CtrlEn{[}1{]} & (unacc.) & Angle {[}7:0{]} & Angle
{[}15:8{]} & Speed {[}7:0{]} & Speed{[}15:8{]} /
RollCntEn{[}0{]}+CksEn{[}1{]}+(gap)+RollCnt{[}4:7{]} (overlap) & Veh\_Spd
{[}7:0{]} & CheckSum \\
\end{longtable}

\begin{quote}
\textbf{Angle conversion (CSV offset=-3000):}
\texttt{physical\_deg\ =\ raw\ ×\ 0.1\ +\ (-3000)}. This encoding places the
±700° physical range at raw values approximately 23000--37000. Raw 0 → -3000°
(outside normal range). This offset is unusual for a signed integer --- 0° does
not map to raw 0. The CSV offset may be a tool artifact; verify against observed
CAN bus values.

\textbf{Byte 5 overlap:} CSV declares \texttt{VCU\_SES\_Tgt\_StrAngleSpd} as
16-bit (bytes 4--5) AND security signals at byte 5 (RollCnt\_Enable,
CheckSum\_Enable, RollCnt). Both are listed in the manufacturer's DBC export.
The EPS-C may internally separate these --- the upper nibble of byte 5 carries
security data while the lower bits carry speed. RT firmware should write the
speed value to bytes 4--5 AND set security fields at byte 5 bits 0--3 + upper
nibble; the EPS-C validates the security portion independently of the speed
portion.

\textbf{Slew rate:} Speed-dependent. RT computes
\texttt{VCU\_SES\_Tgt\_StrAngleSpd} based on speed to ensure smooth steering.
CSV lists range 125--525 °/s. The EPS-C may reject speed commands below 125 °/s.
\end{quote}

\textbf{Internal conversion (architecture, offset=0)}:
\texttt{VCU\_SES\_Tgt\_StrAngle\_raw\ =\ internal\_angle\_mdeg\ /\ 100} (45500
mdeg → 455 raw → 45.5°). CSV declares offset=-3000; if that encoding is used,
the formula would be \texttt{raw\ =\ internal\_angle\_mdeg\ /\ 100\ +\ 30000}
(45500 mdeg → 30455 raw → 45.5°). Verify which encoding the EPS-C actually
expects by observing CAN bus traffic.

\textbf{Security}: If \texttt{roll\_cnt\_enable=0} or
\texttt{checksum\_enable=0}, unit may reject frames. Both must be 1. Checksum
algorithm: \texttt{XOR(bytes{[}0..6{]})\ \^{}\ 0xFF} (verify exact formula
against SYNTREE spec).

\textbf{Slew rate}: Speed-dependent. RT computes
\texttt{VCU\_SES\_Tgt\_StrAngleSpd} based on speed to ensure smooth steering.
Lower speed → lower slew rate for comfort; higher speed → higher slew rate for
responsiveness (within dynamic clamp).

\begin{center}\rule{0.5\linewidth}{0.5pt}\end{center}

\subsection{0x202 --- SES\_ErrInfo (SYNTREE EPS-C Error
Detail)}\label{x202-ses_errinfo-syntree-eps-c-error-detail}

\begin{longtable}[]{@{}ll@{}}
\toprule\noalign{}
Property & Value \\
\midrule\noalign{}
\endhead
\bottomrule\noalign{}
\endlastfoot
\textbf{Sender} & SYNTREE EPS-C (steering module) \\
\textbf{Receiver(s)} & RT ESP32-S3 \\
\textbf{DLC} & 8 \\
\textbf{Period} & 100 ms (10 Hz) \\
\textbf{Endianness} & Motorola LSB (little-endian) \\
\end{longtable}

Detailed fault flags. Each bit is an independent fault indicator. 1 = fault
active.

\begin{longtable}[]{@{}
  >{\raggedright\arraybackslash}p{(\columnwidth - 10\tabcolsep) * \real{0.1429}}
  >{\raggedright\arraybackslash}p{(\columnwidth - 10\tabcolsep) * \real{0.1964}}
  >{\raggedright\arraybackslash}p{(\columnwidth - 10\tabcolsep) * \real{0.0893}}
  >{\raggedright\arraybackslash}p{(\columnwidth - 10\tabcolsep) * \real{0.1071}}
  >{\raggedright\arraybackslash}p{(\columnwidth - 10\tabcolsep) * \real{0.2321}}
  >{\raggedright\arraybackslash}p{(\columnwidth - 10\tabcolsep) * \real{0.2321}}@{}}
\toprule\noalign{}
\begin{minipage}[b]{\linewidth}\raggedright
Signal
\end{minipage} & \begin{minipage}[b]{\linewidth}\raggedright
Start bit
\end{minipage} & \begin{minipage}[b]{\linewidth}\raggedright
Len
\end{minipage} & \begin{minipage}[b]{\linewidth}\raggedright
Type
\end{minipage} & \begin{minipage}[b]{\linewidth}\raggedright
Fault Level
\end{minipage} & \begin{minipage}[b]{\linewidth}\raggedright
Description
\end{minipage} \\
\midrule\noalign{}
\endhead
\bottomrule\noalign{}
\endlastfoot
\texttt{SES\_ECUUnderVolt\_Err} & 0 & 1 & bool & L2 & Controller Under
Voltage \\
\texttt{SES\_ECUOverVolt\_Err} & 1 & 1 & bool & L2 & Controller Over Voltage \\
\texttt{SES\_CanCom\_Err} & 2 & 1 & bool & L1 & CAN Communication Fault \\
\texttt{SES\_ECUTemp\_Err} & 3 & 1 & bool & L1 & Controller Temp Fault \\
\texttt{SES\_Domain\_drive\_SC\_Err} & 4 & 1 & bool & L2 & Domain Drive Short
Circuit \\
\texttt{SES\_Domain\_drive\_V\_Err} & 5 & 1 & bool & L2 & Domain Drive Voltage
Fault \\
\texttt{SES\_Domain\_drive\_T\_Err} & 6 & 1 & bool & L2 & Domain Drive
Temperature Fault \\
\texttt{SES\_TempSensor\_Err} & 7 & 1 & bool & --- & Temperature Sensor Fault \\
\texttt{SES\_AngleSensor\_P\_OC\_Err} & 8 & 1 & bool & \textbf{L3} & Angle
Sensor Pri. Open Circuit \\
\texttt{SES\_AngleSensor\_P\_AF\_Err} & 9 & 1 & bool & \textbf{L3} & Angle
Sensor Pri. Out of Range \\
\texttt{SES\_AngleSensor\_S\_OC\_Err} & 10 & 1 & bool & \textbf{L3} & Angle
Sensor Sec. Open Circuit \\
\texttt{SES\_AngleSensor\_S\_AF\_Err} & 11 & 1 & bool & \textbf{L3} & Angle
Sensor Sec. Out of Range \\
\texttt{SES\_SensorPow\_Err} & 12 & 1 & bool & L2 & Sensor Power Fault \\
\texttt{SES\_Alignment\_Err} & 13 & 1 & bool & L1 & Centering Fault \\
\texttt{SES\_OverAngle\_Err} & 14 & 1 & bool & L2 & Over Angle Fault \\
\texttt{SES\_StrMtr\_Stall\_Err} & 15 & 1 & bool & L1 & Motor Stall Fault \\
\texttt{SES\_MtrCurt\_Err} & 16 & 1 & bool & L2 & Motor Current Fault \\
\texttt{SES\_SensorCL\_Err} & 17 & 1 & bool & L2 & Sensor 5V Power 1/2 Fault \\
\texttt{SES\_TorqSensor\_T1\_OC\_Err} & 18 & 1 & bool & \textbf{L3} & Torque
Sensor T1 Open Circuit \\
\texttt{SES\_TorqSensor\_T1\_AF\_Err} & 19 & 1 & bool & \textbf{L3} & Torque
Sensor T1 Out of Range \\
\texttt{SES\_TorqSensor\_T2\_OC\_Err} & 20 & 1 & bool & \textbf{L3} & Torque
Sensor T2 Open Circuit \\
\texttt{SES\_TorqSensor\_T2\_AF\_Err} & 21 & 1 & bool & \textbf{L3} & Torque
Sensor T2 Out of Range \\
\texttt{SentAngle\_Err} & 22 & 1 & bool & L1 & Angle Error \\
\texttt{SES\_StrMtr\_Idling\_Err} & 23 & 1 & bool & L2 & Motor Idling Fault \\
\texttt{SES\_EPROM\_Err} & 24 & 1 & bool & L2 & EEPROM Fault \\
(reserved) & 25 & 31 & --- & --- & Bits 25--55 (byte 3 bits 1--7 + bytes 4--6).
Not enumerated in CSV. \\
\texttt{SES\_Veh\_Spd\_Value} & 56 & 8 & u8 & --- & Vehicle speed at fault
snapshot (CSV Row 46) \\
\end{longtable}

\textbf{Byte layout} (little-endian):

\begin{longtable}[]{@{}
  >{\raggedright\arraybackslash}p{(\columnwidth - 16\tabcolsep) * \real{0.2000}}
  >{\raggedright\arraybackslash}p{(\columnwidth - 16\tabcolsep) * \real{0.1000}}
  >{\raggedright\arraybackslash}p{(\columnwidth - 16\tabcolsep) * \real{0.1000}}
  >{\raggedright\arraybackslash}p{(\columnwidth - 16\tabcolsep) * \real{0.1000}}
  >{\raggedright\arraybackslash}p{(\columnwidth - 16\tabcolsep) * \real{0.1000}}
  >{\raggedright\arraybackslash}p{(\columnwidth - 16\tabcolsep) * \real{0.1000}}
  >{\raggedright\arraybackslash}p{(\columnwidth - 16\tabcolsep) * \real{0.1000}}
  >{\raggedright\arraybackslash}p{(\columnwidth - 16\tabcolsep) * \real{0.1000}}
  >{\raggedright\arraybackslash}p{(\columnwidth - 16\tabcolsep) * \real{0.1000}}@{}}
\toprule\noalign{}
\begin{minipage}[b]{\linewidth}\raggedright
Byte
\end{minipage} & \begin{minipage}[b]{\linewidth}\raggedright
0
\end{minipage} & \begin{minipage}[b]{\linewidth}\raggedright
1
\end{minipage} & \begin{minipage}[b]{\linewidth}\raggedright
2
\end{minipage} & \begin{minipage}[b]{\linewidth}\raggedright
3
\end{minipage} & \begin{minipage}[b]{\linewidth}\raggedright
4
\end{minipage} & \begin{minipage}[b]{\linewidth}\raggedright
5
\end{minipage} & \begin{minipage}[b]{\linewidth}\raggedright
6
\end{minipage} & \begin{minipage}[b]{\linewidth}\raggedright
7
\end{minipage} \\
\midrule\noalign{}
\endhead
\bottomrule\noalign{}
\endlastfoot
Content & faults {[}7:0{]} & faults {[}15:8{]} & faults {[}23:16{]} & faults
{[}31:24{]} & rsvd & rsvd & rsvd & \texttt{SES\_Veh\_Spd\_Value} \\
\end{longtable}

\begin{quote}
\textbf{Safety note:} 8 faults are L3 (angle sensor primary/secondary
open-circuit/out-of-range × 4 + torque sensor T1/T2 open-circuit/out-of-range ×
4). RT must subscribe and escalate L3 faults to ESTOP. The EPS-C has
dual-redundant angle sensors (P/S) and dual torque sensors (T1/T2) --- L3 on
either channel of a redundant pair indicates critical sensor failure. Note fault
levels differ from SEB: steering CAN comms is L1 (minor) vs brake CAN comms L3
(severe) --- steering comm loss is less immediately dangerous than brake comm
loss.
\end{quote}

\begin{quote}
⚠️ \textbf{CSV Row 40--41 description swap:} Row 40 signal
\texttt{SES\_TorqSensor\_T2\_OC\_Err} is labeled ``Torque Sensor T1 Out of
Range'' --- descriptions appear swapped. Signal names used above are
authoritative; descriptions are corrected.
\end{quote}

\begin{center}\rule{0.5\linewidth}{0.5pt}\end{center}

\subsection{0x203 --- SES\_Version (SYNTREE EPS-C Firmware
Version)}\label{x203-ses_version-syntree-eps-c-firmware-version}

\begin{longtable}[]{@{}ll@{}}
\toprule\noalign{}
Property & Value \\
\midrule\noalign{}
\endhead
\bottomrule\noalign{}
\endlastfoot
\textbf{Sender} & SYNTREE EPS-C (steering module) \\
\textbf{Receiver(s)} & RT ESP32-S3 \\
\textbf{DLC} & 8 \\
\textbf{Period} & 1000 ms (1 Hz) \\
\textbf{Endianness} & Motorola LSB (little-endian) \\
\end{longtable}

\begin{longtable}[]{@{}
  >{\raggedright\arraybackslash}p{(\columnwidth - 18\tabcolsep) * \real{0.1081}}
  >{\raggedright\arraybackslash}p{(\columnwidth - 18\tabcolsep) * \real{0.1486}}
  >{\raggedright\arraybackslash}p{(\columnwidth - 18\tabcolsep) * \real{0.0676}}
  >{\raggedright\arraybackslash}p{(\columnwidth - 18\tabcolsep) * \real{0.0811}}
  >{\raggedright\arraybackslash}p{(\columnwidth - 18\tabcolsep) * \real{0.0946}}
  >{\raggedright\arraybackslash}p{(\columnwidth - 18\tabcolsep) * \real{0.1081}}
  >{\raggedright\arraybackslash}p{(\columnwidth - 18\tabcolsep) * \real{0.0676}}
  >{\raggedright\arraybackslash}p{(\columnwidth - 18\tabcolsep) * \real{0.0676}}
  >{\raggedright\arraybackslash}p{(\columnwidth - 18\tabcolsep) * \real{0.0811}}
  >{\raggedright\arraybackslash}p{(\columnwidth - 18\tabcolsep) * \real{0.1757}}@{}}
\toprule\noalign{}
\begin{minipage}[b]{\linewidth}\raggedright
Signal
\end{minipage} & \begin{minipage}[b]{\linewidth}\raggedright
Start bit
\end{minipage} & \begin{minipage}[b]{\linewidth}\raggedright
Len
\end{minipage} & \begin{minipage}[b]{\linewidth}\raggedright
Type
\end{minipage} & \begin{minipage}[b]{\linewidth}\raggedright
Scale
\end{minipage} & \begin{minipage}[b]{\linewidth}\raggedright
Offset
\end{minipage} & \begin{minipage}[b]{\linewidth}\raggedright
Min
\end{minipage} & \begin{minipage}[b]{\linewidth}\raggedright
Max
\end{minipage} & \begin{minipage}[b]{\linewidth}\raggedright
Unit
\end{minipage} & \begin{minipage}[b]{\linewidth}\raggedright
Description
\end{minipage} \\
\midrule\noalign{}
\endhead
\bottomrule\noalign{}
\endlastfoot
\texttt{SES\_SW\_Version} & 0 & 8 & u8 & 0.01 & 0 & 0 & 255 & --- & Software
version (e.g., 0x64 = 1.00) \\
\texttt{SES\_HW\_Version} & 8 & 8 & u8 & 0.1 & 0 & 0 & 25.5 & --- & Hardware
version (e.g., 0x0D = 1.3) \\
(reserved) & 16 & 48 & --- & --- & --- & --- & --- & --- & Bytes 2--7 (code also
reads bytes 2-3 as extended HW version info --- minor logging discrepancy) \\
\end{longtable}

\textbf{RT usage}: Log on boot for compatibility check
(\texttt{SW=\%02X.\%02X\ HW=\%02X.\%02X}). Report via telemetry to Jetson.

\begin{center}\rule{0.5\linewidth}{0.5pt}\end{center}

\subsection{0x6FA --- SES\_Test (SYNTREE EPS-C
Telemetry)}\label{x6fa-ses_test-syntree-eps-c-telemetry}

\begin{longtable}[]{@{}ll@{}}
\toprule\noalign{}
Property & Value \\
\midrule\noalign{}
\endhead
\bottomrule\noalign{}
\endlastfoot
\textbf{Sender} & SYNTREE EPS-C (steering module) \\
\textbf{Receiver(s)} & RT ESP32-S3 \\
\textbf{DLC} & 8 \\
\textbf{Period} & 10 ms (100 Hz) \\
\textbf{Endianness} & Motorola LSB (little-endian) \\
\end{longtable}

\begin{longtable}[]{@{}
  >{\raggedright\arraybackslash}p{(\columnwidth - 18\tabcolsep) * \real{0.1081}}
  >{\raggedright\arraybackslash}p{(\columnwidth - 18\tabcolsep) * \real{0.1486}}
  >{\raggedright\arraybackslash}p{(\columnwidth - 18\tabcolsep) * \real{0.0676}}
  >{\raggedright\arraybackslash}p{(\columnwidth - 18\tabcolsep) * \real{0.0811}}
  >{\raggedright\arraybackslash}p{(\columnwidth - 18\tabcolsep) * \real{0.0946}}
  >{\raggedright\arraybackslash}p{(\columnwidth - 18\tabcolsep) * \real{0.1081}}
  >{\raggedright\arraybackslash}p{(\columnwidth - 18\tabcolsep) * \real{0.0676}}
  >{\raggedright\arraybackslash}p{(\columnwidth - 18\tabcolsep) * \real{0.0676}}
  >{\raggedright\arraybackslash}p{(\columnwidth - 18\tabcolsep) * \real{0.0811}}
  >{\raggedright\arraybackslash}p{(\columnwidth - 18\tabcolsep) * \real{0.1757}}@{}}
\toprule\noalign{}
\begin{minipage}[b]{\linewidth}\raggedright
Signal
\end{minipage} & \begin{minipage}[b]{\linewidth}\raggedright
Start bit
\end{minipage} & \begin{minipage}[b]{\linewidth}\raggedright
Len
\end{minipage} & \begin{minipage}[b]{\linewidth}\raggedright
Type
\end{minipage} & \begin{minipage}[b]{\linewidth}\raggedright
Scale
\end{minipage} & \begin{minipage}[b]{\linewidth}\raggedright
Offset
\end{minipage} & \begin{minipage}[b]{\linewidth}\raggedright
Min
\end{minipage} & \begin{minipage}[b]{\linewidth}\raggedright
Max
\end{minipage} & \begin{minipage}[b]{\linewidth}\raggedright
Unit
\end{minipage} & \begin{minipage}[b]{\linewidth}\raggedright
Description
\end{minipage} \\
\midrule\noalign{}
\endhead
\bottomrule\noalign{}
\endlastfoot
(reserved) & 0 & 8 & --- & --- & --- & --- & --- & --- & Byte 0 \\
\texttt{SES\_MtrCurt} & 8 & 16 & i16 & 0.0078125 & 0 & 0 & 60 & A & Motor
current. Bytes 1--2. Narrower range than brake SEB\_Test (±255A). \\
\texttt{SES\_ECUTemp} & 24 & 16 & u16 & 0.5 & 0 & 0 & 255 & °C & ECU
temperature. Bytes 3--4. \\
\texttt{SES\_PowVolt} & 40 & 16 & u16 & 0.00390625 & 0 & 0 & 18 & V & Supply
voltage. Bytes 5--6. Narrower range than brake (0--18V vs 0--32V). \\
(reserved) & 56 & 8 & --- & --- & --- & --- & --- & --- & Byte 7 \\
\end{longtable}

\begin{quote}
\textbf{Note:} CSV uses byte-local start-bit numbering for this message.
Converted to absolute Motorola LSB above. Range limits differ from brake
SEB\_Test (0x6FB): steering motor current 0--60A vs brake ±255A; steering supply
voltage 0--18V vs brake 0--32V. \textbf{RT usage:} Monitor \texttt{SES\_MtrCurt}
for mechanical binding / rack damage. \texttt{SES\_ECUTemp} for thermal
throttling.
\end{quote}

\begin{center}\rule{0.5\linewidth}{0.5pt}\end{center}

\subsection{0x302 --- HOST\_LIGHT\_CMD
(forwarded)}\label{x302-host_light_cmd-forwarded}

\begin{longtable}[]{@{}ll@{}}
\toprule\noalign{}
Property & Value \\
\midrule\noalign{}
\endhead
\bottomrule\noalign{}
\endlastfoot
\textbf{Sender} & RT (fwd from Jetson) \\
\textbf{Receiver(s)} & SYS \\
\textbf{DLC} & 1 \\
\textbf{Period} & On change \\
\end{longtable}

\begin{longtable}[]{@{}llll@{}}
\toprule\noalign{}
Signal & Start bit & Len & Type \\
\midrule\noalign{}
\endhead
\bottomrule\noalign{}
\endlastfoot
\texttt{HOST\_LeftTurn} & 0 & 1 & bool \\
\texttt{HOST\_RightTurn} & 1 & 1 & bool \\
\texttt{HOST\_BrakeLight} & 2 & 1 & bool \\
\texttt{HOST\_Headlight} & 3 & 1 & bool \\
(reserved) & 4 & 4 & --- \\
\end{longtable}

\begin{center}\rule{0.5\linewidth}{0.5pt}\end{center}

\subsection{0x600 --- SYS\_DIAG\_RPT}\label{x600-sys_diag_rpt}

\begin{longtable}[]{@{}ll@{}}
\toprule\noalign{}
Property & Value \\
\midrule\noalign{}
\endhead
\bottomrule\noalign{}
\endlastfoot
\textbf{Sender} & SYS \\
\textbf{Receiver(s)} & RT (→ Jetson) \\
\textbf{DLC} & 8 \\
\textbf{Period} & 1 Hz \\
\end{longtable}

\begin{longtable}[]{@{}llll@{}}
\toprule\noalign{}
Signal & Start bit & Len & Type \\
\midrule\noalign{}
\endhead
\bottomrule\noalign{}
\endlastfoot
\texttt{SYS\_DiagMode} & 0 & 8 & u8 \\
\texttt{SYS\_DiagBrakeEngaged} & 8 & 8 & u8 \\
\texttt{SYS\_DiagHeartbeatOk} & 16 & 8 & u8 \\
\texttt{SYS\_DiagEstopActive} & 24 & 8 & u8 \\
\texttt{SYS\_DiagFreeHeapKb} & 32 & 16 & u16 \\
\texttt{SYS\_DiagTec} & 48 & 8 & u8 \\
\texttt{SYS\_DiagRec} & 56 & 8 & u8 \\
\end{longtable}

Byte layout (big-endian): Byte 0=mode, 1=brake, 2=hb\_ok, 3=estop, 4-5=heap,
6=tec, 7=rec.

\begin{center}\rule{0.5\linewidth}{0.5pt}\end{center}

\subsection{0x7B9 --- VCU\_SEB\_REQ (SYNTREE SEB Brake
Command)}\label{x7b9-vcu_seb_req-syntree-seb-brake-command}

\begin{longtable}[]{@{}ll@{}}
\toprule\noalign{}
Property & Value \\
\midrule\noalign{}
\endhead
\bottomrule\noalign{}
\endlastfoot
\textbf{Sender} & SYS ESP32-S3 \\
\textbf{Receiver(s)} & SYNTREE SEB (brake module) \\
\textbf{DLC} & 8 \\
\textbf{Period} & 20 ms (50 Hz) --- \textbf{continuous, every frame} \\
\textbf{Endianness} & Motorola LSB (little-endian) \\
\end{longtable}

\begin{longtable}[]{@{}
  >{\raggedright\arraybackslash}p{(\columnwidth - 18\tabcolsep) * \real{0.1081}}
  >{\raggedright\arraybackslash}p{(\columnwidth - 18\tabcolsep) * \real{0.1486}}
  >{\raggedright\arraybackslash}p{(\columnwidth - 18\tabcolsep) * \real{0.0676}}
  >{\raggedright\arraybackslash}p{(\columnwidth - 18\tabcolsep) * \real{0.0811}}
  >{\raggedright\arraybackslash}p{(\columnwidth - 18\tabcolsep) * \real{0.0946}}
  >{\raggedright\arraybackslash}p{(\columnwidth - 18\tabcolsep) * \real{0.1081}}
  >{\raggedright\arraybackslash}p{(\columnwidth - 18\tabcolsep) * \real{0.0676}}
  >{\raggedright\arraybackslash}p{(\columnwidth - 18\tabcolsep) * \real{0.0676}}
  >{\raggedright\arraybackslash}p{(\columnwidth - 18\tabcolsep) * \real{0.0811}}
  >{\raggedright\arraybackslash}p{(\columnwidth - 18\tabcolsep) * \real{0.1757}}@{}}
\toprule\noalign{}
\begin{minipage}[b]{\linewidth}\raggedright
Signal
\end{minipage} & \begin{minipage}[b]{\linewidth}\raggedright
Start bit
\end{minipage} & \begin{minipage}[b]{\linewidth}\raggedright
Len
\end{minipage} & \begin{minipage}[b]{\linewidth}\raggedright
Type
\end{minipage} & \begin{minipage}[b]{\linewidth}\raggedright
Scale
\end{minipage} & \begin{minipage}[b]{\linewidth}\raggedright
Offset
\end{minipage} & \begin{minipage}[b]{\linewidth}\raggedright
Min
\end{minipage} & \begin{minipage}[b]{\linewidth}\raggedright
Max
\end{minipage} & \begin{minipage}[b]{\linewidth}\raggedright
Unit
\end{minipage} & \begin{minipage}[b]{\linewidth}\raggedright
Description
\end{minipage} \\
\midrule\noalign{}
\endhead
\bottomrule\noalign{}
\endlastfoot
\texttt{VCU\_SEB\_Alignment\_Enable} & 0 & 1 & bool & 1 & 0 & 0 & 1 & --- &
Calibration enable (CSV Row 2) \\
\texttt{VCU\_SEB\_Control\_Enable} & 1 & 1 & bool & 1 & 0 & 0 & 1 & --- & Active
control enable (CSV Row 3) \\
\texttt{VCU\_SEB\_Control\_Mode} & 2 & 1 & bool & 1 & 0 & 0 & 1 & enum &
0=Stroke, 1=Pressure (CSV Row 4) \\
\texttt{VCU\_SEB\_AutoBrake} & 3 & 1 & bool & 1 & 0 & 0 & 1 & --- & Auto-brake /
emergency trigger (CSV Row 5) \\
(unaccounted) & 4 & 4 & --- & --- & --- & --- & --- & --- & Byte 0 bits 4--7 ---
not enumerated in CSV \\
(unaccounted) & --- & 8 & --- & --- & --- & --- & --- & --- & Byte 1 --- not
enumerated in CSV \\
\texttt{VCU\_SEB\_Stroke\_Value\_Req} & 16 & 16 & u16 & 0.05 & -30 & -5 & 27 &
mm & Requested stroke position (CSV Row 6) \\
\texttt{VCU\_SEB\_Pre\_Value\_Req} & 24 & 8 & u8 & 0.05 & 0 & 0 & 5 & MPa &
Requested pressure (CSV Row 7). Raw = kPa × 0.02. Overlaps Stroke{[}15:8{]} at
byte 3 --- mode-dependent: Stroke uses full 16-bit in Mode 0, byte 3 carries
pressure in Mode 1. \\
(unaccounted) & 32 & 16 & --- & --- & --- & --- & --- & --- & Bytes 4--5 --- not
enumerated in CSV \\
\texttt{VCU\_SEB\_RollCnt\_Enable} & 48 & 1 & bool & 1 & 0 & 0 & 1 & --- & Life
Signal Validity --- \textbf{Must be 1} (CSV Row 8) \\
\texttt{VCU\_SEB\_CheckSum\_Enable} & 49 & 1 & bool & 1 & 0 & 0 & 1 & --- &
Checksum Validity --- \textbf{Must be 1} (CSV Row 9) \\
(unaccounted) & 50 & 2 & --- & --- & --- & --- & --- & --- & Byte 6 bits 2--3
--- not enumerated in CSV \\
\texttt{VCU\_SEB\_RollCnt} & 52 & 4 & u8 & 1 & 0 & 0 & 15 & --- & Life Signal
rolling counter. Increment every frame. (CSV Row 10) \\
\texttt{VCU\_SEB\_CheckSum} & 56 & 8 & u8 & 1 & 0 & 0 & 255 & --- & Checksum =
XOR(bytes 0--6) \^{} 0xFF (CSV Row 11) \\
\end{longtable}

\textbf{Byte layout} (little-endian, CSV as source of truth):

\begin{longtable}[]{@{}
  >{\raggedright\arraybackslash}p{(\columnwidth - 16\tabcolsep) * \real{0.2000}}
  >{\raggedright\arraybackslash}p{(\columnwidth - 16\tabcolsep) * \real{0.1000}}
  >{\raggedright\arraybackslash}p{(\columnwidth - 16\tabcolsep) * \real{0.1000}}
  >{\raggedright\arraybackslash}p{(\columnwidth - 16\tabcolsep) * \real{0.1000}}
  >{\raggedright\arraybackslash}p{(\columnwidth - 16\tabcolsep) * \real{0.1000}}
  >{\raggedright\arraybackslash}p{(\columnwidth - 16\tabcolsep) * \real{0.1000}}
  >{\raggedright\arraybackslash}p{(\columnwidth - 16\tabcolsep) * \real{0.1000}}
  >{\raggedright\arraybackslash}p{(\columnwidth - 16\tabcolsep) * \real{0.1000}}
  >{\raggedright\arraybackslash}p{(\columnwidth - 16\tabcolsep) * \real{0.1000}}@{}}
\toprule\noalign{}
\begin{minipage}[b]{\linewidth}\raggedright
Byte
\end{minipage} & \begin{minipage}[b]{\linewidth}\raggedright
0
\end{minipage} & \begin{minipage}[b]{\linewidth}\raggedright
1
\end{minipage} & \begin{minipage}[b]{\linewidth}\raggedright
2
\end{minipage} & \begin{minipage}[b]{\linewidth}\raggedright
3
\end{minipage} & \begin{minipage}[b]{\linewidth}\raggedright
4
\end{minipage} & \begin{minipage}[b]{\linewidth}\raggedright
5
\end{minipage} & \begin{minipage}[b]{\linewidth}\raggedright
6
\end{minipage} & \begin{minipage}[b]{\linewidth}\raggedright
7
\end{minipage} \\
\midrule\noalign{}
\endhead
\bottomrule\noalign{}
\endlastfoot
Content & Align{[}0{]}+CtrlEn{[}1{]}+Mode{[}2{]}+AutoBrk{[}3{]} & (unaccounted)
& Stroke\_Req {[}7:0{]} & Stroke\_Req {[}15:8{]} / Pre\_Req {[}7:0{]}
(mode-muxed) & (unaccounted) & (unaccounted) &
RollCntEn{[}0{]}+CksEn{[}1{]}+(gap)+RollCnt{[}4:7{]} & CheckSum \\
\end{longtable}

\begin{quote}
\textbf{Byte 3 multiplexing:} In Stroke Mode (Mode=0), bytes 2--3 carry the
16-bit stroke value (\texttt{VCU\_SEB\_Stroke\_Value\_Req}). In Pressure Mode
(Mode=1), byte 3 carries the 8-bit pressure value
(\texttt{VCU\_SEB\_Pre\_Value\_Req}). Both signals are declared in the CSV at
overlapping positions --- the SEB interprets byte 3 based on the active mode
bit. Bytes 1, 4, and 5 are not enumerated in the CSV; the SEB may ignore them or
use them for undocumented functions.
\end{quote}

\textbf{Stroke conversion}: \texttt{raw\ =\ (physical\_mm\ +\ 30.0)\ /\ 0.05}

\begin{longtable}[]{@{}lll@{}}
\toprule\noalign{}
Physical & Raw & Use case \\
\midrule\noalign{}
\endhead
\bottomrule\noalign{}
\endlastfoot
-5 mm & 500 & Min \\
0 mm & 600 & Released \\
15 mm & 900 & Manual lever pressed \\
27 mm & 1140 & ESTOP full brake \\
\end{longtable}

\textbf{Security}: Rolling counter must increment 0→15 every frame. Same value
twice → SEB rejects (assumes frozen controller). Checksum =
\texttt{XOR(bytes{[}0..6{]})\ \^{}\ 0xFF} (verify against SYNTREE spec).

\textbf{Mode 0 (Stroke)}: Command a specific pushrod position in mm. Best for
mimicking pedal travel / ESTOP full brake / manual lever. \textbf{Mode 1
(Pressure)}: Command hydraulic pressure in MPa. SEB's internal PID maintains
target. Best for autonomous deceleration control (compensates for pad wear,
temperature).

\begin{center}\rule{0.5\linewidth}{0.5pt}\end{center}

\subsection{0x721 --- SEB\_STATUS (SYNTREE SEB Brake
Feedback)}\label{x721-seb_status-syntree-seb-brake-feedback}

\begin{longtable}[]{@{}ll@{}}
\toprule\noalign{}
Property & Value \\
\midrule\noalign{}
\endhead
\bottomrule\noalign{}
\endlastfoot
\textbf{Sender} & SYNTREE SEB (brake module) \\
\textbf{Receiver(s)} & SYS ESP32-S3 \\
\textbf{DLC} & 8 \\
\textbf{Period} & 10 ms (100 Hz) \\
\textbf{Endianness} & Motorola LSB (little-endian) \\
\end{longtable}

\begin{longtable}[]{@{}
  >{\raggedright\arraybackslash}p{(\columnwidth - 18\tabcolsep) * \real{0.1081}}
  >{\raggedright\arraybackslash}p{(\columnwidth - 18\tabcolsep) * \real{0.1486}}
  >{\raggedright\arraybackslash}p{(\columnwidth - 18\tabcolsep) * \real{0.0676}}
  >{\raggedright\arraybackslash}p{(\columnwidth - 18\tabcolsep) * \real{0.0811}}
  >{\raggedright\arraybackslash}p{(\columnwidth - 18\tabcolsep) * \real{0.0946}}
  >{\raggedright\arraybackslash}p{(\columnwidth - 18\tabcolsep) * \real{0.1081}}
  >{\raggedright\arraybackslash}p{(\columnwidth - 18\tabcolsep) * \real{0.0676}}
  >{\raggedright\arraybackslash}p{(\columnwidth - 18\tabcolsep) * \real{0.0676}}
  >{\raggedright\arraybackslash}p{(\columnwidth - 18\tabcolsep) * \real{0.0811}}
  >{\raggedright\arraybackslash}p{(\columnwidth - 18\tabcolsep) * \real{0.1757}}@{}}
\toprule\noalign{}
\begin{minipage}[b]{\linewidth}\raggedright
Signal
\end{minipage} & \begin{minipage}[b]{\linewidth}\raggedright
Start bit
\end{minipage} & \begin{minipage}[b]{\linewidth}\raggedright
Len
\end{minipage} & \begin{minipage}[b]{\linewidth}\raggedright
Type
\end{minipage} & \begin{minipage}[b]{\linewidth}\raggedright
Scale
\end{minipage} & \begin{minipage}[b]{\linewidth}\raggedright
Offset
\end{minipage} & \begin{minipage}[b]{\linewidth}\raggedright
Min
\end{minipage} & \begin{minipage}[b]{\linewidth}\raggedright
Max
\end{minipage} & \begin{minipage}[b]{\linewidth}\raggedright
Unit
\end{minipage} & \begin{minipage}[b]{\linewidth}\raggedright
Description
\end{minipage} \\
\midrule\noalign{}
\endhead
\bottomrule\noalign{}
\endlastfoot
\texttt{SEB\_Alignment\_Status} & 0 & 1 & bool & 1 & 0 & 0 & 1 & --- & Alignment
Info Feedback. 1 = aligned. (CSV Row 12) \\
\texttt{SEB\_Control\_Enable\_Status} & 1 & 1 & bool & 1 & 0 & 0 & 1 & --- &
Control Enable Feedback (CSV Row 13) \\
\texttt{SEB\_Control\_Mode\_Status} & 2 & 2 & u8 & 1 & 0 & 0 & 3 & enum &
Control Mode Feedback: 0=?, 1=Stroke?, 2=Pressure?, 3=? (CSV Row 14) \\
\texttt{SEB\_AutoBrake\_Status} & 4 & 1 & bool & 1 & 0 & 0 & 1 & --- & Auto
Brake Status Feedback (CSV Row 15) \\
(unaccounted) & 5 & 1 & --- & --- & --- & --- & --- & --- & Byte 0 bit 5 --- not
enumerated in CSV \\
\texttt{SEB\_Error\_Status} & 6 & 2 & u8 & 1 & 0 & 0 & 3 & enum & 0=No fault,
1=L1 minor, 2=L2 general, 3=L3 severe (CSV Row 16) \\
(unaccounted) & --- & 8 & --- & --- & --- & --- & --- & --- & Byte 1 --- not
enumerated in CSV \\
\texttt{SEB\_Stroke\_Value} & 16 & 16 & u16 & 0.05 & -30 & -5 & 27 & mm & Stroke
Value Feedback (CSV Row 17) \\
\texttt{SEB\_Pressure\_Value} & 24 & 8 & u8 & 0.05 & 0 & 0 & 5 & MPa & Pressure
Value Feedback (CSV Row 18). Overlaps Stroke{[}15:8{]} at byte 3 ---
mode-dependent. \\
\texttt{SEB\_Angle\_Value} & 40 & 16 & i16 & 0.5 & 0 & -150 & 840 & --- & Angle
Feedback (CSV Row 19). Overlaps security echo bits at byte 6 --- see note
below. \\
\texttt{SEB\_RollCnt\_Enable\_Status} & 48 & 1 & bool & 1 & 0 & 0 & 1 & --- &
Life Signal Status Feedback (CSV Row 20). Overlaps Angle\_Value{[}15:8{]}. \\
\texttt{SEB\_CheckSum\_Enable\_Status} & 49 & 1 & bool & 1 & 0 & 0 & 1 & --- &
Checksum Status Feedback (CSV Row 21). Overlaps Angle\_Value{[}15:8{]}. \\
(unaccounted) & 50 & 2 & --- & --- & --- & --- & --- & --- & Byte 6 bits 2--3
--- not enumerated in CSV \\
\texttt{SEB\_RollCnt\_Status} & 52 & 4 & u8 & 1 & 0 & 0 & 15 & --- & Life Signal
Feedback --- echoes received rolling counter (CSV Row 22) \\
\texttt{SEB\_CheckSum\_Status} & 56 & 8 & u8 & 1 & 0 & 0 & 255 & --- & Checksum
Feedback (CSV Row 23) \\
\end{longtable}

\textbf{Byte layout} (little-endian, CSV as source of truth):

\begin{longtable}[]{@{}
  >{\raggedright\arraybackslash}p{(\columnwidth - 16\tabcolsep) * \real{0.2000}}
  >{\raggedright\arraybackslash}p{(\columnwidth - 16\tabcolsep) * \real{0.1000}}
  >{\raggedright\arraybackslash}p{(\columnwidth - 16\tabcolsep) * \real{0.1000}}
  >{\raggedright\arraybackslash}p{(\columnwidth - 16\tabcolsep) * \real{0.1000}}
  >{\raggedright\arraybackslash}p{(\columnwidth - 16\tabcolsep) * \real{0.1000}}
  >{\raggedright\arraybackslash}p{(\columnwidth - 16\tabcolsep) * \real{0.1000}}
  >{\raggedright\arraybackslash}p{(\columnwidth - 16\tabcolsep) * \real{0.1000}}
  >{\raggedright\arraybackslash}p{(\columnwidth - 16\tabcolsep) * \real{0.1000}}
  >{\raggedright\arraybackslash}p{(\columnwidth - 16\tabcolsep) * \real{0.1000}}@{}}
\toprule\noalign{}
\begin{minipage}[b]{\linewidth}\raggedright
Byte
\end{minipage} & \begin{minipage}[b]{\linewidth}\raggedright
0
\end{minipage} & \begin{minipage}[b]{\linewidth}\raggedright
1
\end{minipage} & \begin{minipage}[b]{\linewidth}\raggedright
2
\end{minipage} & \begin{minipage}[b]{\linewidth}\raggedright
3
\end{minipage} & \begin{minipage}[b]{\linewidth}\raggedright
4
\end{minipage} & \begin{minipage}[b]{\linewidth}\raggedright
5
\end{minipage} & \begin{minipage}[b]{\linewidth}\raggedright
6
\end{minipage} & \begin{minipage}[b]{\linewidth}\raggedright
7
\end{minipage} \\
\midrule\noalign{}
\endhead
\bottomrule\noalign{}
\endlastfoot
Content &
Align{[}0{]}+CtrlEn{[}1{]}+Mode{[}2:3{]}+AutoBrk{[}4{]}+(gap)+Error{[}6:7{]} &
(unacc.) & Stroke {[}7:0{]} & Stroke {[}15:8{]} / Pressure {[}7:0{]}
(mode-muxed) & (unacc.) & Angle {[}7:0{]} & Angle{[}15:8{]} /
RollCntEn{[}0{]}+CksEn{[}1{]}+(gap)+RollCnt{[}4:7{]} (overlap) & CksSum\_Stat \\
\end{longtable}

\begin{quote}
\textbf{Byte 3 multiplexing:} Same pattern as command frame ---
\texttt{SEB\_Stroke\_Value} uses full 16-bit at bytes 2--3 in Stroke Mode;
\texttt{SEB\_Pressure\_Value} uses byte 3 in Pressure Mode. The SEB reports
whichever is active.

\textbf{Byte 6 overlap:} CSV lists \texttt{SEB\_Angle\_Value} as 16-bit (bytes
5--6) AND security echo bits at byte 6 (bits 48--49, 52--55). These overlap. The
CSV (manufacturer DBC export) declares both --- they may represent different
firmware versions or the Angle\_Value may be 8-bit in practice (byte 5 only).
Trust the CSV's declaration and handle in firmware by reading Angle as 16-bit
from bytes 5--6, understanding that the upper byte may carry security echo data
in some SEB firmware revisions.
\end{quote}

\textbf{SYS usage}: Boot sync --- read \texttt{SEB\_Stroke\_Value} as initial
command target. Active --- confirm \texttt{SEB\_Alignment\_Status\ ==\ 1}.
\texttt{SEB\_Error\_Status\ \textgreater{}\ 0} → log and report via
\texttt{0x011}. Subscribe to \texttt{0x731\ SEB\_ErrInfo} for detailed fault
flags --- escalate L3 faults to ESTOP.

\begin{center}\rule{0.5\linewidth}{0.5pt}\end{center}

\subsection{0x731 --- SEB\_ErrInfo (SYNTREE SEB Error
Detail)}\label{x731-seb_errinfo-syntree-seb-error-detail}

\begin{longtable}[]{@{}ll@{}}
\toprule\noalign{}
Property & Value \\
\midrule\noalign{}
\endhead
\bottomrule\noalign{}
\endlastfoot
\textbf{Sender} & SYNTREE SEB (brake module) \\
\textbf{Receiver(s)} & SYS ESP32-S3 \\
\textbf{DLC} & 8 \\
\textbf{Period} & 100 ms (10 Hz) \\
\textbf{Endianness} & Motorola LSB (little-endian) \\
\end{longtable}

Detailed fault flags. Each bit is an independent fault indicator. 1 = fault
active.

\begin{longtable}[]{@{}
  >{\raggedright\arraybackslash}p{(\columnwidth - 10\tabcolsep) * \real{0.1429}}
  >{\raggedright\arraybackslash}p{(\columnwidth - 10\tabcolsep) * \real{0.1964}}
  >{\raggedright\arraybackslash}p{(\columnwidth - 10\tabcolsep) * \real{0.0893}}
  >{\raggedright\arraybackslash}p{(\columnwidth - 10\tabcolsep) * \real{0.1071}}
  >{\raggedright\arraybackslash}p{(\columnwidth - 10\tabcolsep) * \real{0.2321}}
  >{\raggedright\arraybackslash}p{(\columnwidth - 10\tabcolsep) * \real{0.2321}}@{}}
\toprule\noalign{}
\begin{minipage}[b]{\linewidth}\raggedright
Signal
\end{minipage} & \begin{minipage}[b]{\linewidth}\raggedright
Start bit
\end{minipage} & \begin{minipage}[b]{\linewidth}\raggedright
Len
\end{minipage} & \begin{minipage}[b]{\linewidth}\raggedright
Type
\end{minipage} & \begin{minipage}[b]{\linewidth}\raggedright
Fault Level
\end{minipage} & \begin{minipage}[b]{\linewidth}\raggedright
Description
\end{minipage} \\
\midrule\noalign{}
\endhead
\bottomrule\noalign{}
\endlastfoot
\texttt{SEB\_ECUUnderVolt\_Err} & 0 & 1 & bool & L2 & Controller Undervoltage \\
\texttt{SEB\_ECUOverVolt\_Err} & 1 & 1 & bool & L2 & Controller Overvoltage \\
\texttt{SEB\_CanCom\_Err} & 2 & 1 & bool & \textbf{L3} & CAN Communication
Fault \\
\texttt{SEB\_ECUTemp\_Err} & 3 & 1 & bool & \textbf{L3} & Controller Temperature
Fault \\
\texttt{SEB\_Domain\_drive\_SC\_Err} & 4 & 1 & bool & \textbf{L3} & Domain Drive
Short Circuit \\
\texttt{SEB\_Domain\_drive\_V\_Err} & 5 & 1 & bool & \textbf{L3} & Domain Drive
Voltage Fault \\
\texttt{SEB\_Domain\_drive\_T\_Err} & 6 & 1 & bool & \textbf{L3} & Domain Drive
Temperature Fault \\
\texttt{SEB\_AngleSensor\_P\_OC\_Err} & 7 & 1 & bool & \textbf{L3} & Angle
Sensor P Open Circuit \\
\texttt{SEB\_AngleSensor\_P\_AF\_Err} & 8 & 1 & bool & \textbf{L3} & Angle
Sensor P Mainboard Abnormal \\
\texttt{SEB\_AngleSensor\_S\_OC\_Err} & 9 & 1 & bool & \textbf{L3} & Angle
Sensor S Open Circuit \\
\texttt{SEB\_AngleSensor\_S\_AF\_Err} & 10 & 1 & bool & \textbf{L3} & Angle
Sensor S Sub-board Abnormal \\
\texttt{SEB\_NOPreSensor\_Err} & 11 & 1 & bool & \textbf{L3} & Unconnected Oil
Pressure Sensor \\
(reserved) & 12 & 1 & --- & --- & \\
\texttt{SEB\_SensorUCL\_Err} & 13 & 1 & bool & \textbf{L3} & Sensor Plausibility
Fault \\
\texttt{SEB\_Alignment\_Err} & 14 & 1 & bool & L2 & Alignment Fault \\
\texttt{SEB\_AngleOver\_Err} & 15 & 1 & bool & L2 & Angle Out of Bounds \\
(reserved) & 16 & 1 & --- & --- & \\
\texttt{SEB\_Mtr\_Stall\_Err} & 17 & 1 & bool & \textbf{L3} & Motor Stall
Fault \\
\texttt{SEB\_MtrDC\_Err} & 18 & 1 & bool & \textbf{L3} & Motor Disconnect
Fault \\
\texttt{SEB\_Oil\_Err} & 19 & 1 & bool & L2 & Oil Pressure Error \\
\texttt{SEB\_InitOil\_Err} & 20 & 1 & bool & \textbf{L3} & Initial Oil Pressure
Fault \\
\texttt{SEB\_SentValue\_Err} & 21 & 1 & bool & \textbf{L3} & Send Value Error \\
\texttt{SEB\_Mtr\_NoLoad\_Err} & 22 & 1 & bool & \textbf{L3} & Motor No-load
Fault \\
(reserved) & 23 & 1 & --- & --- & \\
\texttt{SEB\_PreSensorOver\_Err} & 24 & 1 & bool & L2 & Oil Pressure Sensor
Overvoltage \\
\texttt{SEB\_LowVolt\_Charging\_Err} & 25 & 1 & bool & L2 & Low Voltage Charging
Failure \\
(reserved) & 26 & 38 & --- & --- & Bits 26--63 (byte 3 bits 2--7 + bytes 4--7).
Not enumerated in CSV. \\
\end{longtable}

\textbf{Byte layout} (little-endian):

\begin{longtable}[]{@{}
  >{\raggedright\arraybackslash}p{(\columnwidth - 16\tabcolsep) * \real{0.2000}}
  >{\raggedright\arraybackslash}p{(\columnwidth - 16\tabcolsep) * \real{0.1000}}
  >{\raggedright\arraybackslash}p{(\columnwidth - 16\tabcolsep) * \real{0.1000}}
  >{\raggedright\arraybackslash}p{(\columnwidth - 16\tabcolsep) * \real{0.1000}}
  >{\raggedright\arraybackslash}p{(\columnwidth - 16\tabcolsep) * \real{0.1000}}
  >{\raggedright\arraybackslash}p{(\columnwidth - 16\tabcolsep) * \real{0.1000}}
  >{\raggedright\arraybackslash}p{(\columnwidth - 16\tabcolsep) * \real{0.1000}}
  >{\raggedright\arraybackslash}p{(\columnwidth - 16\tabcolsep) * \real{0.1000}}
  >{\raggedright\arraybackslash}p{(\columnwidth - 16\tabcolsep) * \real{0.1000}}@{}}
\toprule\noalign{}
\begin{minipage}[b]{\linewidth}\raggedright
Byte
\end{minipage} & \begin{minipage}[b]{\linewidth}\raggedright
0
\end{minipage} & \begin{minipage}[b]{\linewidth}\raggedright
1
\end{minipage} & \begin{minipage}[b]{\linewidth}\raggedright
2
\end{minipage} & \begin{minipage}[b]{\linewidth}\raggedright
3
\end{minipage} & \begin{minipage}[b]{\linewidth}\raggedright
4
\end{minipage} & \begin{minipage}[b]{\linewidth}\raggedright
5
\end{minipage} & \begin{minipage}[b]{\linewidth}\raggedright
6
\end{minipage} & \begin{minipage}[b]{\linewidth}\raggedright
7
\end{minipage} \\
\midrule\noalign{}
\endhead
\bottomrule\noalign{}
\endlastfoot
Content & faults {[}7:0{]} & faults {[}15:8{]} & faults {[}23:16{]} & faults
{[}31:24{]} & rsvd & rsvd & rsvd & rsvd \\
\end{longtable}

\begin{quote}
\textbf{Safety note:} 14 of 23 documented faults are L3 (severe, request
shutdown). SYS must subscribe to this message and escalate any L3 fault to ESTOP
via \texttt{0x001}. The summary \texttt{SEB\_Error\_Status} in \texttt{0x721}
only provides a 2-bit aggregate level; this message reveals \emph{which} sensor
or subsystem failed and at what severity.
\end{quote}

\begin{center}\rule{0.5\linewidth}{0.5pt}\end{center}

\subsection{0x741 --- SEB\_Version (SYNTREE SEB Firmware
Version)}\label{x741-seb_version-syntree-seb-firmware-version}

\begin{longtable}[]{@{}ll@{}}
\toprule\noalign{}
Property & Value \\
\midrule\noalign{}
\endhead
\bottomrule\noalign{}
\endlastfoot
\textbf{Sender} & SYNTREE SEB (brake module) \\
\textbf{Receiver(s)} & SYS ESP32-S3 \\
\textbf{DLC} & 8 \\
\textbf{Period} & 1000 ms (1 Hz) \\
\textbf{Endianness} & Motorola LSB (little-endian) \\
\end{longtable}

\begin{longtable}[]{@{}
  >{\raggedright\arraybackslash}p{(\columnwidth - 18\tabcolsep) * \real{0.1081}}
  >{\raggedright\arraybackslash}p{(\columnwidth - 18\tabcolsep) * \real{0.1486}}
  >{\raggedright\arraybackslash}p{(\columnwidth - 18\tabcolsep) * \real{0.0676}}
  >{\raggedright\arraybackslash}p{(\columnwidth - 18\tabcolsep) * \real{0.0811}}
  >{\raggedright\arraybackslash}p{(\columnwidth - 18\tabcolsep) * \real{0.0946}}
  >{\raggedright\arraybackslash}p{(\columnwidth - 18\tabcolsep) * \real{0.1081}}
  >{\raggedright\arraybackslash}p{(\columnwidth - 18\tabcolsep) * \real{0.0676}}
  >{\raggedright\arraybackslash}p{(\columnwidth - 18\tabcolsep) * \real{0.0676}}
  >{\raggedright\arraybackslash}p{(\columnwidth - 18\tabcolsep) * \real{0.0811}}
  >{\raggedright\arraybackslash}p{(\columnwidth - 18\tabcolsep) * \real{0.1757}}@{}}
\toprule\noalign{}
\begin{minipage}[b]{\linewidth}\raggedright
Signal
\end{minipage} & \begin{minipage}[b]{\linewidth}\raggedright
Start bit
\end{minipage} & \begin{minipage}[b]{\linewidth}\raggedright
Len
\end{minipage} & \begin{minipage}[b]{\linewidth}\raggedright
Type
\end{minipage} & \begin{minipage}[b]{\linewidth}\raggedright
Scale
\end{minipage} & \begin{minipage}[b]{\linewidth}\raggedright
Offset
\end{minipage} & \begin{minipage}[b]{\linewidth}\raggedright
Min
\end{minipage} & \begin{minipage}[b]{\linewidth}\raggedright
Max
\end{minipage} & \begin{minipage}[b]{\linewidth}\raggedright
Unit
\end{minipage} & \begin{minipage}[b]{\linewidth}\raggedright
Description
\end{minipage} \\
\midrule\noalign{}
\endhead
\bottomrule\noalign{}
\endlastfoot
\texttt{SEB\_SW\_Version} & 0 & 8 & u8 & 0.01 & 0 & 0 & 25.5 & --- & Software
version (e.g., 0xC8 = 2.00) \\
\texttt{SEB\_HW\_Version} & 8 & 8 & u8 & 0.1 & 0 & 0 & 25.5 & --- & Hardware
version (e.g., 0x0D = 1.3) \\
(reserved) & 16 & 48 & --- & --- & --- & --- & --- & --- & Bytes 2--7 \\
\end{longtable}

\textbf{SYS usage}: Log on boot for compatibility check and field diagnostics.
Report via \texttt{0x600\ SYS\_DIAG\_RPT}.

\begin{center}\rule{0.5\linewidth}{0.5pt}\end{center}

\subsection{0x6FB --- SEB\_Test (SYNTREE SEB
Telemetry)}\label{x6fb-seb_test-syntree-seb-telemetry}

\begin{longtable}[]{@{}ll@{}}
\toprule\noalign{}
Property & Value \\
\midrule\noalign{}
\endhead
\bottomrule\noalign{}
\endlastfoot
\textbf{Sender} & SYNTREE SEB (brake module) \\
\textbf{Receiver(s)} & SYS ESP32-S3 \\
\textbf{DLC} & 8 \\
\textbf{Period} & 10 ms (100 Hz) \\
\textbf{Endianness} & Motorola LSB (little-endian) \\
\end{longtable}

\begin{longtable}[]{@{}
  >{\raggedright\arraybackslash}p{(\columnwidth - 18\tabcolsep) * \real{0.1081}}
  >{\raggedright\arraybackslash}p{(\columnwidth - 18\tabcolsep) * \real{0.1486}}
  >{\raggedright\arraybackslash}p{(\columnwidth - 18\tabcolsep) * \real{0.0676}}
  >{\raggedright\arraybackslash}p{(\columnwidth - 18\tabcolsep) * \real{0.0811}}
  >{\raggedright\arraybackslash}p{(\columnwidth - 18\tabcolsep) * \real{0.0946}}
  >{\raggedright\arraybackslash}p{(\columnwidth - 18\tabcolsep) * \real{0.1081}}
  >{\raggedright\arraybackslash}p{(\columnwidth - 18\tabcolsep) * \real{0.0676}}
  >{\raggedright\arraybackslash}p{(\columnwidth - 18\tabcolsep) * \real{0.0676}}
  >{\raggedright\arraybackslash}p{(\columnwidth - 18\tabcolsep) * \real{0.0811}}
  >{\raggedright\arraybackslash}p{(\columnwidth - 18\tabcolsep) * \real{0.1757}}@{}}
\toprule\noalign{}
\begin{minipage}[b]{\linewidth}\raggedright
Signal
\end{minipage} & \begin{minipage}[b]{\linewidth}\raggedright
Start bit
\end{minipage} & \begin{minipage}[b]{\linewidth}\raggedright
Len
\end{minipage} & \begin{minipage}[b]{\linewidth}\raggedright
Type
\end{minipage} & \begin{minipage}[b]{\linewidth}\raggedright
Scale
\end{minipage} & \begin{minipage}[b]{\linewidth}\raggedright
Offset
\end{minipage} & \begin{minipage}[b]{\linewidth}\raggedright
Min
\end{minipage} & \begin{minipage}[b]{\linewidth}\raggedright
Max
\end{minipage} & \begin{minipage}[b]{\linewidth}\raggedright
Unit
\end{minipage} & \begin{minipage}[b]{\linewidth}\raggedright
Description
\end{minipage} \\
\midrule\noalign{}
\endhead
\bottomrule\noalign{}
\endlastfoot
(reserved) & 0 & 8 & --- & --- & --- & --- & --- & --- & Byte 0 \\
\texttt{SEB\_MtrCurr} & 8 & 16 & i16 & 0.0078125 & 0 & -255 & 255 & A & Motor
current. Bytes 1--2. ±1.99 A range. \\
\texttt{SEB\_ECUTemp} & 24 & 16 & u16 & 0.5 & 0 & -40 & 215 & °C & ECU
temperature. Bytes 3--4. \\
\texttt{SEB\_PowVolt} & 40 & 16 & u16 & 0.00390625 & 0 & 0 & 32 & V & Power
supply voltage. Bytes 5--6. \\
(reserved) & 56 & 8 & --- & --- & --- & --- & --- & --- & Byte 7 \\
\end{longtable}

\begin{quote}
\textbf{Note:} CSV uses byte-local start-bit numbering for this message (start
bit = offset within start byte). Converted to absolute Motorola LSB above.
\textbf{SYS usage:} Monitor \texttt{SEB\_MtrCurr} for mechanical binding / motor
degradation trends. \texttt{SEB\_ECUTemp} for over-temperature early warning.
\end{quote}

\begin{center}\rule{0.5\linewidth}{0.5pt}\end{center}

\subsection{0x7FD --- RT\_HEARTBEAT
(low-level)}\label{x7fd-rt_heartbeat-low-level}

\begin{longtable}[]{@{}ll@{}}
\toprule\noalign{}
Property & Value \\
\midrule\noalign{}
\endhead
\bottomrule\noalign{}
\endlastfoot
\textbf{Sender} & RT \\
\textbf{Receiver(s)} & SYS \\
\textbf{DLC} & 1 \\
\textbf{Period} & 2 Hz (500 ms) \\
\textbf{Timeout} & 1000ms (2 missed frames) → SYS triggers ESTOP (AUTO only) \\
\end{longtable}

\begin{longtable}[]{@{}
  >{\raggedright\arraybackslash}p{(\columnwidth - 8\tabcolsep) * \real{0.1860}}
  >{\raggedright\arraybackslash}p{(\columnwidth - 8\tabcolsep) * \real{0.2558}}
  >{\raggedright\arraybackslash}p{(\columnwidth - 8\tabcolsep) * \real{0.1163}}
  >{\raggedright\arraybackslash}p{(\columnwidth - 8\tabcolsep) * \real{0.1395}}
  >{\raggedright\arraybackslash}p{(\columnwidth - 8\tabcolsep) * \real{0.3023}}@{}}
\toprule\noalign{}
\begin{minipage}[b]{\linewidth}\raggedright
Signal
\end{minipage} & \begin{minipage}[b]{\linewidth}\raggedright
Start bit
\end{minipage} & \begin{minipage}[b]{\linewidth}\raggedright
Len
\end{minipage} & \begin{minipage}[b]{\linewidth}\raggedright
Type
\end{minipage} & \begin{minipage}[b]{\linewidth}\raggedright
Description
\end{minipage} \\
\midrule\noalign{}
\endhead
\bottomrule\noalign{}
\endlastfoot
\texttt{alive\_ctr} & 0 & 8 & u8 & Increments every frame (wraps at 255). Frozen
= hung RT. \\
\end{longtable}

\texttt{0x204} staleness check at 200ms provides faster detection. RT sends
\texttt{0x7FD} independently on both buses (per-bus, NOT bridged; separate
counters per bus).

\begin{center}\rule{0.5\linewidth}{0.5pt}\end{center}

\subsection{0x7FE --- SYS\_HEARTBEAT
(low-level)}\label{x7fe-sys_heartbeat-low-level}

\begin{longtable}[]{@{}
  >{\raggedright\arraybackslash}p{(\columnwidth - 2\tabcolsep) * \real{0.5882}}
  >{\raggedright\arraybackslash}p{(\columnwidth - 2\tabcolsep) * \real{0.4118}}@{}}
\toprule\noalign{}
\begin{minipage}[b]{\linewidth}\raggedright
Property
\end{minipage} & \begin{minipage}[b]{\linewidth}\raggedright
Value
\end{minipage} \\
\midrule\noalign{}
\endhead
\bottomrule\noalign{}
\endlastfoot
\textbf{Sender} & SYS \\
\textbf{Receiver(s)} & RT \\
\textbf{DLC} & 1 \\
\textbf{Period} & 10 Hz (100 ms) \\
\textbf{Timeout} & 200ms (2 missed frames at 10 Hz) → RT takes over
\texttt{0x7B9} (stroke=max) + sends CAN \texttt{0x001}. Total brake gap
≤220ms. \\
\end{longtable}

\begin{longtable}[]{@{}
  >{\raggedright\arraybackslash}p{(\columnwidth - 8\tabcolsep) * \real{0.1860}}
  >{\raggedright\arraybackslash}p{(\columnwidth - 8\tabcolsep) * \real{0.2558}}
  >{\raggedright\arraybackslash}p{(\columnwidth - 8\tabcolsep) * \real{0.1163}}
  >{\raggedright\arraybackslash}p{(\columnwidth - 8\tabcolsep) * \real{0.1395}}
  >{\raggedright\arraybackslash}p{(\columnwidth - 8\tabcolsep) * \real{0.3023}}@{}}
\toprule\noalign{}
\begin{minipage}[b]{\linewidth}\raggedright
Signal
\end{minipage} & \begin{minipage}[b]{\linewidth}\raggedright
Start bit
\end{minipage} & \begin{minipage}[b]{\linewidth}\raggedright
Len
\end{minipage} & \begin{minipage}[b]{\linewidth}\raggedright
Type
\end{minipage} & \begin{minipage}[b]{\linewidth}\raggedright
Description
\end{minipage} \\
\midrule\noalign{}
\endhead
\bottomrule\noalign{}
\endlastfoot
\texttt{alive\_ctr} & 0 & 8 & u8 & Increments every frame (wraps at 255). Frozen
= hung SYS. \\
\end{longtable}

SYS heartbeat never leaves low bus. Startup grace period: 3 seconds (both
heartbeat monitors).

\begin{center}\rule{0.5\linewidth}{0.5pt}\end{center}

\section{2. High-Level CAN Bus}\label{high-level-can-bus}

Nodes: Jetson Orin, RT ESP32-S3 (MCP2515 SPI).

\begin{center}\rule{0.5\linewidth}{0.5pt}\end{center}

\subsection{0x001 --- SAFETY\_ESTOP
(high-level)}\label{x001-safety_estop-high-level}

\begin{longtable}[]{@{}ll@{}}
\toprule\noalign{}
Property & Value \\
\midrule\noalign{}
\endhead
\bottomrule\noalign{}
\endlastfoot
\textbf{Sender} & Jetson or RT (fwd from low) \\
\textbf{Receiver(s)} & Jetson, RT \\
\textbf{DLC} & 0 \\
\textbf{Period} & On event \\
\end{longtable}

Bridged by RT between buses.

\begin{center}\rule{0.5\linewidth}{0.5pt}\end{center}

\subsection{0x011 --- SYS\_SAFETY\_STS
(forwarded)}\label{x011-sys_safety_sts-forwarded}

Forwarded from low-level by RT. Same payload layout as §1 \texttt{0x011}.

\begin{center}\rule{0.5\linewidth}{0.5pt}\end{center}

\subsection{0x120 --- SYS\_THROTTLE\_STS
(forwarded)}\label{x120-sys_throttle_sts-forwarded}

Forwarded from low-level by RT. Same payload layout as §1 \texttt{0x120}.

\begin{center}\rule{0.5\linewidth}{0.5pt}\end{center}

\subsection{0x210 --- RT\_STATE\_RPT}\label{x210-rt_state_rpt}

\begin{longtable}[]{@{}ll@{}}
\toprule\noalign{}
Property & Value \\
\midrule\noalign{}
\endhead
\bottomrule\noalign{}
\endlastfoot
\textbf{Sender} & RT \\
\textbf{Receiver(s)} & Jetson \\
\textbf{DLC} & 3 \\
\textbf{Period} & 10 Hz \\
\end{longtable}

\begin{longtable}[]{@{}lllllll@{}}
\toprule\noalign{}
Signal & Start bit & Len & Type & Min & Max & Unit \\
\midrule\noalign{}
\endhead
\bottomrule\noalign{}
\endlastfoot
\texttt{RT\_Mode} & 0 & 8 & u8 (enum) & 0 & 2 & --- \\
\texttt{RT\_SteerValid} & 8 & 8 & u8 (bool) & 0 & 1 & --- \\
\texttt{RT\_Reversing} & 16 & 8 & u8 (bool) & 0 & 1 & --- \\
\end{longtable}

Byte layout (big-endian): Byte 0=mode, 1=steer\_valid, 2=reversing.

\begin{center}\rule{0.5\linewidth}{0.5pt}\end{center}

\subsection{0x220 --- RT\_PID\_RPT}\label{x220-rt_pid_rpt}

\begin{longtable}[]{@{}ll@{}}
\toprule\noalign{}
Property & Value \\
\midrule\noalign{}
\endhead
\bottomrule\noalign{}
\endlastfoot
\textbf{Sender} & RT \\
\textbf{Receiver(s)} & Jetson \\
\textbf{DLC} & 6 \\
\textbf{Period} & 10 Hz \\
\end{longtable}

\begin{longtable}[]{@{}lllll@{}}
\toprule\noalign{}
Signal & Start bit & Len & Type & Unit \\
\midrule\noalign{}
\endhead
\bottomrule\noalign{}
\endlastfoot
\texttt{RT\_PidSetpoint} & 0 & 16 & i16 & mm/s \\
\texttt{RT\_PidMeasured} & 16 & 16 & i16 & mm/s \\
\texttt{RT\_PidOutput} & 32 & 16 & i16 & --- \\
\end{longtable}

Byte layout (big-endian): Bytes 0-1=sp, 2-3=meas, 4-5=out.

\begin{center}\rule{0.5\linewidth}{0.5pt}\end{center}

\subsection{0x300 --- HOST\_DRIVE\_CMD}\label{x300-host_drive_cmd}

\begin{longtable}[]{@{}ll@{}}
\toprule\noalign{}
Property & Value \\
\midrule\noalign{}
\endhead
\bottomrule\noalign{}
\endlastfoot
\textbf{Sender} & Jetson \\
\textbf{Receiver(s)} & RT \\
\textbf{DLC} & 8 \\
\textbf{Period} & ≤100 Hz \\
\end{longtable}

\begin{longtable}[]{@{}lllllllll@{}}
\toprule\noalign{}
Signal & Start bit & Len & Type & Scale & Offset & Min & Max & Unit \\
\midrule\noalign{}
\endhead
\bottomrule\noalign{}
\endlastfoot
\texttt{HOST\_DriveSpeed} & 0 & 32 & i32 & 1 & 0 & -500 & 3000 & mm/s \\
\texttt{HOST\_YawRate} & 32 & 24 & i24 & 1 & 0 & -3000 & 3000 & mrad/s \\
\texttt{HOST\_Gear} & 56 & 8 & u8 & 1 & 0 & 0 & 3 & enum (0=N,1=D,2=S,3=R) \\
\end{longtable}

Byte layout (big-endian): Bytes 0-3=speed, 4-6=yaw{[}23:0{]}, 7=gear.

ROS 2 conversion: \texttt{speed\_mmps\ =\ linear.x\ ×\ 1000},
\texttt{yaw\_rate\_mrad\_s\ =\ angular.z\ ×\ 1000}.

\begin{center}\rule{0.5\linewidth}{0.5pt}\end{center}

\subsection{0x301 --- HOST\_BRAKE\_REQ}\label{x301-host_brake_req}

\begin{longtable}[]{@{}ll@{}}
\toprule\noalign{}
Property & Value \\
\midrule\noalign{}
\endhead
\bottomrule\noalign{}
\endlastfoot
\textbf{Sender} & Jetson \\
\textbf{Receiver(s)} & RT \\
\textbf{DLC} & 4 \\
\textbf{Period} & On demand \\
\end{longtable}

\begin{longtable}[]{@{}lllll@{}}
\toprule\noalign{}
Signal & Start bit & Len & Type & Unit \\
\midrule\noalign{}
\endhead
\bottomrule\noalign{}
\endlastfoot
\texttt{HOST\_BrakePressure} & 0 & 32 & i32 & kPa \\
\end{longtable}

Byte layout (big-endian): Bytes 0-3. RT arbitrates: max(RT\_computed,
HOST\_request). Result forwarded to SYS via \texttt{0x205} RT\_BRAKE\_CMD (i32
kPa, 50 Hz).

\begin{center}\rule{0.5\linewidth}{0.5pt}\end{center}

\subsection{0x302 --- HOST\_LIGHT\_CMD}\label{x302-host_light_cmd}

\begin{longtable}[]{@{}ll@{}}
\toprule\noalign{}
Property & Value \\
\midrule\noalign{}
\endhead
\bottomrule\noalign{}
\endlastfoot
\textbf{Sender} & Jetson \\
\textbf{Receiver(s)} & RT (→ SYS) \\
\textbf{DLC} & 1 \\
\textbf{Period} & On change \\
\end{longtable}

Layout identical to low-level \texttt{0x302}. RT forwards transparently.

\begin{center}\rule{0.5\linewidth}{0.5pt}\end{center}

\subsection{0x400 --- RT\_OBSTACLE\_RPT}\label{x400-rt_obstacle_rpt}

\begin{longtable}[]{@{}ll@{}}
\toprule\noalign{}
Property & Value \\
\midrule\noalign{}
\endhead
\bottomrule\noalign{}
\endlastfoot
\textbf{Sender} & RT \\
\textbf{Receiver(s)} & Jetson \\
\textbf{DLC} & 4 \\
\textbf{Period} & 10 Hz \\
\end{longtable}

\begin{longtable}[]{@{}lllllll@{}}
\toprule\noalign{}
Signal & Start bit & Len & Type & Min & Max & Unit \\
\midrule\noalign{}
\endhead
\bottomrule\noalign{}
\endlastfoot
\texttt{RT\_ObstacleDistance} & 0 & 32 & u32 & 0 & 2³²−1 & mm \\
\end{longtable}

UINT32\_MAX = no reading / timeout.

RT reports min obstacle distance to Jetson at 10 Hz. Jetson perception
(LiDAR/camera) sends distance via internal path to RT, which rebroadcasts on the
high bus for logging/monitoring.

\begin{center}\rule{0.5\linewidth}{0.5pt}\end{center}

\subsection{0x600 --- SYS\_DIAG\_RPT
(forwarded)}\label{x600-sys_diag_rpt-forwarded}

Forwarded from low-level by RT. Same layout as §1 \texttt{0x600}.

\begin{center}\rule{0.5\linewidth}{0.5pt}\end{center}

\subsection{0x7FD --- RT\_HEARTBEAT
(high-level)}\label{x7fd-rt_heartbeat-high-level}

\begin{longtable}[]{@{}ll@{}}
\toprule\noalign{}
Property & Value \\
\midrule\noalign{}
\endhead
\bottomrule\noalign{}
\endlastfoot
\textbf{Sender} & RT \\
\textbf{Receiver(s)} & Jetson \\
\textbf{DLC} & 1 \\
\textbf{Period} & 2 Hz (500 ms) \\
\textbf{Timeout} & 1500ms (3 missed frames) → Jetson stops publishing
\texttt{/cmd\_vel} \\
\end{longtable}

\begin{longtable}[]{@{}
  >{\raggedright\arraybackslash}p{(\columnwidth - 8\tabcolsep) * \real{0.1860}}
  >{\raggedright\arraybackslash}p{(\columnwidth - 8\tabcolsep) * \real{0.2558}}
  >{\raggedright\arraybackslash}p{(\columnwidth - 8\tabcolsep) * \real{0.1163}}
  >{\raggedright\arraybackslash}p{(\columnwidth - 8\tabcolsep) * \real{0.1395}}
  >{\raggedright\arraybackslash}p{(\columnwidth - 8\tabcolsep) * \real{0.3023}}@{}}
\toprule\noalign{}
\begin{minipage}[b]{\linewidth}\raggedright
Signal
\end{minipage} & \begin{minipage}[b]{\linewidth}\raggedright
Start bit
\end{minipage} & \begin{minipage}[b]{\linewidth}\raggedright
Len
\end{minipage} & \begin{minipage}[b]{\linewidth}\raggedright
Type
\end{minipage} & \begin{minipage}[b]{\linewidth}\raggedright
Description
\end{minipage} \\
\midrule\noalign{}
\endhead
\bottomrule\noalign{}
\endlastfoot
\texttt{alive\_ctr} & 0 & 8 & u8 & Increments every frame (wraps at 255). Frozen
= hung RT. \\
\end{longtable}

RT sends \texttt{0x7FD} independently on both buses (per-bus, NOT bridged).

\begin{center}\rule{0.5\linewidth}{0.5pt}\end{center}

\subsection{0x7FC --- JETSON\_HEARTBEAT
(high-level)}\label{x7fc-jetson_heartbeat-high-level}

\begin{longtable}[]{@{}
  >{\raggedright\arraybackslash}p{(\columnwidth - 2\tabcolsep) * \real{0.5882}}
  >{\raggedright\arraybackslash}p{(\columnwidth - 2\tabcolsep) * \real{0.4118}}@{}}
\toprule\noalign{}
\begin{minipage}[b]{\linewidth}\raggedright
Property
\end{minipage} & \begin{minipage}[b]{\linewidth}\raggedright
Value
\end{minipage} \\
\midrule\noalign{}
\endhead
\bottomrule\noalign{}
\endlastfoot
\textbf{Sender} & Jetson \\
\textbf{Receiver(s)} & RT \\
\textbf{DLC} & 1 \\
\textbf{Period} & 2 Hz (500 ms) \\
\textbf{Timeout} & 1500ms (3 missed frames) → RT zeroes \texttt{0x204} + stops
\texttt{0x169} (controlled stop) \\
\end{longtable}

\begin{longtable}[]{@{}
  >{\raggedright\arraybackslash}p{(\columnwidth - 8\tabcolsep) * \real{0.1860}}
  >{\raggedright\arraybackslash}p{(\columnwidth - 8\tabcolsep) * \real{0.2558}}
  >{\raggedright\arraybackslash}p{(\columnwidth - 8\tabcolsep) * \real{0.1163}}
  >{\raggedright\arraybackslash}p{(\columnwidth - 8\tabcolsep) * \real{0.1395}}
  >{\raggedright\arraybackslash}p{(\columnwidth - 8\tabcolsep) * \real{0.3023}}@{}}
\toprule\noalign{}
\begin{minipage}[b]{\linewidth}\raggedright
Signal
\end{minipage} & \begin{minipage}[b]{\linewidth}\raggedright
Start bit
\end{minipage} & \begin{minipage}[b]{\linewidth}\raggedright
Len
\end{minipage} & \begin{minipage}[b]{\linewidth}\raggedright
Type
\end{minipage} & \begin{minipage}[b]{\linewidth}\raggedright
Description
\end{minipage} \\
\midrule\noalign{}
\endhead
\bottomrule\noalign{}
\endlastfoot
\texttt{alive\_ctr} & 0 & 8 & u8 & Increments every frame (wraps at 255). Frozen
= hung Jetson. \\
\end{longtable}

Jetson is QM, not safety-critical. Heartbeat loss triggers controlled stop, not
ESTOP.

\begin{center}\rule{0.5\linewidth}{0.5pt}\end{center}

\section{3. CAN ID Summary}\label{can-id-summary}

\subsection{Low-level bus}\label{low-level-bus}

\begin{longtable}[]{@{}llllll@{}}
\toprule\noalign{}
ID & Name & Sender & Receiver & DLC & Rate \\
\midrule\noalign{}
\endhead
\bottomrule\noalign{}
\endlastfoot
\texttt{0x001} & SAFETY\_ESTOP & RT, SYS & All & 0 & Event \\
\texttt{0x011} & SYS\_SAFETY\_STS & SYS & RT (→Jetson) & 2 & 5 Hz \\
\texttt{0x012} & SYS\_DCDC\_CMD & SYS & DC-DC & 1 & Change \\
\texttt{0x110} & SYS\_MODE\_CMD & SYS & RT & 1 & Change \\
\texttt{0x120} & SYS\_THROTTLE\_STS & MTR & RT (→Jetson) & 2 & 100 Hz \\
\texttt{0x169} & VCU\_SES\_REQ & RT & EPS-C & 8 & \textbf{50 Hz} \\
\texttt{0x201} & SES\_STATUS & EPS-C & RT & 8 & 100 Hz \\
\texttt{0x202} & SES\_ErrInfo & EPS-C & RT & 8 & 10 Hz \\
\texttt{0x203} & SES\_Version & EPS-C & RT & 8 & 1 Hz \\
\texttt{0x204} & RT\_DRIVE\_CMD & RT & SYS, MTR & 5 & 100 Hz \\
\texttt{0x205} & RT\_BRAKE\_CMD & RT & SYS & 4 & \textbf{50 Hz} \\
\texttt{0x302} & HOST\_LIGHT\_CMD & RT (fwd) & SYS & 1 & Change \\
\texttt{0x600} & SYS\_DIAG\_RPT & SYS & RT (→Jetson) & 8 & 1 Hz \\
\texttt{0x6FA} & SES\_Test & EPS-C & RT & 8 & 100 Hz \\
\texttt{0x6FB} & SEB\_Test & SEB & SYS & 8 & 100 Hz \\
\texttt{0x721} & SEB\_STATUS & SEB & SYS & 8 & 100 Hz \\
\texttt{0x731} & SEB\_ErrInfo & SEB & SYS & 8 & 10 Hz \\
\texttt{0x741} & SEB\_Version & SEB & SYS & 8 & 1 Hz \\
\texttt{0x7B9} & VCU\_SEB\_REQ & SYS & SEB & 8 & \textbf{50 Hz} \\
\texttt{0x7FD} & RT\_HEARTBEAT & RT & SYS & 1 & 2 Hz \\
\texttt{0x7FE} & SYS\_HEARTBEAT & SYS & RT & 1 & 10 Hz \\
\end{longtable}

\subsection{High-level bus}\label{high-level-bus}

\begin{longtable}[]{@{}llllll@{}}
\toprule\noalign{}
ID & Name & Sender & Receiver & DLC & Rate \\
\midrule\noalign{}
\endhead
\bottomrule\noalign{}
\endlastfoot
\texttt{0x001} & SAFETY\_ESTOP & Jetson, RT & Jetson, RT & 0 & Event \\
\texttt{0x011} & SYS\_SAFETY\_STS & RT (fwd) & Jetson & 2 & 5 Hz \\
\texttt{0x120} & SYS\_THROTTLE\_STS & RT (fwd) & Jetson & 2 & 100 Hz \\
\texttt{0x210} & RT\_STATE\_RPT & RT & Jetson & 3 & 10 Hz \\
\texttt{0x220} & RT\_PID\_RPT & RT & Jetson & 6 & --- (RESERVED, inactive) \\
\texttt{0x300} & HOST\_DRIVE\_CMD & Jetson & RT & 8 & ≤100 Hz \\
\texttt{0x301} & HOST\_BRAKE\_REQ & Jetson & RT & 4 & Demand \\
\texttt{0x302} & HOST\_LIGHT\_CMD & Jetson & RT (→SYS) & 1 & Change \\
\texttt{0x400} & RT\_OBSTACLE\_RPT & RT & Jetson & 4 & 10 Hz \\
\texttt{0x600} & SYS\_DIAG\_RPT & RT (fwd) & Jetson & 8 & 1 Hz \\
\texttt{0x7FD} & RT\_HEARTBEAT & RT & Jetson & 1 & 2 Hz \\
\texttt{0x7FC} & JETSON\_HEARTBEAT & Jetson & RT & 1 & 2 Hz \\
\end{longtable}

\begin{center}\rule{0.5\linewidth}{0.5pt}\end{center}

\section{4. Forwarding Rules (RT Gateway)}\label{forwarding-rules-rt-gateway}

RT is the only dual-bus node. Every CAN message falls into exactly one of three
categories:

\subsection{Category 1: Transparent forward (same ID, same
payload)}\label{category-1-transparent-forward-same-id-same-payload}

\begin{longtable}[]{@{}ll@{}}
\toprule\noalign{}
Direction & IDs \\
\midrule\noalign{}
\endhead
\bottomrule\noalign{}
\endlastfoot
Low → High & \texttt{0x001}, \texttt{0x011}, \texttt{0x120}, \texttt{0x206},
\texttt{0x600} \\
High → Low & \texttt{0x001}, \texttt{0x302} \\
\end{longtable}

\subsection{Category 2: Consumed by RT → different message
generated}\label{category-2-consumed-by-rt-different-message-generated}

\begin{longtable}[]{@{}
  >{\raggedright\arraybackslash}p{(\columnwidth - 6\tabcolsep) * \real{0.3103}}
  >{\raggedright\arraybackslash}p{(\columnwidth - 6\tabcolsep) * \real{0.1724}}
  >{\raggedright\arraybackslash}p{(\columnwidth - 6\tabcolsep) * \real{0.3448}}
  >{\raggedright\arraybackslash}p{(\columnwidth - 6\tabcolsep) * \real{0.1724}}@{}}
\toprule\noalign{}
\begin{minipage}[b]{\linewidth}\raggedright
Inbound
\end{minipage} & \begin{minipage}[b]{\linewidth}\raggedright
Bus
\end{minipage} & \begin{minipage}[b]{\linewidth}\raggedright
Outbound
\end{minipage} & \begin{minipage}[b]{\linewidth}\raggedright
Bus
\end{minipage} \\
\midrule\noalign{}
\endhead
\bottomrule\noalign{}
\endlastfoot
\texttt{0x300} HOST\_DRIVE\_CMD & High & \texttt{0x204} RT\_DRIVE\_CMD +
\texttt{0x169} VCU\_SES\_REQ & Low \\
\texttt{0x301} HOST\_BRAKE\_REQ & High & \texttt{0x205} RT\_BRAKE\_CMD & Low \\
\end{longtable}

\subsection{Category 3: Bus-local (never forwarded, never
regenerated)}\label{category-3-bus-local-never-forwarded-never-regenerated}

\begin{longtable}[]{@{}
  >{\raggedright\arraybackslash}p{(\columnwidth - 2\tabcolsep) * \real{0.5000}}
  >{\raggedright\arraybackslash}p{(\columnwidth - 2\tabcolsep) * \real{0.5000}}@{}}
\toprule\noalign{}
\begin{minipage}[b]{\linewidth}\raggedright
Bus
\end{minipage} & \begin{minipage}[b]{\linewidth}\raggedright
IDs
\end{minipage} \\
\midrule\noalign{}
\endhead
\bottomrule\noalign{}
\endlastfoot
Low only & \texttt{0x012}, \texttt{0x110}, \texttt{0x169}, \texttt{0x202},
\texttt{0x203}, \texttt{0x204}, \texttt{0x205}, \texttt{0x6FA}, \texttt{0x6FB},
\texttt{0x721}, \texttt{0x731}, \texttt{0x741}, \texttt{0x7B9} \\
Low only & \texttt{0x201} (SYNTREE EPS-C feedback) \\
High only & \texttt{0x210}, \texttt{0x220}, \texttt{0x400} (RT telemetry) \\
Both independent & \texttt{0x7FD}, \texttt{0x7FE}, \texttt{0x7FC} (per-node
heartbeat --- NOT bridged) \\
\end{longtable}

\begin{center}\rule{0.5\linewidth}{0.5pt}\end{center}

\section{5. Priority Groups}\label{priority-groups}

\begin{longtable}[]{@{}
  >{\raggedright\arraybackslash}p{(\columnwidth - 4\tabcolsep) * \real{0.4000}}
  >{\raggedright\arraybackslash}p{(\columnwidth - 4\tabcolsep) * \real{0.4000}}
  >{\raggedright\arraybackslash}p{(\columnwidth - 4\tabcolsep) * \real{0.2000}}@{}}
\toprule\noalign{}
\begin{minipage}[b]{\linewidth}\raggedright
Priority
\end{minipage} & \begin{minipage}[b]{\linewidth}\raggedright
ID Range
\end{minipage} & \begin{minipage}[b]{\linewidth}\raggedright
IDs
\end{minipage} \\
\midrule\noalign{}
\endhead
\bottomrule\noalign{}
\endlastfoot
Highest & \texttt{0x001} & ESTOP \\
Very High & \texttt{0x010}--\texttt{0x01F} & SAFETY\_STATUS, DCDC\_CMD \\
High & \texttt{0x100}--\texttt{0x11F} & MODE\_CMD \\
Medium & \texttt{0x120}--\texttt{0x3FF} & THROTTLE, DRIVE, MTR\_MOTOR\_FBK,
SES\_STATUS/REQ/ErrInfo/Version, DRIVE\_CMD, BRAKE\_REQ, LIGHT\_CMD \\
Low & \texttt{0x400}--\texttt{0x5FF} & OBSTACLE, STATE\_REPORT, PID\_FEEDBACK \\
Lowest & \texttt{0x600}--\texttt{0x7FF} & DIAG, SES\_Test (\texttt{0x6FA}),
SEB\_Test (\texttt{0x6FB}), SEB\_STATUS/ErrInfo/Version
(\texttt{0x721}/\texttt{0x731}/\texttt{0x741}), SEB\_REQ (\texttt{0x7B9}),
HEARTBEAT (\texttt{0x7FC}--\texttt{0x7FE}) \\
\end{longtable}

\begin{quote}
Lower CAN ID = higher bus arbitration priority. Safety-critical frames occupy
\texttt{0x00X}. SYNTREE IDs (\texttt{0x2XX}, \texttt{0x7XX}) are in
medium/lowest ranges per manufacturer assignment.
\end{quote}
